% Revised from the user-supplied LowerBound_Complete (1).tex.
% Proof-detail revision: 26 September 2026.
% Original equation numbers are retained. Lemma 15.10 is added.
% Main changes: representative algebra and explicit target lifts;
% Walsh-norm ghost leakage; local-sign removal across all incident blocks.
% Build the complete PDF with the companion Proof_Flowcharts.pdf present.
\documentclass[11pt]{article}
\usepackage[margin=1in]{geometry}
\usepackage{amsmath,amssymb,amsthm,mathtools,bm}
\usepackage{enumitem}
\usepackage{pdfpages}
\usepackage{aliascnt}
\usepackage{microtype}
\setlength{\emergencystretch}{2em}
\usepackage[colorlinks=true,linkcolor=blue,citecolor=blue,urlcolor=blue]{hyperref}
\usepackage[nameinlink,noabbrev]{cleveref}

\newtheorem{theorem}{Theorem}[section]
\newaliascnt{lemma}{theorem}
\newtheorem{lemma}[lemma]{Lemma}
\aliascntresetthe{lemma}
\newaliascnt{proposition}{theorem}
\newtheorem{proposition}[proposition]{Proposition}
\aliascntresetthe{proposition}
\newaliascnt{corollary}{theorem}
\newtheorem{corollary}[corollary]{Corollary}
\aliascntresetthe{corollary}
\newaliascnt{remark}{theorem}
\newtheorem{remark}[remark]{Remark}
\aliascntresetthe{remark}
\newaliascnt{definition}{theorem}
\newtheorem{definition}[definition]{Definition}
\aliascntresetthe{definition}
\crefname{lemma}{Lemma}{Lemmas}
\crefname{proposition}{Proposition}{Propositions}
\crefname{corollary}{Corollary}{Corollaries}
\crefname{theorem}{Theorem}{Theorems}
\crefname{remark}{Remark}{Remarks}
\crefname{definition}{Definition}{Definitions}

\DeclareMathOperator{\Rad}{Rad}
\DeclareMathOperator{\TV}{TV}
\DeclareMathOperator{\ev}{ev}
\DeclareMathOperator{\supp}{supp}
\DeclareMathOperator{\sym}{Sym}
\newcommand{\E}{\mathbb{E}}
\newcommand{\PP}{\mathbb{P}}
\newcommand{\R}{\mathbb{R}}
\newcommand{\Torus}{\mathbb{T}}
\newcommand{\cW}{\mathcal{W}}
\newcommand{\cA}{\mathcal{A}}
\newcommand{\cB}{\mathcal{B}}
\newcommand{\cE}{\mathcal{E}}
\newcommand{\cH}{\mathcal{H}}
\newcommand{\cK}{\mathcal{K}}
\newcommand{\cL}{\mathcal{L}}
\newcommand{\cP}{\mathcal{P}}
\newcommand{\cS}{\mathcal{S}}
\newcommand{\cJ}{\mathcal{J}}
\newcommand{\norm}[1]{\lVert #1\rVert}
\newcommand{\abs}[1]{\lvert #1\rvert}
\newcommand{\ip}[2]{\langle #1,#2\rangle}
\newcommand{\od}{\mathbin{\odot}}
\newcommand{\one}{\mathbf{1}}
\newcommand{\ja}{\mathsf{j}_a}
\newcommand{\jb}{\mathsf{j}_b}
\newcommand{\jsh}{\mathsf{j}_\circ}
\newcommand{\jfi}{\mathsf{j}_\bullet}
\newcommand{\dg}{\mathfrak{d}}
\newcommand{\K}{\mathsf{K}}
\newcommand{\real}{\mathrm{real}}
\newcommand{\cf}{\mathfrak{c}}
\newcommand{\realize}{\mathfrak r}

\title{Minimax lower bound for pointwise CATE estimation\\
with binary outcomes and rough design:\\
\large the complete low-smoothness theorem}
\author{Zhenghan Li}
\date{5 September 2026}

\begin{document}
% The three original proof flowcharts are retained as vector PDF pages.
% The companion file is included in the complete source archive.
\hypersetup{pageanchor=false}
\IfFileExists{Proof_Flowcharts.pdf}{%
  \includepdf[pages=-,fitpaper=true,pagecommand={}]{Proof_Flowcharts.pdf}%
}{%
  \IfFileExists{LowerBound_Complete (3).pdf}{%
    \includepdf[pages=1-3,fitpaper=true,pagecommand={}]{LowerBound_Complete (3).pdf}%
  }{%
    \PackageWarningNoLine{LowerBound}{Proof_Flowcharts.pdf not found; compiling the complete mathematical text without the three front flowcharts}%
  }%
}
\clearpage
\setcounter{page}{1}
\hypersetup{pageanchor=true}
\maketitle

\begin{abstract}
We prove a pointwise minimax lower bound for conditional average
treatment-effect estimation in a Bernoulli-outcome model whose covariate
density is assumed only to be bounded above and below.  With
\(A_0=\alpha+\beta\), \(s_1=\{(\alpha\wedge1)+(\beta\wedge1)\}/2\),
\(D_1=d+4s_1+2ds_1/\gamma\) and \(\rho=(A_0+2s_1)/D_1\), in the regime
\(d>A_0(2+d/\gamma)\) the minimax risk at an interior point is at least
\(c_\varepsilon n^{-(\rho+\varepsilon)}\) for every \(\varepsilon>0\); when
both nuisance exponents are at most one the bound holds with
\(\varepsilon=0\), and it holds over the uniform-design submodel.  The
proof is organised in three cases according to which nuisance functions
are rough, with the shared tools collected in Part~I.
\end{abstract}

\tableofcontents

\part{Model, theorem, and common tools}

\section{Model and main result}\label{sec:model}

We observe independent copies of \(Z=(X,A,Y)\), where
\begin{equation}
 X\in[0,1]^d,\qquad
 A\mid X=x\sim\operatorname{Bernoulli}\{\pi(x)\},\qquad
 Y\mid(X=x,A=a)\sim\operatorname{Bernoulli}\{\mu_0(x)+a\tau(x)\}.
 \label{eq:model}
\end{equation}
Fix \(x_0\in(0,1)^d\).  For \(s>0\), \(\cH^s(L)\) denotes the usual
isotropic Hölder ball of radius \(L\) on \([0,1]^d\), with the convention
that its elements are restrictions of functions defined on a fixed
neighbourhood of the cube.  Fix constants
\[
 0<\underline f<1<\overline f<\infty,\qquad 0<\kappa<\tfrac14,
\]
and radii \(L_\pi,L_0,L_\tau\) such that the constants \(1/2\) are interior
points of \(\cH^\alpha(L_\pi)\) and \(\cH^\beta(L_0)\) (e.g.\ \(L_\pi,L_0>1/2\)).
The class \(\cP=\cP(\alpha,\beta,\gamma)\) consists of all laws
\eqref{eq:model} with design density \(f\) and nuisance functions such that
\begin{align*}
 &\underline f\le f\le\overline f,\qquad \int f=1,\\
 &\pi\in\cH^\alpha(L_\pi),\qquad \mu_0\in\cH^\beta(L_0),\qquad \tau\in\cH^\gamma(L_\tau),\\
 &\kappa\le\pi\le1-\kappa,\qquad \kappa\le\mu_0\le1-\kappa,\qquad
 \kappa\le\mu_0+\tau\le1-\kappa .
\end{align*}
No smoothness is imposed on the design density: \(\cP\) is the
\emph{rough-design} class, in which the design is constrained only by the
two-sided bound; the design density is one of the objects that the
lower-bound constructions perturb, and the rate below is specific to
this class.  The pointwise absolute-error minimax risk is
\begin{equation}
 R^*_n(x_0)=\inf_{\widehat\tau}\sup_{P\in\cP}\E_P\abs{\widehat\tau(x_0)-\tau_P(x_0)}.
 \label{eq:risk}
\end{equation}
Write
\begin{equation}
 A_0=\alpha+\beta,\qquad
 s_1=\frac{(\alpha\wedge1)+(\beta\wedge1)}2,\qquad
 D_1=d+4s_1+\frac{2ds_1}\gamma,\qquad
 \rho=\frac{A_0+2s_1}{D_1}.
 \label{eq:rho}
\end{equation}

\begin{theorem}[Rough-design pointwise lower bound]\label{thm:main}
Let \(\alpha,\beta,\gamma>0\) satisfy
\[
 d>A_0\Bigl(2+\frac d\gamma\Bigr).
\]
For every fixed \(\varepsilon>0\) there are \(c_\varepsilon>0\) and
\(n_\varepsilon<\infty\) such that
\begin{equation}
 R^*_n(x_0)\ge c_\varepsilon n^{-(\rho+\varepsilon)},\qquad n\ge n_\varepsilon .
 \label{eq:main-bound}
\end{equation}
If \(\alpha,\beta\le1\) the bound holds with \(\varepsilon=0\), and the two
hypotheses realising it have exactly uniform design, so that
\eqref{eq:main-bound} holds over the uniform-design submodel as well.
\end{theorem}

The proof splits according to which nuisance functions are rough
(Hölder exponent at most one):
\[
 \text{(A)}\ \ \alpha,\beta\le1;\qquad
 \text{(B)}\ \ \alpha,\beta>1;\qquad
 \text{(C)}\ \ \alpha\le1<\beta\ \text{ or }\ \beta\le1<\alpha ,
\]
proved in Parts~A, B and~C respectively; every \((\alpha,\beta)\) falls in
exactly one case.  The parts share the model, the coded Walsh
likelihood, the two-point reduction, the component principle, the
design tilt, the scales, the interpolation system and taper on the
torus, the finite moment simplex, the configuration-space Wiener
estimates, and the polarization identity, all collected in Part~I.
They differ in the construction of the two hypotheses:
\begin{itemize}[leftmargin=*]
\item in case~(A) both nuisances are carried by one rough Rademacher
field at the fine scale \(\ell\), the design is uniform, and the proof is
a direct Walsh-parity computation on packet-overlap components;
\item in case~(B) both nuisances are smooth, the outcome is a cubic
function of the propensity (the cubic bridge), and the design
carries a positive moment carrier whose virtual jitter is made of
independent signs attached to the observation sites;
\item in case~(C) one nuisance is rough and the other is carried by the
blocks; the carrier must then be built conditionally on the rough field,
which makes the design informative about it, and the two sides are
matched after design weighting.
\end{itemize}
The formula \(s_1=\{(\alpha\wedge1)+(\beta\wedge1)\}/2\) is the same in all
cases: \(s_1=(\alpha+\beta)/2\) in (A), \(s_1=1\) in (B), and
\(s_1=(1+\alpha\wedge\beta)/2\) in (C).  The regime
\(d>A_0(2+d/\gamma)\) is the low-smoothness regime of Kennedy,
Balakrishnan, Robins and Wasserman \cite{KBRW}; in case~(A) it reads
\(s_1<(d/4)/(1+d/(2\gamma))\) and the rate is
\(n^{-1/(1+d/(2\gamma)+d/(4s_1))}\), their rate.  The exponent \(1\) is
treated as rough in case~(C); there the strict inequality
\(\alpha\vee\beta>1\) is used only through \eqref{C:eq:jsh-over-jfi}.

\begin{remark}[Scope]
The theorem is stated for an interior target point, absolute-error risk,
and Hölder balls containing the fair-sign baseline \(1/2\) as an interior
point.  The strict inequalities \(\underline f<1<\overline f\) place a
neighbourhood of the uniform density inside the design class; the block
carriers of Parts~B and~C are perturbations of the uniform density at
the block scale, and Part~A uses the uniform design itself.
\end{remark}

\section{Common tools}\label{sec:tools}

\subsection{Coded signs and the Walsh likelihood}
Let \(R=2A-1\), \(T=2Y-1\), \(p=2\pi-1\), \(y=2\mu_0-1\).  Relative to
independent fair signs, the parameter \((p,y,\tau)\) has likelihood
\begin{equation}
 L_{p,y,\tau}=1+Rp+T\{y+\tau(1+p)\}+RT\{py+\tau(1+p)\}.
 \label{eq:walsh}
\end{equation}
Thus equality of the four Walsh coefficients is equality of all four
binary cells.  For a function \(g\) of iid Rademacher signs
\(\lambda=(\lambda_m)\), the Walsh norm \(\norm g_W\) is the \(\ell^1\) norm
of its Walsh coefficients; it is an algebra norm,
\(\norm g_\infty\le\norm g_W\), and \(\abs{\E_\lambda g}\le\norm g_W\).  The
\emph{degree} of a Walsh monomial \(\lambda^S=\prod_{m\in S}\lambda_m\) is
\(\abs S\); since \(\lambda_m^2=1\), \(\lambda^S\lambda^{S'}=\lambda^{S\triangle S'}\),
so the degree of a product of monomials is congruent modulo two to the
sum of the degrees.  A Walsh polynomial is \emph{even} (\emph{odd}) if all
its monomials have even (odd) degree, i.e.\ if it is invariant (changes
sign) under \(\lambda\mapsto-\lambda\); an odd Walsh polynomial has zero
mean, and more generally \(\E_\lambda\lambda^S=\one\{S=\varnothing\}\).

\subsection{Partition of unity, bump, and the rough field}
\begin{lemma}[Smooth quadratic partition of unity]\label{lem:partition}
There exists \(\phi\in C_c^\infty(\R^d)\) such that
\begin{equation}
 \sum_{m\in\mathbb Z^d}\phi(u-m)^2=1,\qquad u\in\R^d .
 \label{eq:quadratic-pu}
\end{equation}
\end{lemma}
\begin{proof}
Choose a nonnegative \(\eta\in C_c^\infty(\R)\) which is positive on
\((-1,1)\), so that \(\sum_{k\in\mathbb Z}\eta(u-k)^2>0\) for all \(u\); the
denominator of \(\phi_1(u)=\eta(u)\{\sum_k\eta(u-k)^2\}^{-1/2}\) is smooth,
strictly positive and periodic, so \(\phi_1\in C_c^\infty(\R)\) and
\(\sum_k\phi_1(u-k)^2=1\).  Put \(\phi(u)=\prod_{j=1}^d\phi_1(u_j)\).
\end{proof}
Fix \(G\in C_c^\infty(\R^d)\) with \(0\le G\le1\), \(G(0)=1\), and put
\(G_h(x)=G((x-x_0)/h)\); since \(x_0\) is interior, \(\supp G_h\subset[0,1]^d\)
for small \(h\).  For a fine scale \(\ell\le h\) the \emph{rough packets}
and the \emph{rough field} are
\begin{equation}
 \psi^\flat_m(x)=G_h(x)\phi\Bigl(\frac{x-x_0}\ell-m\Bigr),\qquad
 \Lambda_\lambda(x)=\sum_{m\in\mathbb Z^d}\lambda_m\psi^\flat_m(x),\qquad
 \lambda_m\stackrel{\rm iid}\sim\Rad(\tfrac12),
 \label{eq:rough}
\end{equation}
where only the \(O((h/\ell)^d)\) packets that are not identically zero
need be retained.  Then \(\sum_m\psi^\flat_m(x)^2=G_h(x)^2\); at most
\(N_0\) packets meet a point; \(\abs{\Lambda_\lambda}\le N_0\); and with
\(J(x)=\{m:\psi^\flat_m(x)\ne0\}\), \(J(x)\cap J(x')\ne\varnothing\) implies
\(x,x'\in\supp G_h\) and \(\norm{x-x'}_\infty\le R_0\ell\), for fixed
\(N_0,R_0\).  For fixed \(x\), \(\Lambda_\lambda(x)\) is a Walsh polynomial
of degree one with \(\norm{\Lambda_\lambda(x)}_W\le N_0\); if
\(J(x_1),\dots,J(x_q)\) are pairwise disjoint, \(\Lambda_\lambda(x_1),\dots,\Lambda_\lambda(x_q)\)
are independent.  The moments \(M_j(x):=\E\Lambda_\lambda(x)^j\) are
\begin{equation}
 M_0=1,\qquad M_j=0\ (j\text{ odd}),\qquad M_2=G_h^2,\qquad
 M_4=\mu_4:=3G_h^4-2\sum_m(\psi^\flat_m)^4=G_h^4\{3-2q_4\},
 \label{eq:rough-moments}
\end{equation}
with \(q_4(x)=\sum_{m\in\mathbb Z^d}\phi((x-x_0)/\ell-m)^4\in(0,1]\): indeed
\(\E\Lambda^4=\sum_m\psi_m^4+3\sum_{m\ne m'}\psi_m^2\psi_{m'}^2
=3(\sum_m\psi_m^2)^2-2\sum_m\psi_m^4\), \(\sum_m\phi_m^4\le(\sum_m\phi_m^2)^2=1\),
and odd moments vanish by symmetry.  The function \(q_4\) is smooth at
scale \(\ell\), \(\ell\)-periodic, and does not depend on \(\lambda\); so
\(\mu_4\in[G_h^4,3G_h^4]\).  The centred square
\(V(x):=\Lambda_\lambda(x)^2-G_h(x)^2=\sum_{m\ne m'}\lambda_m\lambda_{m'}\psi^\flat_m(x)\psi^\flat_{m'}(x)\)
is a Walsh polynomial all of whose monomials have degree two, with
\(\norm{V(x)}_W\le N_0^2\).

\begin{lemma}[Hölder scaling]\label{lem:holder-scaling}
If \(g\in C_c^\infty(\R^d)\) and \(0<q\le1\), then for every fixed \(t>0\)
\(\norm{g((\cdot-x_c)/q)}_{C^t}\le C_{g,t}q^{-t}\), and the same bound holds
for finite products and for bounded-overlap sums of translates of such
functions.  The estimate is valid for noninteger \(t\): derivatives up to
\(\lfloor t\rfloor\) scale in the displayed way, and the remaining Hölder
seminorm gains the fractional power of \(q^{-1}\).
\end{lemma}
In particular the rough field satisfies
\(\norm{\Lambda_\lambda}_{C^t}\le C_t\ell^{-t}\) for every \(t>0\) (it is
\(G_h\) times a bounded-overlap sum of translates of \(\phi(\cdot/\ell)\),
and \(h\ge\ell\)).

\subsection{Two-point reduction and the component principle}
\begin{lemma}[Two-point reduction]\label{lem:two-point}
Let \(\PP_n,\mathbb Q_n\) be the laws of \((X_i,A_i,Y_i)_{i\le n}\) obtained
by drawing a law from a prior supported in \(\cP\) and then \(n\) iid
triples from it, and suppose \(\tau_P(x_0)=\theta_1\) for every law in the
support of the first prior and \(\tau_P(x_0)=\theta_0\) for every law in
the support of the second, with \(\theta_1-\theta_0=\Delta_{\rm sep}>0\).
Then for every estimator
\[
 \sup_{P\in\cP}\E_P\abs{\widehat\tau(x_0)-\tau_P(x_0)}
 \ge\frac{\Delta_{\rm sep}}2\bigl(1-\TV(\PP_n,\mathbb Q_n)\bigr).
\]
\end{lemma}
\begin{proof}
The supremum dominates both mixture risks; with
\(g=\abs{\widehat\tau(x_0)-\theta_1}\), \(g'=\abs{\widehat\tau(x_0)-\theta_0}\),
\(g+g'\ge\Delta_{\rm sep}\), so
\(\E_{\PP_n}g+\E_{\mathbb Q_n}g'\ge\int(g+g')\,d(\PP_n\wedge\mathbb Q_n)\ge\Delta_{\rm sep}(1-\TV)\).
\end{proof}

For probability vectors \(P',Q'\) on a finite set write
\(H^2(P',Q')=\sum_w(\sqrt{P'(w)}-\sqrt{Q'(w)})^2\in[0,2]\).

\begin{lemma}[Hellinger facts]\label{lem:hellinger}
\begin{enumerate}[label=(\alph*),leftmargin=*]
\item \(\TV(P',Q')\le H(P',Q')\).
\item \(H^2(\bigotimes_CP'_C,\bigotimes_CQ'_C)\le\sum_CH^2(P'_C,Q'_C)\).
\item If \(P'(w)\ge c>0\) for all \(w\), then
\(H^2(P',Q')\le c^{-1}\sum_w\abs{P'(w)-Q'(w)}^2\).
\end{enumerate}
\end{lemma}
\begin{proof}
(a) \(2\TV=\sum\abs{\sqrt{P'}-\sqrt{Q'}}(\sqrt{P'}+\sqrt{Q'})
\le H\{\sum(\sqrt{P'}+\sqrt{Q'})^2\}^{1/2}\le2H\).
(b) \(1-H^2/2=\sum\sqrt{P'Q'}\) is multiplicative and
\(\prod_C(1-h_C)\ge1-\sum_Ch_C\) for \(h_C\in[0,1]\).
(c) \((\sqrt{P'}-\sqrt{Q'})^2=(P'-Q')^2/(\sqrt{P'}+\sqrt{Q'})^2\le(P'-Q')^2/P'\).
\end{proof}

\begin{lemma}[Component principle]\label{lem:components}
Let \(\PP_n,\mathbb Q_n\) have the same \(X^n\)-marginal, and let
\(\mathcal E_n\) be an event depending on \(X^n\) only.  Suppose that for
every \(x^n\notin\mathcal E_n\) there is a partition of \(\{1,\dots,n\}\)
into \emph{components} \(C\) such that the conditional laws of
\(w^n=(R_i,T_i)_{i\le n}\) given \(X^n=x^n\) factor as
\(\PP_n(w^n\mid x^n)=\prod_C P_C(w_C)\), \(\mathbb Q_n(w^n\mid x^n)=\prod_CQ_C(w_C)\).
Then, with \(H^2_C=H^2(P_C,Q_C)\),
\[
 \TV(\PP_n,\mathbb Q_n)\le\PP(\mathcal E_n)+\Bigl\{\E\Bigl[\one_{\mathcal E_n^c}\sum_CH^2_C\Bigr]\Bigr\}^{1/2}.
\]
\end{lemma}
\begin{proof}
Since the \(X^n\)-marginals coincide,
\(\TV(\PP_n,\mathbb Q_n)\le\E\,\TV(\PP_n(\cdot\mid X^n),\mathbb Q_n(\cdot\mid X^n))\).
Bound the conditional total variation by one on \(\mathcal E_n\) and, on
\(\mathcal E_n^c\), by \(\{\sum_CH^2_C\}^{1/2}\) (Lemma~\ref{lem:hellinger}(a),(b));
then apply Jensen's inequality to the square root.
\end{proof}

\begin{lemma}[Design tilt]\label{lem:tilt}
Let \(\nu_P,\nu_Q\) be laws of a triple \((H,\lambda,U)\) with the same
\((H,\lambda)\)-marginal, and let \((H,\lambda)\mapsto\widetilde f_{H,\lambda}\)
be a measurable family of nonnegative functions on \([0,1]^d\) with
\(Z(H,\lambda)=\int\widetilde f_{H,\lambda}\in[z_0,z_1]\), \(0<z_0\le z_1<\infty\);
put \(f_{H,\lambda}=\widetilde f_{H,\lambda}/Z(H,\lambda)\).  Define the tilted
priors
\begin{equation}
 d\Pi_{S,n}=\frac{Z(H,\lambda)^n}{\mathcal Z_n}\,d\nu_S,\qquad
 \mathcal Z_n=\int Z^n\,d\nu_P=\int Z^n\,d\nu_Q,\qquad S\in\{P,Q\},
 \label{eq:z-tilt}
\end{equation}
and let \(\PP_n,\mathbb Q_n\) be the mixtures obtained by drawing
\((H,\lambda,U)\sim\Pi_{S,n}\), then \(X_1,\dots,X_n\) iid with density
\(f_{H,\lambda}\), then \(W_i=(R_i,T_i)\) from site likelihoods
\(L^S_i(w_i)\) relative to fair signs, depending on \((x_i,H,\lambda,U)\).
Then the \((H,\lambda)\)-marginals of \(\Pi_{P,n}\) and \(\Pi_{Q,n}\)
coincide, hence so do the laws of \(f_{H,\lambda}\) and the
\(X^n\)-marginals of \(\PP_n\) and \(\mathbb Q_n\); and the joint density of
\((X^n,W^n)\) with respect to Lebesgue measure times counting measure is
\begin{equation}
 \overline p_S(x^n,w^n)=\frac{4^{-n}}{\mathcal Z_n}\int\prod_{i=1}^n\widetilde f_{H,\lambda}(x_i)\prod_{i=1}^nL^S_i(w_i)\,d\nu_S(H,\lambda,U):
 \label{eq:z-cancel}
\end{equation}
the design normaliser \(Z^n\) cancels against the tilt and the untilted
prior \(\nu_S\) reappears.
\end{lemma}
\begin{proof}
\(Z\) is a function of \((H,\lambda)\) alone, so \(\mathcal Z_n\) is the same
on both sides and the density \(Z^n/\mathcal Z_n\) of \(\Pi_{S,n}\) with
respect to \(\nu_S\) tilts only the common \((H,\lambda)\)-marginal.  The
joint density is
\[
 \int\prod_{i=1}^nf_{H,\lambda}(x_i)\prod_{i=1}^n4^{-1}L^S_i(w_i)\,d\Pi_{S,n},
 \qquad
 \prod_{i=1}^nf_{H,\lambda}(x_i)=Z(H,\lambda)^{-n}\prod_{i=1}^n\widetilde f_{H,\lambda}(x_i),
\]
and substituting \eqref{eq:z-tilt} gives \eqref{eq:z-cancel}.
\end{proof}

\section{Scales}\label{sec:scales}

This section fixes the scales of Parts~B and~C.  (Part~A uses only
\(\Delta\), \(h\) and \(\ell\), which it defines directly in
\eqref{A:eq:amplitudes} with \(\varepsilon=0\); the quantities \(r\), \(t\),
\(Q\) and the ledgers below are not used there.)  Fix \(\varepsilon>0\)
small and put
\begin{equation}
 z=\rho+\varepsilon,\quad \Delta=n^{-z},\quad h=\Delta^{1/\gamma},\quad
 r=n^{-\xi},\quad \xi=\frac1d+c\varepsilon,\quad 0<c<\frac{D_1}{dA_0},
 \quad \Lambda=1+\log n .
 \label{eq:scales}
\end{equation}
Define the fine ratio \(t\), the rough scale \(\ell\), and the taper scale
by
\begin{equation}
 t^{2s_1}=\Delta r^{-A_0},\qquad \ell=tr,\qquad t_n=t\Lambda^B,
 \label{eq:fine}
\end{equation}
with a fixed exponent \(B\) chosen in each Part (Lemma~\ref{B:lem:wiener-increment},
Lemma~\ref{C:lem:wiener}).  Note that \(\ell^{2s_1}=\Delta r^{2s_1-A_0}\),
which reduces to \(\ell^{2s_1}=\Delta\) when \(A_0=2s_1\); this is the
relation between \(\ell\) and \(\Delta\) used in Part~A.  Let
\begin{equation}
 Q=Q_\varepsilon=\max\Bigl\{2,\,1+\Bigl\lfloor\frac{[1-zd/\gamma]_+}{cd\varepsilon}\Bigr\rfloor\Bigr\},
 \label{eq:Q}
\end{equation}
the finite matching order of Parts~B and~C; it depends on \(\varepsilon\)
but not on \(n\).  Constants denoted by \(C_Q\) may depend on the fixed
quantities \(d,\alpha,\beta,\gamma,\varepsilon,c\), the Hölder radii, and
the fixed smooth cutoffs and feature system, but never on \(n\) or on
dyadic shell indices.

\begin{lemma}[Ordering]\label{lem:order}
In the regime of Theorem~\ref{thm:main}: \(\gamma>A_0\), \(\rho>A_0/d\),
\(d\rho<A_0+2s_1\), \(\gamma(d+4s_1)>dA_0\), and \(\rho<1\).  Moreover, for
all sufficiently small \(\varepsilon>0\): \(t\to0\), \(\ell<r<h\),
\(nr^d=n^{-cd\varepsilon}\to0\), and \(t/t_n=\Lambda^{-B}\to0\).
\end{lemma}

\begin{proof}
The regime reads \(\gamma d>A_0(2\gamma+d)\), so \(d(\gamma-A_0)>2\gamma
A_0>0\) and \(\gamma>A_0\); also \(\gamma(d+4s_1)>\gamma d>dA_0\).
\(\rho>A_0/d\) iff \(d(A_0+2s_1)>A_0D_1\) iff \(2ds_1>4s_1A_0+2ds_1A_0/\gamma\)
iff \(d>A_0(2+d/\gamma)\).  \(d\rho<A_0+2s_1\) because \(D_1>d\), and
\(\rho<1\) because \(A_0+2s_1<d+4s_1<D_1\) (the regime gives \(d>2A_0\)).
From \eqref{eq:fine},
\[
 \log_nt=-\frac{z-\xi A_0}{2s_1}
 =-\frac{(\rho-A_0/d)+\varepsilon(1-cA_0)}{2s_1}<0
\]
for small \(\varepsilon\), so \(t\to0\) and \(\ell<r\).  \(r<h\) iff
\(\xi>z/\gamma\), which for small \(\varepsilon\) is \(\gamma>d\rho\), i.e.\
\(\gamma D_1>d(A_0+2s_1)\), i.e.\ \(\gamma(d+4s_1)>dA_0\).
\end{proof}

\begin{lemma}[Rate ledgers]\label{lem:ledgers}
\begin{align}
 nh^dt^d\Delta^2&=n^{-\varepsilon(D_1-cdA_0)/(2s_1)}\to0,
 \label{eq:singleton}\\
 n^2h^d\ell^d\Delta^2&=n^{-\varepsilon\{D_1-cd(A_0-2s_1)\}/(2s_1)}\to0,
 \label{eq:collision}\\
 n^{Q+1}h^dr^{dQ}&=n^{1-zd/\gamma-cd\varepsilon Q}\to0.
 \label{eq:high}
\end{align}
Fixed powers of \(\Lambda\) multiplying the left sides do not change the
conclusions.
\end{lemma}

\begin{proof}
By \eqref{eq:fine}, \(d\log_nt=-d(z-\xi A_0)/(2s_1)\), so
\[
 \log_n(nh^dt^d\Delta^2)=1-\frac{zd}\gamma-\frac{d(z-\xi A_0)}{2s_1}-2z
 =1+\frac{A_0}{2s_1}+\frac{cd\varepsilon A_0}{2s_1}-\frac{zD_1}{2s_1}.
\]
Since \(\rho D_1=2s_1+A_0\), substituting \(z=\rho+\varepsilon\) gives
\eqref{eq:singleton}, negative because \(c<D_1/(dA_0)\).  Multiply by
\(nr^d=n^{-cd\varepsilon}\) and use \(t^dr^d=\ell^d\) for
\eqref{eq:collision}.  For \eqref{eq:high}: if \(zd/\gamma\ge1\) the
exponent is negative for \(Q\ge2\); otherwise
\(Q>[1-zd/\gamma]/(cd\varepsilon)\) by the choice of \(Q\).
\end{proof}
Nothing in this section distinguishes cases~(B) and~(C): only \(A_0\), \(s_1\),
\(\gamma\), \(d\) enter.

\section{Interpolation, taper, moment simplex, and shells}\label{sec:torus}

\subsection{An explicit finite interpolation system}

Let
\[
 \iota(u)=(\cos 2\pi u_1,\sin 2\pi u_1,\ldots,
           \cos 2\pi u_d,\sin 2\pi u_d)\in\R^{2d}.
\]
Write
\[
 d_{\circ}(u,v)=\norm{\iota(u)-\iota(v)}
 =2\left\{\sum_{a=1}^d\sin^2\!\bigl(\pi(u_a-v_a)\bigr)\right\}^{1/2}.
\]
Choose \(r_0>0\) small and a smooth nondecreasing truncation
\(\kappa_\delta:[0,\infty)\to[0,1]\) such that
\(\kappa_\delta(r)=r\) for \(r\le r_0\),
\(\kappa_\delta(r)\) is a fixed positive constant for \(r\ge2r_0\), and
\(\kappa_\delta(r)>0\) for \(r>0\).  Put
\(\delta(u,v)=\kappa_\delta\{d_{\circ}(u,v)\}\).  Thus \(\delta\) depends
only on the torus difference, is smooth off the diagonal, is comparable to
\(d_{\circ}\wedge1\), and is constant away from a fixed diagonal
neighborhood.  For a tuple \(T=(u_1,\ldots,u_q)\) of distinct points, define
\begin{equation}
 P_i(T)=\prod_{j\ne i}\delta(u_i,u_j),
 \qquad
 P_{\min}(T)=\min_iP_i(T).
 \label{eq:product-distance}
\end{equation}
Let \(v_Q(u)\) contain all coordinate monomials in \(\iota(u)\) of total
degree at most \(Q-1\), including repetitions and the constant.  Its dimension depends only on
\((d,Q)\).

\begin{lemma}[Uniform Lagrange vectors]
\label{lem:lagrange}
For \(q\le Q\) and any distinct tuple \(T=(u_1,\ldots,u_q)\),
there are vectors \(c_i(T)\) satisfying
\begin{equation}
 v_Q(u_j)^\top c_i(T)=\one\{i=j\},
 \qquad
 \norm{c_i(T)}\le\frac{C_Q}{P_i(T)}.
 \label{eq:lagrange}
\end{equation}
The construction is equivariant under permutations of the tuple.
\end{lemma}

\begin{proof}
For \(i\ne j\), put
\[
 L_{ij}(u;T)
 =\frac{\ip{\iota(u)-\iota(u_j)}{\iota(u_i)-\iota(u_j)}}
        {\norm{\iota(u_i)-\iota(u_j)}^2}.
\]
Then \(L_{ij}(u_j;T)=0\), \(L_{ij}(u_i;T)=1\), and it is a linear
combination of the constant and the coordinates of \(\iota(u)\).  Set
\[
 L_i(u;T)=\prod_{j\ne i}L_{ij}(u;T).
\]
This lies in the span of \(v_Q\) and satisfies
\(L_i(u_j;T)=\one\{i=j\}\).  Each factor has coefficient norm at most
\(C/\delta(u_i,u_j)\); multiplication in the fixed finite-dimensional
feature space therefore gives the second assertion.  The formula is
permutation equivariant.
\end{proof}


Each coordinate of \(c_i(T)\), being a coefficient of \(\prod_{j\ne i}L_{ij}\),
is a finite sum of products containing exactly one of the
\emph{elementary multipliers}
\begin{equation}
 \frac{\{\iota(u_i)-\iota(u_j)\}_a}{\norm{\iota(u_i)-\iota(u_j)}^2},
 \qquad
 \frac{\ip{\iota(u_j)}{\iota(u_i)-\iota(u_j)}}{\norm{\iota(u_i)-\iota(u_j)}^2}
 \label{eq:elementary-lagrange-multipliers}
\end{equation}
for every edge \(\{i,j\}\) incident to \(i\), each of inverse-distance
order one.

\subsection{Taper and close-tuple volume}

With \(\Lambda=1+\log n\) and the fine ratio \(t\) of \eqref{eq:fine},
the taper scale is
\begin{equation}
 t_n=t\Lambda^B,
 \label{eq:taper-scale}
\end{equation}
where the fixed exponent \(B>0\) is chosen in each Part.
For large \(n\), \(t_n<1/4\).  Let
\(\vartheta\in C^\infty([0,\infty);[0,1])\) be zero on \([0,1]\) and one
on \([2,\infty)\), and define the symmetric taper
\begin{equation}
 \chi_Q(T)=\prod_{i=1}^Q
 \vartheta\left(\frac{P_i(T)}{t_n}\right).
 \label{eq:taper}
\end{equation}
The product is extended by zero on the diagonals.  Thus \(\chi_Q=0\) if
\(P_{\min}\le t_n\), and \(\chi_Q=1\) if every \(P_i\ge2t_n\).

\begin{lemma}[Close-tuple volume]
\label{lem:bad-volume}
For each fixed \(2\le q\le Q\) and \(0<\varrho<1/4\), the normalized
\(dq\)-dimensional volume of
\[
 \cB_q(\varrho)=\{T\in(\Torus^d)^q:P_{\min}(T)\le C\varrho\}
\]
satisfies
\begin{equation}
 \abs{\cB_q(\varrho)}
 \le C_Q\varrho^d\{1+\log(1/\varrho)\}^{q-2}.
 \label{eq:bad-volume}
\end{equation}
Consequently, the set on which \(\chi_Q\ne1\) has volume at most
\begin{equation}
 C_Qt^d\Lambda^{M_Q},\qquad M_Q=Bd+Q-2 .
 \label{eq:taper-volume}
\end{equation}
\end{lemma}

\begin{proof}
Fix \(i\) and \(u_i\).  Up to a fixed multiplicative constant, integration
in polar coordinates around \(u_i\) reduces the relevant part to
\[
 \int_{[0,1]^{q-1}}
 \one\left\{\prod_{j=1}^{q-1}r_j\le C\varrho\right\}
 \prod_{j=1}^{q-1}r_j^{d-1}\,dr_j.
\]
With \(w_j=r_j^d\), this is a constant times the volume of
\(\{w\in[0,1]^{q-1}:\prod_jw_j\le C^d\varrho^d\}\), which equals
\[
 O_Q\left(\varrho^d\{1+\log(1/\varrho)\}^{q-2}\right)
\]
by one integration, or induction on \(q\).  Sum over \(i\).  Substituting
\(\varrho=2t_n=2t\Lambda^B\) and using \(\log(1/t)=O(\log n)\) gives
\eqref{eq:taper-volume}.
\end{proof}

\subsection{Wiener notation and an interior coefficient law}

For \(m\ge1\), let \(\cW_m\) be the real Wiener algebra on \(\Torus^m\):
\[
 \norm{g}_{\cW_m}
 =\sum_{\lambda\in\mathbb Z^m}\abs{\widehat g(\lambda)}.
\]
For finite-dimensional vector-valued functions, sum the coordinate norms.
The Wiener algebra is a Banach algebra and
\begin{equation}
 \norm{gh}_{\cW_m}\le\norm g_{\cW_m}\norm h_{\cW_m},
 \qquad
 \norm g_\infty\le\norm g_{\cW_m}.
 \label{eq:wiener-algebra}
\end{equation}

Let \(M\) be the dimension of \(v_Q\), and let
\(\Psi_{4Q}(U)\in\R^{N_Q}\) contain every nonconstant monomial in
\(U\in\R^M\) of total degree at most \(4Q\).

\begin{lemma}[Interior finite moment simplex]
\label{lem:moment-simplex}
There is a finite law \(\Pi_0\), supported in an arbitrarily small open ball
\(\mathbb B_0\subset\R^M\), and a number \(\delta_Q>0\) such that every
vector \(v\in\R^{N_Q}\) with \(\norm v\le\delta_Q\) is realized as
\begin{equation}
 \E_{\Pi_v}\Psi_{4Q}(U)
 =\E_{\Pi_0}\Psi_{4Q}(U)+v
 \label{eq:moment-perturbation}
\end{equation}
by a probability law \(\Pi_v\) supported on the same finite set as
\(\Pi_0\).
\end{lemma}

\begin{proof}
The functions \(1\) and the coordinates of \(\Psi_{4Q}\) are linearly
independent polynomials on every open ball.  Hence one can choose
\(N_Q+1\) points \(u^0,\ldots,u^{N_Q}\in\mathbb B_0\) for which the vectors
\((1,\Psi_{4Q}(u^j))\) are affinely independent; otherwise a nonzero
polynomial of degree at most \(4Q\) would vanish on an open set.  Give all
points strictly positive mass and call the resulting law \(\Pi_0\).  Its
moment vector is an interior point of the resulting \(N_Q\)-dimensional
simplex, so a fixed Euclidean ball around it is contained in that simplex.
This gives \eqref{eq:moment-perturbation}.
\end{proof}


\subsection{Dyadic configuration shells}

We fix a dyadic partition that is also used in Section~\ref{sec:wiener-config}.  Choose an
integer \(m_0\) so large that \(2^{-m_0+3}<r_0\), every pair with
\(\delta(u,v)\le 2^{-m_0+3}\) has a unique difference lift in
\((-1/4,1/4)^d\), and all rescaled fine-shell lifts lie in a fixed
compact set contained in \((-2^{m_0-2},2^{m_0-2})^d\).  Let \(\rho\in C_c^\infty((1/4,4))\) be nonnegative and
satisfy
\[
 \sum_{m\in\mathbb Z}\rho(2^mr)=1,\qquad r>0.
\]
For \(m\ge m_0\), set \(\rho_m(r)=\rho(2^mr)\), and put
\[
 \rho_{\mathrm c}(r)=\sum_{m<m_0}\rho(2^mr).
\]
The latter is smooth on \([0,1]\), vanishes in a neighborhood of zero, and
together with the \(\rho_m\)'s forms a nonnegative partition of unity on
\((0,1]\).  Write
\[
 \mathbb M=\{\mathrm c\}\cup\{m:m\ge m_0\},
 \qquad
 \mu(\mathrm c)=0,\quad \mu(m)=m.
\]
There are fixed constants \(0<c_\rho<C_\rho<\infty\) such that
\begin{equation}
 c_\rho 2^{-\mu(m)}\le r\le C_\rho2^{-\mu(m)}
 \quad\text{on }\supp\rho_m,
 \qquad m\in\mathbb M.
 \label{eq:dyadic-shell-comparability}
\end{equation}
For the edge set
\[
 \mathfrak E_Q=\{\{i,j\}:1\le i<j\le Q\},
 \qquad E_Q:=\abs{\mathfrak E_Q}=\binom Q2,
\]
and a shell assignment \(\bm m=(m_e)_{e\in\mathfrak E_Q}\in
\mathbb M^{\mathfrak E_Q}\), define
\begin{equation}
 \rho_{\bm m}(T)
 =\prod_{e=\{i,j\}\in\mathfrak E_Q}
   \rho_{m_e}\{\delta(u_i,u_j)\}.
 \label{eq:configuration-shell-partition}
\end{equation}
Off the diagonals, \(\sum_{\bm m}\rho_{\bm m}=1\), with a locally finite
sum.


For a shell assignment \(\bm m\) and a vertex \(i\) write
\begin{equation}
 M_i(\bm m)=\sum_{e\in\mathfrak E_Q:\,e\ni i}\mu(m_e).
 \label{eq:vertex-shell-sum}
\end{equation}

\begin{lemma}[Admissible shells]\label{lem:admissible-shells}
If \(\chi_Q\rho_{\bm m}\not\equiv0\), then \(M_i(\bm m)\le\log_2(1/t_n)+C_Q\)
and \(2^{M_i(\bm m)}\le C_Qt_n^{-1}\) for every \(1\le i\le Q\); and the number
of shell assignments with \(\chi_Q\rho_{\bm m}\not\equiv0\) is at most
\(C_Q\{1+\log(1/t_n)\}^{E_Q}\).
\end{lemma}
\begin{proof}
If \(\chi_Q\rho_{\bm m}\) is not identically zero, there is a point of its
support at which every \(P_i>t_n\).  By
\eqref{eq:dyadic-shell-comparability}, throughout that shell
\(P_i(T)\le C_Q2^{-M_i(\bm m)}\), whence the first claim.  In particular
every fine edge index is at most \(\log_2(1/t_n)+C_Q\); since there are
\(E_Q\) edges, the finitely many coarse choices and the bounded overlap
of the dyadic partition give the count.
\end{proof}

\subsection{The polarization identity}
For elements \(f_1,\dots,f_Q\) of a commutative algebra with unit, write
\(f_1\od\cdots\od f_Q=(Q!)^{-1}\sum_{\sigma}f_{\sigma(1)}\otimes\cdots\otimes f_{\sigma(Q)}\)
for the symmetrised tensor product.

\begin{lemma}[Real polarization]\label{lem:polarization-identity}
Let \(g_1,\dots,g_j\) (\(1\le j\le Q\)) be elements of a commutative Banach
algebra with unit and put \(g_\epsilon=j^{-1}\sum_{i=1}^j\epsilon_ig_i\),
\(\epsilon\in\{\pm1\}^j\).  Fix distinct nodes \(t_0,\dots,t_Q\in[-1,1]\) and
let \(\lambda_{jr}\) solve the Vandermonde system
\begin{equation}
 \sum_{r=0}^Q\lambda_{jr}t_r^k=\frac{\one\{k=j\}}{\binom Qj},\qquad 0\le k\le Q .
 \label{eq:vandermonde}
\end{equation}
Then for every \(\theta\ne0\),
\begin{equation}
 g_1\od\cdots\od g_j\od1^{\od(Q-j)}
 =\frac{j^j}{2^jj!\,\theta^j}\sum_{\epsilon\in\{\pm1\}^j}\Bigl(\prod_{i=1}^j\epsilon_i\Bigr)
 \sum_{r=0}^{Q}\lambda_{jr}(1+\theta t_rg_\epsilon)^{\otimes Q}.
 \label{eq:polarization-identity}
\end{equation}
The identity is algebraic: it only uses commutativity, and the tensor
powers on the right are products over distinct slots.
\end{lemma}
\begin{proof}
The real polarization identity gives
\begin{equation}
 g_1\od\cdots\od g_j
 =\frac1{2^jj!}\sum_{\epsilon\in\{\pm1\}^j}\Bigl(\prod_i\epsilon_i\Bigr)
 \Bigl(\sum_i\epsilon_ig_i\Bigr)^{\otimes j}
 =\frac{j^j}{2^jj!}\sum_\epsilon\Bigl(\prod_i\epsilon_i\Bigr)g_\epsilon^{\otimes j},
 \label{eq:real-polarization}
\end{equation}
since expanding \((\sum_i\epsilon_ig_i)^{\otimes j}\) and summing against
\(\prod_i\epsilon_i\) retains exactly the terms in which every \(g_i\)
occurs once.  Expanding \((1+\theta t_rg)^{\otimes Q}=\sum_{k=0}^Q(\theta t_r)^k\binom Qk\,g^{\otimes k}\od1^{\od(Q-k)}\)
and using \eqref{eq:vandermonde},
\begin{equation}
 g^{\otimes j}\od1^{\od(Q-j)}=\theta^{-j}\sum_{r=0}^Q\lambda_{jr}(1+\theta t_rg)^{\otimes Q}.
 \label{eq:degree-extraction}
\end{equation}
Combine \eqref{eq:real-polarization} with \eqref{eq:degree-extraction}
for \(g=g_\epsilon\).
\end{proof}

\section{Configuration-space Wiener estimates}\label{sec:wiener-config}

The purpose of this section is to justify the shellwise estimate used in
\cref{B:lem:wiener-increment} in the full \(dQ\)-variable Wiener algebra.  The
argument treats all pair differences simultaneously; no independence of the
edge variables is assumed.

\begin{lemma}[Uniform dyadic pullback on a finite configuration graph]
\label{lem:configuration-pullback}
Let \(V=\{1,\ldots,Q\}\), let \(E\) be any fixed set of unordered pairs from
\(V\), and fix an integer \(m_0\).  For every \(e=\{i,j\}\in E\), choose an
integer \(m_e\ge m_0\) and put \(L_e=2^{m_e}\).  Let
\(K\Subset(-2^{m_0-2},2^{m_0-2})^d\), and let
\[
 H_{\bm m}\in C^\infty\bigl((\Torus^d)^Q\times(\R^d)^E\bigr)
\]
be periodic in the first \(dQ\) variables and supported, in each Euclidean
variable \(z_e\), inside \(K\).  Suppose that for an integer
\[
 R>\frac{dQ+d\abs E}{2}
\]
one has the uniform derivative bound
\begin{equation}
 \sup_{\bm m}
 \sum_{\abs\alpha\le2R}
 \norm{\partial^\alpha H_{\bm m}}_{L^\infty}
 \le C_H.
 \label{eq:profile-derivative-bound}
\end{equation}
Periodize \(H_{\bm m}\) in the variable \(z_e\) with period \(L_e\), and
define on \((\Torus^d)^Q\)
\begin{equation}
 (\mathfrak P_{\bm m}H_{\bm m})(u_1,\ldots,u_Q)
 =H_{\bm m}^{\mathrm{per}}
 \left(u_1,\ldots,u_Q,
       \{2^{m_e}(u_i-u_j)\}_{e=\{i,j\}\in E}
 \right).
 \label{eq:dyadic-pullback}
\end{equation}
The expression is independent of the chosen representatives of the torus
points.  Moreover,
\begin{equation}
 \norm{\mathfrak P_{\bm m}H_{\bm m}}_{\cW_{dQ}}
 \le C_{d,Q,E,R,K}\,C_H,
 \label{eq:dyadic-pullback-bound}
\end{equation}
uniformly over all dyadic indices \(\bm m\).  The conclusion remains true
for a finite sum of such profiles and for a family depending on auxiliary
parameters ranging over a fixed compact set, provided
\eqref{eq:profile-derivative-bound} is uniform in those parameters.
\end{lemma}

\begin{proof}
Because \(K\Subset(-L_e/2,L_e/2)^d\) for every \(e\), periodization creates
no overlap on the support of \(H_{\bm m}\).  Write the Fourier series of the
periodized profile as
\begin{equation}
 H_{\bm m}^{\mathrm{per}}(u,z)
 =\sum_{k\in\mathbb Z^{dQ}}
   \sum_{\nu\in(\mathbb Z^d)^E}
   b_{\bm m}(k,\nu)
   e^{2\pi i k\cdot u}
   \prod_{e\in E}e^{2\pi i\nu_e\cdot z_e/L_e}.
 \label{eq:profile-fourier-series}
\end{equation}
Put \(p=dQ+d\abs E\) and
\(V_{\bm m}=\prod_{e\in E}L_e^d\).  The Fourier coefficient is normalized by
\(V_{\bm m}^{-1}\).  Apply \((1-\Delta_{u,z})^R\), integrate by parts, and
use the fixed compact support in the \(z\)-variables.  The derivative bound
\eqref{eq:profile-derivative-bound} gives
\begin{equation}
 \abs{b_{\bm m}(k,\nu)}
 \le C_{R,K}C_H V_{\bm m}^{-1}
 \left\{1+\abs k^2+
  \sum_{e\in E}\frac{\abs{\nu_e}^2}{L_e^2}\right\}^{-R}.
 \label{eq:anisotropic-fourier-coefficient}
\end{equation}
We now sum this estimate without losing a scale factor.  For fixed \(k\),
the points \((\nu_e/L_e)_{e\in E}\) form an anisotropic rectangular grid in
\(\R^{d\abs E}\) whose cell volume is \(V_{\bm m}^{-1}\).  If
\(y\) lies in the grid cell adjacent to \((\nu_e/L_e)_e\), then, because
all \(L_e\ge1\), the two quantities are uniformly comparable:
\[
 C_E^{-1}\left(1+\abs k^2+\abs y^2\right)
 \le 1+\abs k^2+\sum_e\abs{\nu_e/L_e}^2
 \le C_E\left(1+\abs k^2+\abs y^2\right).
\]
The left inequality is the direction needed below.  Comparison with the
integral over the disjoint grid cells gives
\begin{align}
 &V_{\bm m}^{-1}
 \sum_{\nu\in(\mathbb Z^d)^E}
 \left\{1+\abs k^2+
  \sum_e\frac{\abs{\nu_e}^2}{L_e^2}\right\}^{-R}
 \nonumber\\
 &\hspace{2cm}\le
 C_E\int_{\R^{d\abs E}}
       (1+\abs k^2+\abs y^2)^{-R}\,dy.
 \label{eq:anisotropic-grid-sum}
\end{align}
The integral on the right is bounded by
\(C(1+\abs k^2)^{-R+d\abs E/2}\).  Summing over
\(k\in\mathbb Z^{dQ}\) is finite precisely under the stated condition
\(R>(dQ+d\abs E)/2\).  Equations
\eqref{eq:anisotropic-fourier-coefficient}--
\eqref{eq:anisotropic-grid-sum} therefore imply
\begin{equation}
 \sum_{k,\nu}\abs{b_{\bm m}(k,\nu)}
 \le C_{d,Q,E,R,K}C_H,
 \label{eq:profile-l1-fourier}
\end{equation}
uniformly in all \(m_e\).

Substituting \eqref{eq:dyadic-pullback} into one term of
\eqref{eq:profile-fourier-series} produces
\[
 e^{2\pi i k\cdot u}
 \prod_{e=\{i,j\}\in E}
 e^{2\pi i\nu_e\cdot(u_i-u_j)},
\]
because \(2^{m_e}/L_e=1\).  This is one character of
\((\Torus^d)^Q\).  Edges sharing a vertex merely add their integer
frequencies in that vertex coordinate, and different dyadic scales have
already disappeared from the character after the matched period choice
\(L_e=2^{m_e}\).  Every pulled-back Fourier term thus has Wiener norm one.
Distinct \((k,\nu)\) may produce the same character, but the triangle
inequality and \eqref{eq:profile-l1-fourier} give
\eqref{eq:dyadic-pullback-bound}.

Finally, changing a representative \(u_i\) by an integer changes the
argument \(2^{m_e}(u_i-u_j)\) by an integer multiple of the period
\(L_e=2^{m_e}\).  Thus \eqref{eq:dyadic-pullback} is well defined.  Finite
sums and compact auxiliary parameter families obey the same estimates.
\end{proof}

\begin{corollary}[Shell-localized Lagrange and product-taper multipliers]
\label{cor:configuration-multipliers}
Use the dyadic partition from
\eqref{eq:configuration-shell-partition}.  Fix a shell assignment
\(\bm m\in\mathbb M^{\mathfrak E_Q}\), and for each edge
\(e=\{i,j\}\) let \(a_e\) be a nonnegative integer bounded by a fixed
constant \(A\).  Let \(M_{\bm m}(T)\) be any one of the finitely many scalar
products obtained from
\begin{enumerate}
\item the shell cutoffs \(\rho_{m_e}\{\delta(u_i,u_j)\}\);
\item the product taper \(\chi_Q(T)\);
\item \(a_e\) factors, at edge \(e\), chosen from the elementary
      Lagrange coefficient multipliers in
      \eqref{eq:elementary-lagrange-multipliers}; and
\item a fixed number of functions from a fixed bounded subset of
      \(C^\infty\) on the torus coordinates.
\end{enumerate}
Then
\begin{equation}
 \norm{M_{\bm m}}_{\cW_{dQ}}
 \le C_{d,Q,A}\,
      2^{\sum_{e\in\mathfrak E_Q}a_e\mu(m_e)}.
 \label{eq:configuration-multiplier-bound}
\end{equation}
The constant is independent of \(n\), \(t_n\), and all shell indices.
\end{corollary}

\begin{proof}
Call an edge fine when \(m_e\ne\mathrm c\).  On a coarse edge,
\(\rho_{\mathrm c}(\delta_e)\) is supported a fixed positive distance from
the diagonal.  Every multiplier in
\eqref{eq:elementary-lagrange-multipliers} is therefore a fixed smooth
function there, with a Wiener norm bounded independently of all scales.  We
absorb every coarse-edge factor into the periodic \(u\)-dependence of a
profile.

Consider a fine edge \(e=\{i,j\}\), write \(m=m_e\), and use a local lift
\(w=u_i-u_j\in(-1/4,1/4)^d\).  Put \(z=2^mw\).  By
\eqref{eq:dyadic-shell-comparability} and the local comparability of
\(\delta\) with the chordal distance, the shell cutoff restricts \(z\) to a
fixed compact annulus
\[
 K_0\Subset\R^d\setminus\{0\},
\]
independently of \(m\).  The families
\begin{align*}
 q_m(u_j,z)
 &=2^m\delta(u_j+2^{-m}z,u_j),\\
 B_{m,a}(u_j,z)
 &=2^{-m}
   \frac{\{\iota(u_j+2^{-m}z)-\iota(u_j)\}_a}
        {\norm{\iota(u_j+2^{-m}z)-\iota(u_j)}^2},\\
 B_{m,0}(u_j,z)
 &=2^{-m}
   \frac{\ip{\iota(u_j)}
   {\iota(u_j+2^{-m}z)-\iota(u_j)}}
        {\norm{\iota(u_j+2^{-m}z)-\iota(u_j)}^2}
\end{align*}
are uniformly \(C^N\) on
\(\Torus^d\times K_0\), for every fixed \(N\).  To verify this, multiply
numerator and denominator by the appropriate powers of \(2^m\) and use the
Taylor formula for the smooth embedding \(\iota\).  Its differential has
full rank, so the rescaled denominator is bounded above and below on
\(K_0\); differentiating the resulting quotient gives bounds independent of
\(m\).  The same argument applies to \(q_m\), and
\(0<c_Q\le q_m\le C_Q\) on the shell.  Thus each order-one multiplier is
\(2^m\) times a uniformly smooth function of \((u_j,z)\), while the shell
cutoff itself is the uniformly smooth function \(\rho\{q_m(u_j,z)\}\).
A product containing \(a_e\) order-one factors at edge \(e\) therefore
contributes the explicit scale \(2^{a_em_e}\) and otherwise has a uniformly
smooth rescaled profile.

It remains to verify the product taper.  For a vertex \(i\), put
\[
 M_i(\bm m)=\sum_{e\ni i}\mu(m_e),
 \qquad
 \lambda_i=t_n^{-1}2^{-M_i(\bm m)}.
\]
On the fixed shell profile,
\begin{equation}
 \frac{P_i(T)}{t_n}
 =\lambda_i
  \prod_{e\ni i}q_{e,m_e}(u,z_e),
 \label{eq:rescaled-product-taper}
\end{equation}
where, for a coarse edge, \(q_{e,\mathrm c}\) denotes the corresponding
fixed positive smooth distance factor.  The product on the right ranges in a
fixed interval \([c_Q,C_Q]\).  If
\(\lambda_iC_Q\le1\), the factor
\(\vartheta(P_i/t_n)\) is identically zero on the shell.  If
\(\lambda_ic_Q\ge2\), it is identically one.  In every remaining case,
\(\lambda_i\) belongs to the fixed compact interval
\([C_Q^{-1},2c_Q^{-1}]\).  Hence, by the chain rule applied to
\eqref{eq:rescaled-product-taper}, the family of all nonconstant transition
profiles has uniformly bounded derivatives of every fixed order.  Taking
the product over the finitely many vertices preserves this uniform bound.
This is the required control of the product taper; no derivative bound
contains a dyadic index or \(t_n\).

To make the pullback exact, choose once and for all a cutoff
\(\kappa_0\in C_c^\infty(\R^d)\) that equals one on a neighborhood of
\(K_0\).  For every fine edge, replace its rescaled factor by the same
factor multiplied by \(\kappa_0(z_e)\); this does not alter the multiplier
on the support of the shell cutoff.  Products of these extended factors,
together with the coarse-edge factors and the taper profiles above, define
functions \(H_{\bm m}(u,z)\) supported in one fixed compact set in each
\(z_e\).  Their derivatives through any prescribed finite order are
uniformly bounded.  After the explicit factor
\(2^{\sum_ea_em_e}\) is extracted, one has the exact identity
\[
 M_{\bm m}(T)
 =2^{\sum_ea_em_e}
   (\mathfrak P_{\bm m}H_{\bm m})(T).
\]
The matched period \(L_e=2^{m_e}\) creates no extraneous copies: modulo
\(L_e\), the point \(2^{m_e}(u_i-u_j)\) lies in the fixed support of
\(\kappa_0\) exactly when the torus difference has the unique fine-shell
lift used above.

The profiles now satisfy the hypotheses of
\cref{lem:configuration-pullback}.  Applying that lemma gives
\eqref{eq:configuration-multiplier-bound}.  Its Fourier proof also explains
why shared vertices do not create an additional loss: all pulled-back edge
frequencies add to one configuration-space character, and the Wiener
triangle inequality is taken only after that pullback.
\end{proof}


\part{Case A: two rough nuisances}

\section{Case A: construction}\label{A:sec:construction}

Throughout Part~A, \(\alpha,\beta\le1\), so \(s_1=(\alpha+\beta)/2=A_0/2\)
and \(\rho=1/\{1+d/(2\gamma)+d/(4s_1)\}\).  Both nuisances are carried by
the single rough field \(\Lambda_\lambda\) of \eqref{eq:rough}, at the fine
scale \(\ell\); there is no block carrier, the design is exactly uniform,
and no \(\varepsilon\)-loss is needed: the scales \(\Delta,h,\ell\) are
defined directly in \eqref{A:eq:amplitudes}, and the quantities \(r,t,Q\)
and the ledgers of Section~\ref{sec:scales} are not used.  The
construction is that of the
uniform-covariate note of Hou, written in the coded notation of
Section~\ref{sec:tools}.

\subsection{Scales and amplitudes}
Put
\begin{equation}
 \begin{gathered}
  \ja=c_a\ell^\alpha,\qquad \jb=c_b\ell^\beta,\qquad
  \Delta:=\ell^{2s_1}=\ell^{A_0},\qquad h=\Delta^{1/\gamma},\\
  d_*(x)=\ja\jb G_h(x)^2=c_ac_b\Delta G_h(x)^2,
 \end{gathered}
 \label{A:eq:amplitudes}
\end{equation}
with equal small constants \(c_a=c_b=c_0>0\) fixed in
Lemma~\ref{A:lem:legal}; the
scale \(\ell\) (equivalently \(\Delta\)) is tuned in
Section~\ref{A:sec:tuning}.  Since \(\gamma>A_0=2s_1\) (Lemma~\ref{lem:order}),
\(\ell=\Delta^{1/(2s_1)}\le\Delta^{1/\gamma}=h\) for small \(\Delta\).  Write
\(\jsh=\ja\wedge\jb\), \(\jfi=\ja\vee\jb\); then \(\jsh\jfi=\ja\jb=c_ac_b\Delta\)
and \(\ja,\jb\le1\).

\subsection{The two hypotheses}
In coded coordinates \((p,y,\tau)=(2\pi-1,2\mu_0-1,\tau)\), the rough
signs \(\lambda\) being drawn from the Rademacher prior on both sides and
the design being \(\operatorname{Unif}([0,1]^d)\) on both sides:
\begin{equation}
\begin{array}{c|ccc}
 &p&y&\tau\\ \hline
 Q&\ja\Lambda_\lambda&\jb\Lambda_\lambda&0\\
 P\ \text{(I: }\alpha\ge\beta)&0&\jb\Lambda_\lambda-d_*&d_*\\
 P\ \text{(II: }\beta\ge\alpha)&\ja\Lambda_\lambda&-d_*&d_*
\end{array}
\label{A:eq:two-sides}
\end{equation}
Under \(Q\) the two nuisances are proportional to the same rough field,
and their product produces the \(RT\)-cell \(\ja\jb\Lambda_\lambda^2\) with mean
\(\ja\jb G_h^2=d_*\); under \(P\) this cell is supplied by the treatment
effect \(d_*\), the rougher of the two nuisances keeps its field, and the
smoother one is shifted by \(-d_*\) to absorb the \(T\)-cell.

\begin{lemma}[Legality and separation]\label{A:lem:legal}
For sufficiently small fixed \(c_a,c_b>0\) and all sufficiently small
\(\Delta\), every law in either row of \eqref{A:eq:two-sides} belongs to
\(\cP\), and \(\tau_P(x_0)-\tau_Q(x_0)=c_ac_b\Delta\).
\end{lemma}
\begin{proof}
By Lemma~\ref{lem:holder-scaling}, \(\norm{\ja\Lambda_\lambda}_{C^\alpha}\le Cc_a\)
and \(\norm{\jb\Lambda_\lambda}_{C^\beta}\le Cc_b\), while
\(\abs{\ja\Lambda_\lambda}\le N_0c_a\), \(\abs{\jb\Lambda_\lambda}\le N_0c_b\).
The treatment effect \(d_*=c_ac_b\Delta G_h^2\) has
\(\norm{d_*}_{C^\gamma}\le Cc_ac_b\Delta h^{-\gamma}=Cc_ac_b\) and
\(\norm{d_*}_{C^\beta}\le C\Delta h^{-\beta}=Ch^{\gamma-\beta}\to0\) since
\(\gamma>A_0\ge\beta\).  Hence \(\pi=(1+p)/2\in\cH^\alpha(L_\pi)\),
\(\mu_0=(1+y)/2\in\cH^\beta(L_0)\), \(\tau\in\cH^\gamma(L_\tau)\) after
choosing \(c_a,c_b\) small, and \(\pi,\mu_0,\mu_0+\tau\in[\kappa,1-\kappa]\)
because \(\abs p,\abs y,\abs\tau\le N_0(c_a+c_b)+o(1)\).  The uniform
density lies in \([\underline f,\overline f]\).  Finally \(G_h(x_0)=1\).
\end{proof}

\subsection{Site likelihoods}
Fix a site \(x\) and write \(\Lambda=\Lambda_\lambda(x)\), \(G=G_h(x)\),
\(d=d_*(x)=\ja\jb G^2\), \(V=\Lambda^2-G^2\) (the centred square of
Section~\ref{sec:tools}, a Walsh polynomial of degree two with
\(\norm V_W\le N_0^2\)).  By \eqref{eq:walsh},
\begin{align}
 L^Q&=1+R\ja\Lambda+T\jb\Lambda+RT\ja\jb\Lambda^2
 =\bigl[1+R\ja\Lambda+T\jb\Lambda+RTd\bigr]+RT\ja\jb V,
 \label{A:eq:LQ}\\
 L^P&=1+T\jb\Lambda+RTd\qquad\text{(I)},\qquad
 L^P=1+R\ja\Lambda+Td\ja\Lambda+RTd\qquad\text{(II)},
 \label{A:eq:LP}
\end{align}
(for (II): \(y+\tau(1+p)=-d+d(1+\ja\Lambda)=d\ja\Lambda\) and
\(py+\tau(1+p)=-d\ja\Lambda+d(1+\ja\Lambda)=d\)).  The difference
\(D:=L^Q-L^P\) is
\begin{equation}
 D=R\ja\Lambda+RT\ja\jb V\quad\text{(I)},\qquad
 D=T\jb\Lambda-Td\ja\Lambda+RT\ja\jb V\quad\text{(II)}.
 \label{A:eq:D}
\end{equation}
Thus \(D\) has no constant term; its \emph{linear part} (monomials of
degree one) has Walsh norm \(\le CN_0\jsh\) in both subcases
(in (I) \(\ja=\jsh\); in (II) \(\jb=\jsh\) and \(d\ja\le\ja\jb\le\jsh\)); its
\emph{quadratic part} \(RT\ja\jb V\) has degree two and Walsh norm
\(\le N_0^2\ja\jb\).  Every part of \(L^P\) is constant, or linear with Walsh
norm \(\le CN_0\jfi\) (in (II) also \(d\ja\le\jfi\)).  In particular
\begin{equation}
 \E_\lambda L^Q=\E_\lambda L^P=1+RTd\qquad\text{in both subcases},
 \label{A:eq:one-site}
\end{equation}
because \(\E\Lambda=0\) and \(\E V=0\): the one-observation mixtures
coincide.

\section{Case A: proof}\label{A:sec:proof}

\subsection{Components}
Condition on \(X^n=x^n\).  Declare \(i\sim j\) if \(J(x_i)\cap J(x_j)\ne\varnothing\)
and let \(\{C\}\) be the components of the transitive closure; the sign
sets \(M_C=\bigcup_{i\in C}J(x_i)\) of distinct components are disjoint,
so the signs of distinct components are independent, and every link
satisfies \(x_i,x_j\in\supp G_h\), \(\norm{x_i-x_j}_\infty\le R_0\ell\).  The
\(X^n\)-marginals coincide (iid uniform), and relative to fair coins the
conditional laws factor over components:
\[
 \PP_n(w^n\mid x^n)=\prod_C\bar p_C(w_C),\qquad
 \mathbb Q_n(w^n\mid x^n)=\prod_C\bar q_C(w_C),
\]
\[
 \bar p_C(w_C)=4^{-\abs C}\E_\lambda\prod_{i\in C}L^P_i,\qquad
 \bar q_C(w_C)=4^{-\abs C}\E_\lambda\prod_{i\in C}L^Q_i .
\]
For a singleton \(C=\{i\}\), \eqref{A:eq:one-site} gives \(\bar p_C=\bar q_C\)
and \(H^2_C=0\).

\begin{lemma}[Component bound]\label{A:lem:component}
For every component \(C\) with \(q=\abs C\) and every cell \(w_C\),
\(\abs{\bar q_C(w_C)-\bar p_C(w_C)}\le4^{-q}C^q\,\ja\jb\), and
\(H^2_C\le C^q(\ja\jb)^2=C^q\Delta^2\).
\end{lemma}
\begin{proof}
Expand
\[
 \E_\lambda\prod_{i\in C}L^Q_i-\E_\lambda\prod_{i\in C}L^P_i
 =\sum_{\varnothing\ne S\subseteq C}\E_\lambda\Bigl[\prod_{i\in S}D_i\prod_{j\in C\setminus S}L^P_j\Bigr],
\]
and expand every \(D_i\) into its linear or quadratic part and every
\(L^P_j\) into its constant or linear part: at most \(C^q\) terms, each a
product of Walsh polynomials whose absolute mean is at most the product
of the Walsh norms of its factors.  A term with nonzero mean must have
even total degree (Section~\ref{sec:tools}), the degree of a product
being the sum of the degrees modulo two.  If the term contains a
quadratic part, its coefficient product is \(\le C^q\ja\jb\) (all other
coefficients are \(\le C\)).  Otherwise all its nonconstant parts are
linear; it contains the linear part of some \(D_i\), \(i\in S\)
(coefficient \(\le C\jsh\)), and by parity at least one further linear
part (coefficient \(\le C\jfi\)); so the coefficient product is
\(\le C^q\jsh\jfi=C^q\ja\jb\).  Summing over the terms gives the first
claim.  For the second, \(L^P_i\ge1-\abs p-\abs{y+\tau(1+p)}-\abs{py+\tau(1+p)}\ge\frac12\)
for small \(c_a,c_b\) (Lemma~\ref{A:lem:legal}), so \(\bar p_C\ge8^{-q}\), and
Lemma~\ref{lem:hellinger}(c) gives
\(H^2_C\le8^q\sum_{w_C}(4^{-q}C^q\ja\jb)^2\le C^q(\ja\jb)^2\).
\end{proof}

\begin{lemma}[Component counting]\label{A:lem:counting}
For every \(C_0>0\) there is \(C\) such that
\[
 \E\sum_{C:\abs C=q}C_0^q\le C^qn^qh^d\ell^{d(q-1)},\qquad q\ge2.
\]
Consequently, if \(n\ell^d\le c_0\) with \(c_0\) sufficiently small,
\(\E\sum_{C:\abs C\ge2}C_0^{\abs C}\le Cn^2h^d\ell^d\).
\end{lemma}
\begin{proof}
It suffices to count connected labelled \(q\)-sets, which may overcount
components.  Every such set contains a spanning tree; fix a labelled tree
and a root.  Each edge forces both endpoints into \(\supp G_h\)
(volume \(O(h^d)\)), and each child into a cube of volume \(O(\ell^d)\)
around its parent; since the \(X_i\) are iid uniform, the tree is realised
with probability \(\le Ch^d\ell^{d(q-1)}\).  By Cayley's formula there are
\(q^{q-2}\) labelled trees, and
\(\binom nq q^{q-2}\le\frac{n^q}{q!}q^{q-2}\le\frac{e^q}{q^2}n^q\le C^qn^q\)
(using \(q!\ge(q/e)^q\)); this gives the first bound.  Summing over \(q\ge2\) gives
\(Cn^2h^d\ell^d\sum_{q\ge2}(Cn\ell^d)^{q-2}\), a bounded geometric series
when \(n\ell^d\le c_0\).
\end{proof}

\subsection{Tuning and conclusion}\label{A:sec:tuning}
By Lemmas~\ref{lem:components} (with \(\mathcal E_n=\varnothing\)),
\ref{A:lem:component} and~\ref{A:lem:counting},
\begin{equation}
 \TV(\PP_n,\mathbb Q_n)^2\le\E\sum_{C:\abs C\ge2}H^2_C\le Cn^2h^d\ell^d\Delta^2
 =Cn^2\Delta^{d/\gamma+d/(2s_1)+2}=C\bigl(n\Delta^{1/\rho}\bigr)^2,
 \label{A:eq:hellinger-total}
\end{equation}
provided \(n\ell^d\le c_0\); here \(h^d=\Delta^{d/\gamma}\), \(\ell^d=\Delta^{d/(2s_1)}\),
and \(2+d/\gamma+d/(2s_1)=2/\rho\) because \(1/\rho=D_1/(A_0+2s_1)=1+d/(2\gamma)+d/(4s_1)\)
when \(A_0=2s_1\).  Choose
\begin{equation}
 \Delta=c_{\star}n^{-\rho},
 \label{A:eq:tuning}
\end{equation}
with a small fixed constant \(c_{\star}>0\).  Then
\(n\ell^d=c_{\star}^{d/(2s_1)}n^{1-\rho d/(2s_1)}\to0\), because
\(\rho>A_0/d=2s_1/d\) (Lemma~\ref{lem:order}); so \(n\ell^d\le c_0\) for large
\(n\), and \eqref{A:eq:hellinger-total} gives
\(\TV(\PP_n,\mathbb Q_n)\le C^{1/2}c_{\star}^{1/\rho}\le\frac12\) for \(c_{\star}\)
small.  Every law in the support of the \(P\)-prior has
\(\tau(x_0)=d_*(x_0)=c_ac_b\Delta\) and every law in the support of the
\(Q\)-prior has \(\tau(x_0)=0\) (Lemma~\ref{A:lem:legal}), so
Lemma~\ref{lem:two-point} gives
\[
 R^*_n(x_0)\ge\frac{c_ac_bc_{\star}}{4}\,n^{-\rho}
\]
for all large \(n\); smaller \(n\) are absorbed into the constant.  This
proves Theorem~\ref{thm:main} in case~(A), with \(\varepsilon=0\), and the
two priors are supported on laws with uniform design.  \qed

\begin{remark}
Case~(A) needs no design carrier: the shared rough field produces the
one-site matching \eqref{A:eq:one-site} exactly, and \(n\ell^d\to0\) makes
the packet-overlap components sparse.
\end{remark}


\part{Case B: two smooth nuisances}

\section{Case B: packets and carrier boxes}\label{B:sec:packets}

Throughout Part~B, \(\alpha,\beta>1\), so \(s_1=1\) and the scales of
Section~\ref{sec:scales} read \(t^2=\Delta r^{-A_0}\), \(\ell=tr\),
\(t_n=t\Lambda^B\); the fine ratio \(t\) is the relative size of the
virtual jitter below.
The rough field of Section~\ref{sec:tools} is not used in Part~B: both
nuisances are smooth, and the role of the rough field is played by
\emph{virtual} Rademacher signs attached to the observation sites.

Let \(\phi\) be the quadratic partition of unity of
Lemma~\ref{lem:partition}, \eqref{eq:quadratic-pu}, and let \(G_h\) be the
bump of Section~\ref{sec:tools}.  Define
\[
 x_k=x_0+rk,\qquad
 u_k(x)=\frac{x-x_k}{r},\qquad
 \psi_k(x)=G_h(x)\phi(u_k(x)).
\]
Then
\begin{equation}
 \sum_k\psi_k(x)^2=G_h(x)^2.
 \label{B:eq:quadratic-partition}
\end{equation}
(In Part~B the word \emph{packet} refers to these block bumps at scale
\(r\), not to the rough packets of Section~\ref{sec:tools}.)  Only a fixed
number of packets meet a point, and there are \(O(h^dr^{-d})\) active
packets.  Enlarge the support of each active packet
to a cube \(S_k\) of side \(C r\), with \(C\) fixed.  The cubes have bounded
overlap and their union has volume \(O(h^d)\).  Each \(S_k\) is identified,
by an affine map \(\theta_k\), with the normalized torus \(\Torus^d\).
The packet support is contained in a fixed compact subcube of that torus.

\section{Case B: the exact aggregate cubic bridge}

\subsection{Fields and virtual jitter}

Let
\begin{equation}
 a=c_ar^\alpha,\qquad
 b=c_br^\beta,\qquad
 d_*(x)=a\,b\,t^2G_h(x)^2.
 \label{B:eq:amplitudes}
\end{equation}
For bounded coefficient vectors \(U_k\), define
\begin{equation}
 P(x)=\sum_k\psi_k(x)v_Q(\theta_k(x))^\top U_k,
 \qquad p(x)=aP(x).
 \label{B:eq:p-field}
\end{equation}
For a separated tuple of sites in a block, use \cref{lem:lagrange} and put
\begin{equation}
 \widetilde U_k
 =U_k+t\sum_i c_{ki}(T)\zeta_{ki},
 \qquad
 \zeta_{ki}\stackrel{\mathrm{ind}}{\sim}\Rad(\tfrac12),
 \label{B:eq:virtual-shift}
\end{equation}
where \(\Rad(1/2)\) denotes the symmetric \(\{\pm1\}\) law.  At an
observed point,
\begin{equation}
 p_i'=p_i+\delta_i,
 \qquad
 \delta_i=a\,t\sum_k\psi_k(x_i)\zeta_{ki}.
 \label{B:eq:delta}
\end{equation}
Conditional on the baseline field, the \(\delta_i\)'s are independent and
symmetric.  By \eqref{B:eq:quadratic-partition},
\begin{align}
 v(x)&:=\E\delta(x)^2=a^2t^2G_h(x)^2,
 \label{B:eq:variance}\\
 m_4(x)&:=\E\delta(x)^4
 =a^4t^4G_h(x)^4\{3-2q_4(x)\},
 \label{B:eq:fourth-moment}\\
 q_4(x)&=\sum_k\phi(u_k(x))^4.
\end{align}

\subsection{Cubic identity}

Define
\begin{equation}
 K_x(z)
 =\frac ba\left[\{1+\eta(x)\}z-\frac{z^3}{3}\right],
 \qquad
 \eta(x)=\frac{a^2t^2G_h(x)^2}{3}\{3-2q_4(x)\}.
 \label{B:eq:cubic}
\end{equation}
This is also the definition when \(G_h=0\).  Put
\begin{equation}
 u_x(p)=\E_\delta K_x(p+\delta)=K_x(p)-d_*(x)p.
 \label{B:eq:u}
\end{equation}
The two fuzzy sides use
\begin{equation}
\begin{array}{c|ccc}
 &2\pi-1&2\mu_0-1&\tau\\ \hline
 P&p&u_x(p)-d_*(1+p)&d_*\\
 Q&p'&K_x(p')&0.
\end{array}
\label{B:eq:two-sides}
\end{equation}

\begin{lemma}[Cubic covariance identity]
\label{B:lem:cubic}
For every fixed baseline value \(p\),
\begin{equation}
 \operatorname{Cov}\{p+\delta,K_x(p+\delta)\}
 =d_*(x)(1-p^2).
 \label{B:eq:cubic-cov}
\end{equation}
Consequently,
\begin{equation}
 \E_\delta(1,p',K_x(p'),p'K_x(p'))
 =(1,p,u_x(p),pu_x(p)+d_*(1-p^2)).
 \label{B:eq:walsh-vector}
\end{equation}
\end{lemma}

\begin{proof}
Write \(K_x(z)=c_1z+c_3z^3\), where
\[
 c_3=-\frac b{3a}=-\frac{d_*}{3v},
 \qquad
 c_1=\frac ba(1+\eta)
     =\frac{d_*}{v}+\frac{d_*m_4}{3v^2}
\]
when \(v>0\); the identities extend continuously to \(v=0\).  Symmetry
of \(\delta\) gives
\[
 \operatorname{Cov}\{p+\delta,K_x(p+\delta)\}
 =c_1v+c_3(3p^2v+m_4)
 =d_*(1-p^2).
\]
It also gives
\[
 \E K_x(p+\delta)
 =K_x(p)+3c_3pv=K_x(p)-d_*p.
\]
Substitution proves \eqref{B:eq:walsh-vector}.
\end{proof}

The right side of \eqref{B:eq:walsh-vector} is exactly the pair of outcome
Walsh coefficients generated by the \(P\)-row of \eqref{B:eq:two-sides}.
Conditional independence of the site jitters therefore gives equality of all
\(4^q\) joint cells for every finite tuple when every active block is
virtually shifted.  The identity uses the aggregate \(p\), so overlapping
smooth packets cause no additional term.

\section{Case B: a positive pre-data moment carrier}

\subsection{The tapered ideal increment and its Wiener norm}

For a \(Q\)-tuple \(T\), write
\[
 C_T(\zeta)=\sum_{i=1}^Qc_i(T)\zeta_i.
\]
On the support of \(\chi_Q\), \cref{lem:lagrange} gives
\[
 t\norm{C_T(\zeta)}\le C_Q\frac{t}{P_{\min}(T)}
 \le C_Q\Lambda^{-B}.
\]
Choose the support ball in \cref{lem:moment-simplex} strictly inside a fixed
legal coefficient ball.  For large \(n\), the displayed virtual shifts stay
inside that legal ball wherever \(\chi_Q\ne0\).  Define the vector-valued
symmetric tensor
\begin{equation}
 J_n(T)=\chi_Q(T)
 \left[
 \E_{\Pi_0,\zeta}\Psi_{4Q}\{U+tC_T(\zeta)\}
 -\E_{\Pi_0}\Psi_{4Q}(U)
 \right],
 \label{B:eq:ideal-increment}
\end{equation}
with the product extended by zero across the diagonals.
\begin{lemma}[Inverse-square Wiener estimate]
\label{B:lem:wiener-increment}
With \(E_Q=\binom Q2\),
\begin{equation}
 \norm{J_n}_{\cW_{dQ}}
 \le C_Q\{1+\log(1/t_n)\}^{E_Q}
       \sum_{j=2}^{4Q}\left(\frac{t}{t_n}\right)^j.
 \label{B:eq:wiener-increment-strong}
\end{equation}
Consequently,
\begin{equation}
 \norm{J_n}_{\cW_{dQ}}
 \le C_Q\Lambda^{E_Q-2B}=:\delta_n.
 \label{B:eq:wiener-increment}
\end{equation}
In particular, any fixed choice
\begin{equation}
 B>\frac{E_Q+2}{2}
 \label{B:eq:B-choice}
\end{equation}
makes \(\delta_n=o(1)\).  All constants in this lemma are independent of
\(n\) and of the shell indices, with the fixed-dependence convention stated
after \eqref{eq:Q}.
\end{lemma}

\begin{proof}
Fix one coordinate monomial \(U^\nu\) of \(\Psi_{4Q}\), with
\(1\le\abs\nu\le4Q\), and expand
\((U+tC_T(\zeta))^\nu-U^\nu\) by the multinomial theorem.  After integrating
\(U\sim\Pi_0\), every coefficient is bounded by a constant depending only
on \(Q\) and the fixed support of \(\Pi_0\).  A term containing \(j\)
shift factors is a finite linear combination of expressions
\begin{equation}
 t^j\prod_{r=1}^j c_{i_r,a_r}(T),
 \label{B:eq:generic-shift-term}
\end{equation}
where \(c_{i,a}\) denotes a coordinate of \(c_i\).  The expectation over the
independent symmetric signs is zero unless every site index occurs an even
number of times among \(i_1,\ldots,i_j\).  Thus every surviving term has
\(j\ge2\); this is what gives an inverse-square rather than an
inverse-first-order gain.

Insert the partition \eqref{eq:configuration-shell-partition} into a
surviving term, and let \(M_i(\bm m)\) be as in \eqref{eq:vertex-shell-sum}.
Since every coordinate of \(c_i(T)\) is a finite sum of products
containing exactly one elementary multiplier
\eqref{eq:elementary-lagrange-multipliers} for every edge incident to
\(i\), in the product in \eqref{B:eq:generic-shift-term} edge \(e\) occurs
with some inverse order \(a_e\), and
\begin{equation}
 \sum_{e\in\mathfrak E_Q}a_e\mu(m_e)
 =\sum_{r=1}^jM_{i_r}(\bm m).
 \label{B:eq:inverse-order-ledger}
\end{equation}

Apply \cref{cor:configuration-multipliers} to the product of the shell
cutoffs, the product-distance taper, and the elementary multipliers.  It
gives
\begin{align}
 &\norm{
   \chi_Q\rho_{\bm m}
   \prod_{r=1}^jc_{i_r,a_r}
  }_{\cW_{dQ}}
 \le C_Q2^{\sum_rM_{i_r}(\bm m)}.
 \label{B:eq:shellwise-c-bound}
\end{align}
There is no additional scale loss from shared vertices, different edge
scales, or the product taper; those points are exactly what
\cref{cor:configuration-multipliers} proves.  On a shell with
\(\chi_Q\rho_{\bm m}\not\equiv0\), \cref{lem:admissible-shells} yields
\[
 2^{M_i(\bm m)}\le C_Qt_n^{-1},
\]
and hence
\begin{equation}
 \norm{
   \chi_Q\rho_{\bm m}
   t^j\prod_{r=1}^jc_{i_r,a_r}
  }_{\cW_{dQ}}
 \le C_Q\left(\frac{t}{t_n}\right)^j.
 \label{B:eq:one-shell-shift-bound}
\end{equation}

There are only finitely many coordinates, monomials, and multinomial terms,
with their number depending on \(Q\) but not on \(n\).  Sum
\eqref{B:eq:one-shell-shift-bound} over them, over \(j=2,\ldots,4Q\), and over
the \(\le C_Q\{1+\log(1/t_n)\}^{E_Q}\) shell assignments of
\cref{lem:admissible-shells}.  This proves
\eqref{B:eq:wiener-increment-strong}.  Finally,
\(t_n=t\Lambda^B\), so
\[
 \sum_{j=2}^{4Q}(t/t_n)^j
 =\sum_{j=2}^{4Q}\Lambda^{-Bj}
 \le C_Q\Lambda^{-2B},
\]
and \(1+\log(1/t_n)\le C_\varepsilon\Lambda\) because \(t\) is a fixed
negative power of \(n\).  This proves \eqref{B:eq:wiener-increment}.  Under
\eqref{B:eq:B-choice}, \(E_Q-2B<-2\), and therefore \(\delta_n=o(1)\).
\end{proof}

\subsection{A fixed positive polarization operator}

\begin{lemma}[Positive Wiener polarization]
\label{B:lem:positive-polarization}
Fix \(Q\) and \(0<\eta_f<1/2\).  There is a countable dictionary
\(\{F_m:m\ge0\}\subset\cW_d\), fixed independently of the tensor to be
decomposed, with
\begin{equation}
 \int_{\Torus^d}F_m=1,
 \qquad
 1-\eta_f\le F_m\le1+\eta_f,
 \qquad
 \sup_m\norm{F_m}_{\cW_d}<\infty,
 \label{B:eq:positive-atoms}
\end{equation}
and bounded linear maps
\(A\mapsto a_m(A)\) from the symmetric real Wiener tensors on
\((\Torus^d)^Q\) to any fixed finite-dimensional real vector space \(V\),
such that
\begin{align}
 A&=\sum_{m\ge0}a_m(A)F_m^{\otimes Q},
 \label{B:eq:polarization-expansion}\\
 \sum_{m\ge0}\norm{a_m(A)}
 &\le C_{Q,\eta_f}\norm A_{\cW_{dQ}},
 \label{B:eq:polarization-bound}\\
 \sum_{m\ge0}\norm{a_m(A)-a_m(A')}
 &\le C_{Q,\eta_f}\norm{A-A'}_{\cW_{dQ}}.
 \label{B:eq:polarization-lipschitz}
\end{align}
The series converges absolutely in Wiener norm.
\end{lemma}

\begin{proof}
Use the real trigonometric basis consisting of \(1\), sines, and cosines,
normalized so that every nonconstant basis element \(w\) has mean zero,
\(\norm w_\infty\le1\), and \(\norm w_{\cW_d}=1\).  Passing from the
complex Fourier basis to this real basis changes the \(\ell^1\) coefficient
norm by only a fixed factor depending on \(Q\).

It suffices to decompose a symmetrized basis tensor
\[
 w_1\od\cdots\od w_j\od1^{\od(Q-j)},
 \qquad 0\le j\le Q,
\]
where each \(w_i\) is nonconstant.  For \(j\ge1\), apply
\cref{lem:polarization-identity} in the Wiener algebra \(\cW_d\) with
\(g_i=w_i\) and a fixed \(0<\theta\le\eta_f/2\): every symmetrized basis
tensor is a fixed finite linear combination of \(F^{\otimes Q}\) with
\(F=1+\theta t_rg_\epsilon\), \(g_\epsilon=j^{-1}\sum_i\epsilon_iw_i\).  Since
\(\int g_\epsilon=0\), \(\norm{g_\epsilon}_\infty\le1\) and
\(\norm{g_\epsilon}_{\cW_d}\le1\), such \(F\) has integral one, lies between
\(1-\eta_f\) and \(1+\eta_f\), and has uniformly bounded Wiener norm.  The
case \(j=0\) uses \(F=1\).

Apply this fixed stencil coefficientwise to the absolutely summable real
Fourier expansion of \(A\).  Index all resulting atoms by a single
countable index \(m\).  The total absolute coefficient generated by one
basis tensor is bounded by a constant depending only on
\((Q,\eta_f)\).  Summing the Fourier coefficients proves
\eqref{B:eq:polarization-bound}.  The construction is linear, so applying the
same estimate to \(A-A'\) proves \eqref{B:eq:polarization-lipschitz}.  Absolute
Wiener convergence gives \eqref{B:eq:polarization-expansion}.
\end{proof}

\subsection{The self-consistent positive carrier}

\begin{proposition}[Positive pre-data carrier]
\label{B:prop:carrier}
For all sufficiently large \(n\), there are joint block laws
\(\nu_{P,k}(dH_k,dU_k)\) and \(\nu_{Q,k}(dH_k,dU_k)\) with the following
properties.

\begin{enumerate}
\item Their full \(H_k\)-marginals are identical.  The label \(H_k\)
selects a factor \(F_{H_k}\) satisfying
\begin{equation}
 \int_{\Torus^d}F_{H_k}=1,
 \qquad
 1-\eta_f\le F_{H_k}\le1+\eta_f.
 \label{B:eq:local-density-factor}
\end{equation}
All coefficient draws lie in one fixed legal ball.

\item If
\begin{equation}
 D_Q(T)=\E F_{H_k}(u_1)\cdots F_{H_k}(u_Q),
 \label{B:eq:DQ}
\end{equation}
then the same density law that generates \(D_Q\) satisfies the exact moment
identity
\begin{align}
 &\E_{\nu_{Q,k}}
 \left[F_{H_k}^{\otimes Q}(T)\Psi_{4Q}(U_k)\right]
 -\E_{\nu_{P,k}}
 \left[F_{H_k}^{\otimes Q}(T)\Psi_{4Q}(U_k)\right]
 \nonumber\\
 &\hspace{2cm}=D_Q(T)J_n(T).
 \label{B:eq:master-moment}
\end{align}
The laws are chosen before observing the data and are not tuple-indexed.
\end{enumerate}
\end{proposition}

\begin{proof}
For a symmetric scalar tensor \(D\), set
\[
 A(D)=D J_n.
\]
Apply \cref{B:lem:positive-polarization} coordinatewise with
\(V=\R^{N_Q}\):
\[
 A(D)=\sum_m a_m(D)F_m^{\otimes Q}.
\]
Let \(C_{\mathrm{pol}}\) be a constant valid in both
\eqref{B:eq:polarization-bound} and \eqref{B:eq:polarization-lipschitz}.  Choose
\(K_0>2/\delta_Q\), where \(\delta_Q\) is from
\cref{lem:moment-simplex}, and define
\[
 \omega_m(D)=K_0\norm{a_m(D)},
 \qquad
 \Phi(D)=
 \left(1-\sum_m\omega_m(D)\right)1^{\otimes Q}
 +\sum_m\omega_m(D)F_m^{\otimes Q}.
\]
Set
\[
 C_F=\sup_m
 \norm{F_m^{\otimes Q}-1^{\otimes Q}}_{\cW_{dQ}}<\infty,
 \qquad
 C_0=K_0C_FC_{\mathrm{pol}},
 \qquad
 C_1=K_0C_{\mathrm{pol}}.
\]
For any symmetric \(D,D'\), the reverse triangle inequality,
\eqref{B:eq:polarization-lipschitz}, and the Wiener algebra inequality give
\begin{align}
 \norm{\Phi(D)-\Phi(D')}_{\cW_{dQ}}
 &\le K_0C_F\sum_m
   \abs{\norm{a_m(D)}-\norm{a_m(D')}}\nonumber\\
 &\le C_0\norm{(D-D')J_n}_{\cW_{dQ}}\nonumber\\
 &\le C_0\delta_n\norm{D-D'}_{\cW_{dQ}}.
 \label{B:eq:phi-contraction}
\end{align}

Define the closed ball in the symmetric Wiener tensors
\begin{equation}
 \mathbb B_n
 =\left\{D=D^{\mathrm{sym}}:
   \norm{D-1^{\otimes Q}}_{\cW_{dQ}}\le R_n\right\},
 \qquad
 R_n=2C_0\delta_n.
 \label{B:eq:fixed-point-ball}
\end{equation}
For \(D\in\mathbb B_n\),
\begin{align}
 \sum_m\omega_m(D)
 &\le C_1\norm{DJ_n}_{\cW_{dQ}}
 \le C_1(1+R_n)\delta_n
 \label{B:eq:weight-bound}\\
 \norm{\Phi(D)-1^{\otimes Q}}_{\cW_{dQ}}
 &\le C_F\sum_m\omega_m(D)
 \le C_0(1+R_n)\delta_n.
 \label{B:eq:phi-self-map-bound}
\end{align}
Since \(\delta_n=o(1)\), choose \(n\) so large that
\begin{equation}
 R_n\le1,
 \qquad C_0\delta_n\le\tfrac12,
 \qquad 2C_1\delta_n<1.
 \label{B:eq:fixed-point-smallness}
\end{equation}
Then \eqref{B:eq:phi-self-map-bound} is at most
\(2C_0\delta_n=R_n\), so \(\Phi(\mathbb B_n)\subseteq\mathbb B_n\).
Equation \eqref{B:eq:phi-contraction} has contraction constant
\(q_n=C_0\delta_n\le1/2\), and
\eqref{B:eq:weight-bound} gives \(\sum_m\omega_m(D)<1\) throughout the ball.
The weights are nonnegative by definition.  The ball \(\mathbb B_n\) is a
closed subset of a Banach space, so Banach's theorem yields a unique fixed
point
\begin{equation}
 D_Q=\Phi(D_Q)\in\mathbb B_n.
 \label{B:eq:D-fixed-point}
\end{equation}

Use the probabilities \(\omega_m(D_Q)\) on labels carrying \(F_m\), and put
the remaining probability on a label carrying \(F=1\).  This common label
law is the \(H_k\)-marginal on both sides and, by
\eqref{B:eq:D-fixed-point}, its \(Q\)-fold density moment is exactly \(D_Q\).
Under \(P\), conditionally on every label, let \(U_k\sim\Pi_0\).  If
\(\omega_m>0\), define
\[
 v_m=\frac{a_m(D_Q)}{\omega_m},
 \qquad \norm{v_m}=K_0^{-1}<\delta_Q,
\]
and under \(Q\), conditionally on label \(m\), use the law \(\Pi_{v_m}\)
from \cref{lem:moment-simplex}; on the remaining label use \(\Pi_0\).  If
\(\omega_m=0\), no conditional law for that null label is needed.  All
positive-mass conditional laws have the same fixed legal support.  Therefore
\[
 \E_Q[F^{\otimes Q}\Psi]-\E_P[F^{\otimes Q}\Psi]
 =\sum_m\omega_mF_m^{\otimes Q}v_m
 =\sum_ma_m(D_Q)F_m^{\otimes Q}
 =D_QJ_n,
\]
which is \eqref{B:eq:master-moment}.  There is no circularity: the fixed
point determines one common label law, and this law both generates
\(D_Q\) and carries the moment perturbation.
\end{proof}

\subsection{Evaluated projectivity and ghost leakage}

For \(q\le Q\), let
\[
 D_q(T_q)=\E F_{H_k}(u_1)\cdots F_{H_k}(u_q).
\]
Since every density atom has integral one,
\begin{equation}
 D_q(T_q)=\int D_Q(T_q,Z)\,dZ,
 \label{B:eq:D-projectivity}
\end{equation}
where \(Z\in(\Torus^d)^{Q-q}\) and Haar measure is normalized.
Define
\begin{equation}
 \overline\chi_q(T_q)
 =\frac{\int D_Q(T_q,Z)\chi_Q(T_q,Z)\,dZ}{D_q(T_q)}.
 \label{B:eq:chi-bar}
\end{equation}
Because every \(F\) is between fixed positive constants,
\(D_q\) is bounded above and below, and \(0\le\overline\chi_q\le1\).

A likelihood evaluation polynomial for one block means a polynomial in
\(U_k\) of degree at most \(4Q\) that depends on \(U_k\) only through the
retained linear evaluations
\[
 \psi_k(x_i)v_Q(\theta_k(x_i))^\top U_k,
 \qquad i=1,\ldots,q.
\]
Such polynomials are exactly the block factors obtained after expanding a
component Walsh coefficient.

\begin{proposition}[Evaluated lower-order projectivity and ghost leakage]
\label{B:prop:evaluated-projectivity}
Let \(\mathcal L_{T_q}\) be any likelihood evaluation functional as above.
If \(b_{\mathcal L,q}(T_q)\) denotes the corresponding ideal virtual moment
increment at the retained sites, then
\begin{equation}
 \mathcal L_{T_q}
 \left\{
 \E_Q[F^{\otimes q}\Psi]-\E_P[F^{\otimes q}\Psi]
 \right\}
 =D_q(T_q)\overline\chi_q(T_q)b_{\mathcal L,q}(T_q).
 \label{B:eq:evaluated-projectivity}
\end{equation}
Moreover,
\begin{equation}
 \int D_q(T_q)\{1-\overline\chi_q(T_q)\}\,dT_q
 \le C_Qt^d\Lambda^{M_Q}.
 \label{B:eq:ghost-mass}
\end{equation}
Projectivity is used only after evaluation; the tapered coefficient
tensor itself need not be projective.
\end{proposition}

\begin{proof}
Integrating \eqref{B:eq:master-moment} in the last \(Q-q\) coordinates gives
the exact lower-order coefficient moment.  For
\(T=(T_q,Z)\) with \(\chi_Q(T)>0\), the Lagrange identities imply
\[
 v_Q(u_j)^\top\left\{U+t\sum_{i=1}^Qc_i(T)\zeta_i\right\}
 =v_Q(u_j)^\top U+t\zeta_j,
 \qquad 1\le j\le q.
\]
Thus, after applying \(\mathcal L_{T_q}\), the ideal increment in
\eqref{B:eq:ideal-increment} is independent of the ghost tuple \(Z\) and is
exactly \(b_{\mathcal L,q}(T_q)\).  Pull it through the integral and use
\eqref{B:eq:chi-bar}; this proves \eqref{B:eq:evaluated-projectivity}.

For \eqref{B:eq:ghost-mass}, use Fubini and
\eqref{B:eq:D-projectivity}:
\begin{align*}
 \int D_q(1-\overline\chi_q)
 &=\iint D_Q(T_q,Z)\{1-\chi_Q(T_q,Z)\}\,dZ\,dT_q\\
 &\le C_Q\abs{\{T\in(\Torus^d)^Q:\chi_Q(T)\ne1\}}.
\end{align*}
The result follows from \cref{lem:bad-volume}.
\end{proof}

\section{Case B: global priors and parameter legality}

\subsection{The two parameter fields}

Take independent copies of the local carrier from \cref{B:prop:carrier} over
all active blocks.  For a coefficient realization, form \(p\) by
\eqref{B:eq:p-field}, and define the two sides by \eqref{B:eq:two-sides}.  On the
\(Q\)-side the actual coefficient draws are those supplied by
\(\nu_{Q,k}\); the virtual shifts are used only to identify the finite
moments of the mixed likelihood.

\begin{lemma}[Hölder legality and positivity]
\label{B:lem:legality}
For sufficiently small fixed \(c_a,c_b>0\), every parameter draw on both
sides belongs to \(\cP(\alpha,\beta,\gamma)\), apart from the design density,
which is handled in the next subsection.  Moreover,
\begin{equation}
 \tau_P(x_0)-\tau_Q(x_0)=c_ac_b\Delta.
 \label{B:eq:separation}
\end{equation}
\end{lemma}

\begin{proof}
By the Hölder scaling rule (\cref{lem:holder-scaling}), all packet fields
used below satisfy, for each fixed \(t\),
\begin{equation}
 \norm P_{C^t}+\norm{P'}_{C^t}+\norm{q_4}_{C^t}
 \le C_{t,Q}r^{-t},
 \qquad
 \norm{G_h}_{C^t}\le C_t h^{-t}\le C_tr^{-t}.
 \label{B:eq:packet-holder-scaling}
\end{equation}

Since \(a=c_ar^\alpha\), \eqref{B:eq:packet-holder-scaling} gives
\[
 \norm p_{C^\alpha}+\norm{p'}_{C^\alpha}
 \le C_Qc_a.
\]
Thus the two propensity fields lie in the prescribed Hölder ball after
choosing \(c_a\) sufficiently small.

Next write \(p'=aP'\).  From \eqref{B:eq:cubic},
\begin{equation}
 K_x(p')
 =b\left[\{1+\eta(x)\}P'-\frac{a^2(P')^3}{3}\right].
 \label{B:eq:k-smoothness}
\end{equation}
The fixed-overlap packet calculus and \eqref{B:eq:packet-holder-scaling} give
\[
 \norm{P'}_{C^\beta}+\norm{(P')^3}_{C^\beta}
 \le C_Qr^{-\beta}.
\]
Moreover,
\[
 \norm\eta_\infty\le C_Qa^2t^2,
 \qquad
 \norm\eta_{C^\beta}\le C_Qa^2t^2r^{-\beta},
\]
because \(r<h\), \(q_4\) varies at scale \(r\), and \(G_h\) varies at
scale \(h\).  The Hölder product inequality therefore yields
\begin{align*}
 \norm{K_x(p')}_{C^\beta}
 &\le C_Qbr^{-\beta}\{1+a^2+a^2t^2\}\\
 &\le C_Qc_b,
\end{align*}
since \(b=c_br^\beta\), \(a\le1\), and \(t\to0\).  The same estimate holds
for \(K_x(p)\).

The identity \(a\,b\,t^2=c_ac_b\Delta\) gives
\begin{equation}
 d_*(x)=c_ac_b\Delta G_h(x)^2,
 \qquad
 \norm{d_*}_{C^\gamma}
 \le Cc_ac_b\Delta h^{-\gamma}=Cc_ac_b.
 \label{B:eq:dstar-smooth}
\end{equation}
Since \(\gamma>A_0>\beta\), also
\[
 \norm{d_*}_{C^\beta}
 \le Cc_ac_b\Delta h^{-\beta}
 =Cc_ac_bh^{\gamma-\beta}.
\]
Although a draw of \(p\) is smooth, its \(C^\beta\) norm can grow when
\(\beta>\alpha\); the small target amplitude compensates for this.  Indeed,
using the Hölder product inequality and
\eqref{B:eq:packet-holder-scaling},
\begin{align*}
 \norm{d_*p}_{C^\beta}
 &\le C_Q\left\{
   \Delta a r^{-\beta}+a\Delta h^{-\beta}
   \right\}\\
 &\le C_Q\left\{c_at^2r^{2\alpha}
          +a h^{\gamma-\beta}\right\}
 =o(1).
\end{align*}
Here
\(\Delta a r^{-\beta}=c_at^2r^{2\alpha}\), because
\(t^2=\Delta r^{-A_0}\) and \(A_0=\alpha+\beta\).  Consequently
\[
 u_x(p)-d_*(1+p)=K_x(p)-d_*(1+2p)
\]
has \(C^\beta\) norm at most \(C_Qc_b+o(1)\).  Choose \(c_b\)
first, then \(n\) large, so that this lies within the remaining radius of
the prescribed outcome ball about \(1/2\).  This proves the required
smoothness of both baseline outcome regressions, while
\eqref{B:eq:dstar-smooth} proves the \(C^\gamma\) bound for the treatment
effect.

All coded fields have sup norm tending to zero with \(c_a,c_b\) fixed and
small: \(p=O(a)\), \(K_x(p)=O(b)\), and \(d_*=O(\Delta)\).  First choose
\(c_a,c_b\) small enough to fit the fixed Hölder radii, and then take \(n\)
large enough.  The treatment propensity and both Bernoulli outcome means are
then contained in \([\kappa,1-\kappa]\) on every draw.  Finally,
\(G_h(x_0)=1\), \(\tau_P=d_*\), and \(\tau_Q=0\), so
\[
 \tau_P(x_0)-\tau_Q(x_0)=c_ac_b\Delta.
\]
\end{proof}

\subsection{Design density and tilt}

Transfer each torus factor to its physical carrier cube and set it equal to
one outside that cube.  Write this physical factor again as \(F_{H_k}(x)\).
It satisfies
\[
 \int_{S_k}\{F_{H_k}(x)-1\}\,dx=0.
\]
Define
\begin{equation}
 \widetilde f_H(x)
 =\prod_{k:x\in S_k}F_{H_k}(x),
 \qquad
 Z(H)=\int_{[0,1]^d}\widetilde f_H(x)\,dx,
 \qquad
 f_H=\widetilde f_H/Z(H).
 \label{B:eq:density}
\end{equation}
Let \(N_1\) be the maximal number of carrier boxes \(S_k\) containing a
point (a fixed number).  Then
\begin{equation}
 \frac{(1-\eta_f)^{N_1}}{(1+\eta_f)^{N_1}}
 \le f_H(x)\le
 \frac{(1+\eta_f)^{N_1}}{(1-\eta_f)^{N_1}}.
 \label{B:eq:density-bounds}
\end{equation}
Choose \(\eta_f>0\) sufficiently small that these bounds lie inside
\([\underline f,\overline f]\).  Let \(\nu_S=\bigotimes_k\nu_{S,k}\),
\(S\in\{P,Q\}\); their \(H\)-marginals are identical
(\cref{B:prop:carrier}).  Apply \cref{lem:tilt} (with \(\lambda\) trivial):
the tilted priors \eqref{eq:z-tilt} induce the same law of \(f_H\) and the
same \(X^n\)-marginal, and the mixed joint density is given by
\eqref{eq:z-cancel} with the product law \(\nu_S\) restored.

\section{Case B: fixed-sample information bound}

\subsection{Partial-jitter stability}

For a subset \(\cS\) of the active blocks at a site, define
\[
 \delta_{\cS}
 =a\,t\sum_{k\in\cS}\psi_k\zeta_k,
 \qquad
 v_{\cS}=\E\delta_{\cS}^2\le v,
 \qquad
 m_{4,\cS}=\E\delta_{\cS}^4\le3v_{\cS}^2\le3v^2.
\]
The exact partial moments are
\begin{align}
 \E K_x(p+\delta_{\cS})
 &=K_x(p)-\frac ba p v_{\cS},
 \label{B:eq:partial-k}\\
 \E\{(p+\delta_{\cS})K_x(p+\delta_{\cS})\}
 &=pK_x(p)+c_1v_{\cS}
   +c_3\{6p^2v_{\cS}+m_{4,\cS}\}.
 \label{B:eq:partial-pk}
\end{align}

\begin{lemma}[Uniform stability under partial block shifts]
\label{B:lem:partial-stability}
Fix \(q\le Q\), a baseline coefficient realization, and any subset of the
blocks meeting the \(q\) sites.  Average only the virtual jitters in that
subset.  Every joint \(q\)-site Walsh coefficient differs from the
corresponding \(P\)-side coefficient by at most \(C_Q\Delta\).  If all
active blocks are shifted, the difference is zero.
\end{lemma}

\begin{proof}
The propensity coefficient at a site remains \(p\).  Since
\(v=a^2t^2G_h^2\), \(d_*=(b/a)v\), and \(\abs p\le Ca\),
\eqref{B:eq:partial-k} differs from the full-jitter value by at most
\[
 \frac ba\abs p\,(v-v_{\cS})\le C d_*.
\]
Using \(c_1=(b/a)(1+\eta)\), \(c_3=-b/(3a)\), and
\(m_{4,\cS}\le3v^2\), the difference between
\eqref{B:eq:partial-pk} and the full-jitter value is also at most \(Cd_*\).
Here the terms are bounded respectively by
\((b/a)v=d_*\), \((b/a)p^2v\le Ca^2d_*\), and
\((b/a)v^2=d_*v\).  Since \(d_*(x)\le C\Delta\), each one-site outcome
Walsh coefficient differs from its target value by \(C\Delta\).

Conditional on the baseline field, site signs are independent.  Hence a
joint Walsh coefficient is a product of one-site coefficients, and a
finite telescoping bound gives \(C_Q\Delta\).  A product containing only
propensity characters is unchanged exactly: each site propensity appears
only once, and every nonempty jitter term has an odd site sign and averages
to zero.  When all active blocks are shifted,
\cref{B:lem:cubic} gives exact equality.
\end{proof}

\subsection{Component factorization and the local discrepancy}

Condition on the common design \(X^n=x^n\), and write
\(\mathcal K_i=\{k:x_i\in S_k\}\).  Relative to counting measure on the
binary signs, the tilted mixed joint density has the exact form
\begin{equation}
 \overline p_S(x^n,w^n)
 =\frac{4^{-n}}{\E Z(H)^n}
  \int
  \prod_k\left\{\prod_{i:k\in\mathcal K_i}F_{H_k}(x_i)\right\}
  \prod_{i=1}^nL_{S,i}(w_i\mid x_i,U_{\mathcal K_i})
  \,d\nu_S(H,U),
 \label{B:eq:mixed-joint-factorization}
\end{equation}
where \(S\in\{P,Q\}\), and the common \(X^n\) density is obtained by
removing the likelihood factors.  Formula
\eqref{B:eq:mixed-joint-factorization} follows directly from
\eqref{eq:z-cancel}; in particular, no residual \(Z(H)\)-dependence remains.

Form the bipartite incidence graph whose vertices are the observations and
the active block labels, with an edge \(i\sim k\) when
\(k\in\mathcal K_i\).  Equivalently, connect observations \(i,j\) when
\(\mathcal K_i\cap\mathcal K_j\ne\varnothing\).  Distinct connected
components have disjoint block-label sets.  Since \(\nu_S\) is a product
over blocks, both the integral in \eqref{B:eq:mixed-joint-factorization} and
the analogous design integral factor exactly over the connected components.
The common global factor \(\{\E Z(H)^n\}^{-1}\) cancels when conditioning on
\(X^n=x^n\).  Hence the conditional mixed laws of \((R_i,T_i)\) factor
exactly over the component graph, including when physical carrier boxes
overlap.

Consider a component of size \(q\le Q\).  Expand each of its finitely many
Walsh coefficients in block coefficient monomials.  At one observation the
degree in the aggregate propensity is at most four: \(p\) is linear,
\(K(p)\) is cubic, and \(pK(p)\) is quartic.  Therefore the degree in every
block is at most \(4q\le4Q\), so \cref{B:prop:carrier} applies.

For a block meeting \(q_k\) sites in the component,
\cref{B:prop:evaluated-projectivity} says that its lower-order moment vector
is exactly the convex interpolation, with coefficient
\(\overline\chi_{q_k}\), between the baseline moment and the ideal shifted
moment.  Multiplying over the finitely many blocks shows that each component
Walsh coefficient under the \(Q\)-mixture is a convex combination of the
partial-shift Walsh coefficients from \cref{B:lem:partial-stability}.  The
weight of the fully shifted term is the product of the block
\(\overline\chi_{q_k}\)'s.  Consequently, for every binary cell \(w\),
\begin{equation}
 \abs{p_{Q,C}(w\mid x_C)-p_{P,C}(w\mid x_C)}
 \le C_Q\Delta
 \sum_{k\in\mathcal K(C)}
 \{1-\overline\chi_{q_k}(T_k)\}.
 \label{B:eq:component-cell-difference}
\end{equation}
All fixed-parameter binary cells are bounded below by a positive constant,
so every \(q\)-site mixed cell is bounded below by a constant depending only
on \(Q\).  The Walsh transform is finite-dimensional and orthogonal.
Therefore
\begin{equation}
 H_C^2(x_C)
 \le C_Q\Delta^2
 \sum_{k\in\mathcal K(C)}
 \{1-\overline\chi_{q_k}(T_k)\}.
 \label{B:eq:component-hellinger-local}
\end{equation}
We used \((\sum b_k)^2\le C_Q\sum b_k\), valid because the number of blocks
in a \(q\)-component is bounded and \(0\le b_k\le1\).

\subsection{Geometry and low components}

The common mixed design is not necessarily iid, but every fixed
\(q\)-coordinate marginal density is at most \(\overline f^{\,q}\), because
conditionally on \(H\) the density is bounded by \(\overline f\).  A
labelled connected \(q\)-tuple has one point in the \(O(h^d)\) active region
and, after choosing a spanning tree, every other point in an \(O(r^d)\)
neighborhood of its parent.

We insert the ghost factor explicitly.  Fix a labelled connected
\(q\)-tuple, one of its incident blocks \(k_*\), and let
\(I_*\subseteq\{1,\ldots,q\}\) be the \(q_*\ge1\) sites lying in
\(S_{k_*}\).  There are \(O(h^dr^{-d})\) possible active choices of
\(k_*\).  In affine coordinates on \(S_{k_*}\), the integral of the
corresponding leakage factor satisfies
\begin{align}
 &\int_{S_{k_*}^{q_*}}
   \{1-\overline\chi_{q_*}(T_{k_*})\}
   \,dx_{I_*}
 \nonumber\\
 &\hspace{2cm}\le
 C_Qr^{dq_*}t^d\Lambda^{M_Q}.
 \label{B:eq:physical-ghost-integral}
\end{align}
Indeed, the affine Jacobian is \(C r^{dq_*}\), and
\eqref{B:eq:ghost-mass} applies after using the fixed positive lower bound on
\(D_{q_*}\).  Because the observation component is connected, the remaining
\(q-q_*\) sites can be attached to the set \(I_*\) by a labelled forest with
exactly \(q-q_*\) edges.  Each attachment forces the child into an
\(O(r^d)\) neighborhood of its parent.  Thus, for fixed \(k_*\), the
remaining integration contributes at most \(C_Qr^{d(q-q_*)}\).  Multiplying
by the number \(O(h^dr^{-d})\) of choices of \(k_*\) gives exactly
\[
 C_Qh^dr^{d(q-1)}t^d\Lambda^{M_Q}.
\]
The number of labelled forests and the number of incident blocks in a
\(q\)-component are bounded in terms of \(Q\), and the mixed-design
\(q\)-marginal is bounded by \(\overline f^{\,q}\).  Summing
\eqref{B:eq:component-hellinger-local} therefore gives, for
\(2\le q\le Q\),
\begin{align}
 \E\sum_{\substack{C:\,\abs C=q}}H_C^2
 &\le C_Qn^q h^dr^{d(q-1)}t^d\Delta^2\Lambda^{M_Q}
 \nonumber\\
 &=C_Q n^q h^dr^{d(q-2)}\ell^d\Delta^2\Lambda^{M_Q}
 \nonumber\\
 &=C_Q\{n^2h^d\ell^d\Delta^2\}
       (nr^d)^{q-2}\Lambda^{M_Q}
 =o(1).
 \label{B:eq:low-hellinger}
\end{align}
For singleton components, fix an incident block \(k_*\), sum over its
\(O(h^dr^{-d})\) possible locations, and use
\eqref{B:eq:physical-ghost-integral} with \(q_*=1\).  There is no forest edge,
so the same calculation gives
\begin{equation}
 \E\sum_{\substack{C:\,\abs C=1}}H_C^2
 \le C_Qnh^dt^d\Delta^2\Lambda^{M_Q}=o(1).
 \label{B:eq:single-hellinger}
\end{equation}
The final equalities use \cref{lem:ledgers} and \(nr^d=o(1)\).  Fixed
powers of \(\Lambda\) are absorbed by the strictly negative powers of
\(n\).

\subsection{High components}

Every component larger than \(Q\) contains a connected labelled subset of
size \(Q+1\).  The spanning-tree volume bound and the mixed-design upper
bound give
\begin{equation}
 \PP\{\text{some component has size }>Q\}
 \le C_Qn^{Q+1}h^dr^{dQ}=o(1)
 \label{B:eq:high-components}
\end{equation}
by \eqref{eq:high}.  These configurations are handled as the
exceptional event in Lemma~\ref{lem:components}; no sum of uncontrolled
high-component Hellinger distances is needed.

\subsection{Total variation and Le Cam}

Let \(\overline P_n,\overline Q_n\) denote the two prior mixtures.  Their
\(X^n\)-marginals coincide, and given the design the conditional laws
factor over the component graph; on the event of no high component every
component has at most \(Q\) sites.  \Cref{lem:components} with
\eqref{B:eq:low-hellinger}, \eqref{B:eq:single-hellinger}, and
\eqref{B:eq:high-components} gives
\[
 \TV(\overline P_n,\overline Q_n)
 \le \PP\{\text{a high component}\}
 +\Bigl\{\E\sum_{C:\,|C|\le Q} H_C^2\Bigr\}^{1/2}
 =o(1).
\]
In particular,
\begin{equation}
 \TV(\overline P_n,\overline Q_n)\le\tfrac12
 \label{B:eq:tv-bound}
\end{equation}
for all sufficiently large \(n\).

Every draw from the \(P\)-prior has target value
\(c_ac_b\Delta\), and every draw from the \(Q\)-prior has target value zero
(\cref{B:lem:legality}).  \Cref{lem:two-point} gives
\[
 \sup_{P\in\cP}\E_P\abs{\widehat\tau(x_0)-\tau_P(x_0)}
 \ge \frac{c_ac_b\Delta}{2}\{1-\TV(\overline P_n,\overline Q_n)\}
 \ge \frac{c_ac_b}{4}\,n^{-(\rho+\varepsilon)}.
\]
This proves \cref{thm:main} in case~(B) for all sufficiently small fixed
\(\varepsilon>0\).  For any prescribed \(\varepsilon>0\), choose a fixed
\(0<\varepsilon_0\le\varepsilon\) in that range and use
\(n^{-(\rho+\varepsilon_0)}\ge n^{-(\rho+\varepsilon)}\), changing the
constant if necessary.


\part{Case C: mixed smoothness}

\section{Case C: construction}\label{C:sec:construction}

Throughout Part~C exactly one of \(\alpha,\beta\) exceeds one:
case~(a) \(\alpha\le1<\beta\), case~(b) \(\beta\le1<\alpha\); thus
\(s_1=(1+\alpha\wedge\beta)/2\) and the scales are those of
Section~\ref{sec:scales}.  The degree
\[
 \dg=\begin{cases}1&\text{in case (a)},\\3&\text{in case (b)}\end{cases}
\]
is the degree of the likelihood cells in the carried coefficients and the
degree of the carrier algebra of Section~\ref{C:sec:carrier}.

\subsection{Amplitudes}
In both cases the rough nuisance has amplitude \(c\,\ell^{\text{(its
exponent)}}\), the smooth nuisance has amplitude \(c\,r^{\text{(its
exponent)}}\) and receives a virtual jitter of relative size \(t\):
\begin{equation}
\begin{array}{c|ccc|c}
 &\text{rough amplitude}&\text{smooth amplitude}&\text{jitter}&\\ \hline
 \text{(a) }\alpha\le1<\beta&\ja=c_a\ell^\alpha&b=c_br^\beta&\jb=bt=c_btr^\beta&\jsh=\jb,\ \jfi=\ja\\
 \text{(b) }\beta\le1<\alpha&\jb=c_b\ell^\beta&a=c_ar^\alpha&\ja=at=c_atr^\alpha&\jsh=\ja,\ \jfi=\jb
\end{array}
\label{C:eq:amplitudes}
\end{equation}
Here \(\jsh\) (the \emph{shift amplitude}) is the jitter of the smooth
nuisance and \(\jfi\) (the \emph{field amplitude}) the amplitude of the
rough one.  In both cases
\begin{equation}
 \ja\jb=c_ac_b\Delta,\qquad
 d_*(x):=\ja\jb G_h(x)^2=c_ac_b\Delta G_h(x)^2,\qquad
 \frac{\jsh}{\jfi}=\frac{c_\circ}{c_\bullet}\,t^{1-(\alpha\wedge\beta)}r^{\abs{\alpha-\beta}}\to0,
 \label{C:eq:jsh-over-jfi}
\end{equation}
where \(G_h(x)=G((x-x_0)/h)\) for a fixed \(G\in C_c^\infty\) with
\(G(0)=1\), \(0\le G\le1\); the first identity is
\(\ell^{\alpha\wedge\beta}\,t\,r^{\alpha\vee\beta}=t^{1+\alpha\wedge\beta}r^{A_0}=t^{2s_1}r^{A_0}=\Delta\),
and the last uses \(\alpha\vee\beta>1\ge\alpha\wedge\beta\): the factor
\(r^{\abs{\alpha-\beta}}\) tends to zero because \(\alpha\ne\beta\), and
\(t^{1-\alpha\wedge\beta}\le1\).  Hence, for large \(n\),
\begin{equation}
 \jsh\le\jfi,\qquad \jsh^2\le\jsh\jfi=c_ac_b\Delta .
 \label{C:eq:jsh-square}
\end{equation}

\subsection{Blocks and rough packets}
\emph{Blocks.}  Let \(\phi_{\rm lin}\in C_c^\infty(\R^d)\) be a
nonnegative partition-of-unity bump: \(\phi_{\rm lin}\ge0\),
\(\operatorname{supp}\phi_{\rm lin}\subseteq[-1,1]^d\), \(\phi_{\rm lin}>0\)
on \((-1,1)^d\), and \(\sum_{k\in\mathbb Z^d}\phi_{\rm lin}(u-k)\equiv1\)
(take \(\phi_{\rm lin}=\vartheta_0/\sum_k\vartheta_0(\cdot-k)\) for a
nonnegative \(\vartheta_0\in C_c^\infty\) positive on \((-1,1)^d\) and
supported in \([-1,1]^d\); the denominator is bounded below by
\(\vartheta_0\) at the nearest lattice point).  By compactness there is a
constant \(c_d\in(0,1]\), depending only on \(d\) and the fixed bump, with
\begin{equation}
 \phi_{\rm lin}\ge c_d\quad\text{on }[-3/4,3/4]^d .
 \label{C:eq:cd}
\end{equation}
Put \(x_k=x_0+rk\) and
\(\psi_k(x)=\phi_{\rm lin}((x-x_k)/r)\), so \(0\le\psi_k\le1\),
\(\sum_k\psi_k\equiv1\), and \(\psi_k\ge c_d\) on \(x_k+[-3r/4,3r/4]^d\).
Enlarge each support to a cube \(S_k\) of side \(Cr\), identified with
\(\Torus^d\) by an affine chart \(\theta_k\); \(\operatorname{supp}\psi_k\)
lies in a fixed compact subcube.  Let \(N_1\) be the maximal number of
enlarged cubes \(S_k\) containing a given point (a fixed number, \(\le C^d\)); a
fortiori at most \(N_1\) supports \(\operatorname{supp}\psi_k\) meet a point.
Active blocks are those meeting \(\operatorname{supp}G_h\); there are
\(O(h^dr^{-d})\).  Lagrange vectors \(c_{ki}(T)=c_i(T)\) for tuples in
block \(k\) are those of Lemma~\ref{lem:lagrange}, and the taper
is \(\chi_Q(T)=\prod_{i=1}^Q\vartheta(P_i(T)/t_n)\) with the fixed cutoff
\(\vartheta\) of Section~\ref{sec:torus} (\(\vartheta=0\) on \([0,1]\), \(=1\) on \([2,\infty)\)).

\emph{Rough packets.}  The rough packets \(\psi^\flat_m\), the rough field
\(\Lambda_\lambda\), the packet sets \(J(x)\), the constants \(N_0,R_0\), and
the moments \(M_j=\E\Lambda_\lambda^j\) (in particular \(M_2=G_h^2\),
\(M_4=\mu_4=G_h^4\{3-2q_4\}\)) are those of Section~\ref{sec:tools},
\eqref{eq:rough}--\eqref{eq:rough-moments}, at the fine scale \(\ell=tr\)
of \eqref{eq:fine}.  Case~(a) uses only \(M_0,M_1,M_2\); case~(b) also
uses \(M_3=0\) and \(M_4\).

\subsection{Assignment of sites to blocks}\label{C:sec:assign}
The block partition \(\{\psi_k\}\) is smooth with constant overlap, which
is required for the smoothness of the carried field below.  The carrier
of Section~\ref{C:sec:carrier}, however, needs every site to receive its
virtual jitter from a \emph{single} block, up to a set of sites of small
relative volume.  Let \(C_k=x_k+[-r/2,r/2)^d\) be the cubes of the lattice
\(\{x_k\}\), so that \(\{C_k\}\) partitions \(\R^d\); recall from
\eqref{C:eq:cd} that \(\psi_k\ge c_d\) on the enlarged cube
\(x_k+[-3r/4,3r/4]^d\).  Let
\begin{equation}
 w=t^d
 \label{C:eq:layer-width}
\end{equation}
be the relative transition width, and construct the assignment partition
as a tensor product of one-dimensional ones: let \(\varrho_w\) be a fixed
smooth probability density on \(\R\) supported in \([-w/2,w/2]\) (a fixed
bump rescaled to width \(w\)), \(\sigma=\varrho_w*\one_{[-1/2,1/2)}\), and
\begin{equation}
 \phi_k(x)=\prod_{j=1}^d\sigma\Bigl(\frac{x^{(j)}-x_k^{(j)}}r\Bigr),\qquad x\in\R^d .
 \label{C:eq:phi-def}
\end{equation}
Since \(\sum_{l\in\mathbb Z}\one_{[-1/2,1/2)}(y-l)\equiv1\), also
\(\sum_l\sigma(y-l)\equiv1\), hence \(\sum_k\phi_k\equiv1\); moreover
\(0\le\phi_k\le1\), \(\operatorname{supp}\phi_k\subseteq x_k+[-(1+w)r/2,(1+w)r/2]^d\subseteq x_k+[-3r/4,3r/4]^d\),
and \(\phi_k=1\) on \(\{x\in C_k:\operatorname{dist}_\infty(x,\partial C_k)\ge wr/2\}\).
Sites with \(\phi_k(x)=1\) for some \(k\) are \emph{interior} and are
\emph{assigned} to that block, \(k(x)=k\); the remaining set
\(\mathcal L=\{x:0<\phi_k(x)<1\text{ for some }k\}\) is the
\emph{transition layer}.  It is contained in the union over \(k\) of the
\(wr\)-neighbourhoods of \(\partial C_k\), so its relative volume in every
cube \(C_k\) is at most \(2dw=2dt^d\).  Put
\begin{equation}
 m_k(x)=\frac{\phi_k(x)}{\psi_k(x)}\one\{\phi_k(x)>0\},\qquad 0\le m_k\le c_d^{-1} ,
 \label{C:eq:multiplier}
\end{equation}
which is smooth because \(\psi_k\ge c_d\) on \(\operatorname{supp}\phi_k\)
by \eqref{C:eq:cd}; \(c_d^{-1}\) is a fixed constant and enters the bounds
below only through the constant \(N_m\) of \eqref{C:eq:Nm}.

\begin{lemma}[Wiener norm of the multiplier]\label{C:lem:multiplier}
In the chart of block \(k\), \(\norm{m_k\circ\theta_k^{-1}}_{\cW_d}\le C_d\{1+\log(1/w)\}^d\le C_d\Lambda^d\).
\end{lemma}
\begin{proof}
\emph{One-dimensional step.}  On the circle of length one identify
\(\sigma(\cdot)\) (after the affine chart, the cube \(S_k\) of side \(Cr\)
becomes \(\Torus^d\) and \(x\mapsto(x-x_k)/r\) becomes a fixed dilation, so
we may work with \(\sigma\) on an interval of fixed relative length and
transition width \(w'=w/C\)).  Its Fourier coefficients are
\(\widehat\sigma(\nu)=\widehat{\one}(\nu)\widehat{\varrho}(\nu)\) with
\(\abs{\widehat{\one}(\nu)}\le C/\abs\nu\) (\(\nu\ne0\)) and, by smoothness of
the bump, \(\abs{\widehat\varrho(\nu)}\le C_K(1+w'\abs\nu)^{-K}\) for every
\(K\).  Hence
\[
 \sum_\nu\abs{\widehat\sigma(\nu)}\le C+\sum_{1\le\abs\nu\le1/w'}\frac C{\abs\nu}
 +\sum_{\abs\nu>1/w'}\frac{C_K}{\abs\nu(w'\abs\nu)^K}\le C\{1+\log(1/w')\}.
\]
\emph{Tensor product.}  In the chart, \(\phi_k\circ\theta_k^{-1}\) is the
product over the \(d\) coordinates of such one-dimensional functions, and
the Wiener norm of a tensor product is the product of the Wiener norms;
so \(\norm{\phi_k\circ\theta_k^{-1}}_{\cW_d}\le C_d\{1+\log(1/w)\}^d\).
\emph{The factor \(1/\psi_k\).}  Let \(\zeta\) be a fixed smooth function
on \(\Torus^d\) equal to \(1/(\psi_k\circ\theta_k^{-1})\) on
\(\{\psi_k\circ\theta_k^{-1}\ge c_d\}\) (which contains the support of
\(\phi_k\circ\theta_k^{-1}\)); \(\zeta\) does not depend on \(k\) or \(n\)
because \(\psi_k\circ\theta_k^{-1}\) is the same fixed bump for every block,
so \(\norm\zeta_{\cW_d}\le C\).  Then \(m_k\circ\theta_k^{-1}=(\phi_k\circ\theta_k^{-1})\zeta\)
and the algebra property of \(\cW_d\) gives the claim.  Finally, by
\eqref{eq:fine},
\[
 \log(1/w)=d\log(1/t)=\frac{d(z-\xi A_0)}{2s_1}\log n=O(\log n)=O(\Lambda).
\]
\end{proof}

\subsection{Fields, virtual shift, and the two hypotheses}\label{C:sec:fields}
Only active blocks carry coefficients: \(U=(U_k)_{k\ \rm active}\),
\(U_k\in\R^M\), and the \emph{legal ball} is \(\{\norm{U_k}\le R_U\ \forall k\}\)
for a fixed \(R_U\) with \(2\mathbb B_0\subseteq\{\norm U\le R_U\}\), where
\(\mathbb B_0\) is the support ball of \(\Pi_0\).  The \emph{carried
field} (unit amplitude) is
\begin{equation}
 \cf(x)=\sum_{k\ \rm active}\psi_k(x)v_Q(\theta_k(x))^\top U_k .
 \label{C:eq:carried}
\end{equation}
Let \(T_k=(x_1,\dots,x_Q)\) be a full \(Q\)-tuple in block \(k\), with
\(\chi_Q(\theta_k(T_k))>0\); the Lagrange vectors below are evaluated at
its chart coordinates.  The \emph{virtual shift} of block \(k\) is
\begin{equation}
 \widetilde U_k=U_k+t\sum_{i=1}^Qm_k(x_i)\Lambda_\lambda(x_i)\,c_{ki}(T_k)
 +(\text{moment corrections, Section~\ref{C:sec:target}}).
 \label{C:eq:shift}
\end{equation}
For \(q<Q\) retained sites, the remaining \(Q-q\) sites are auxiliary
integration variables, removed after likelihood evaluation as in
Section~\ref{C:sec:projectivity}.  In particular, \(\chi_Q\) is never
applied to a tuple of length less than \(Q\).
\(\widetilde U_k\) is a bookkeeping expression, not a random variable of
the construction: the actual parameters below contain only \(U_k\), and
\eqref{C:eq:shift} records which moments the carrier reproduces.
By \(v_Q(\theta_k(x_j))^\top c_{ki}(T_k)=\one\{i=j\}\), the first sum
changes \(\cf(x_j)\) by
\[
 \psi_k(x_j)\,t\,m_k(x_j)\Lambda_\lambda(x_j)=t\phi_k(x_j)\Lambda_\lambda(x_j);
\]
summing over the blocks meeting \(x_j\) and using \(\sum_k\phi_k=1\), the
total change of \(\cf(x_j)\) at a fully shifted site is
\(t\Lambda_\lambda(x_j)\), coming from the assigned block alone when the
site is interior.  The fields and the two hypotheses are, in coded
coordinates:

\emph{Case (a)} (\(\alpha\le1<\beta\)): \(p_\lambda=\ja\Lambda_\lambda\),
\(y_0=b\,\cf\), site jitter \(\delta=bt\Lambda_\lambda=\jb\Lambda_\lambda\), and
\begin{equation}
\begin{array}{c|ccc}
 &p=2\pi-1&y=2\mu_0-1&\tau\\ \hline
 P&p_\lambda&y_0-d_*&d_*\\
 Q&p_\lambda&y_0&0\\ \hline
 Q\text{ (virtual)}&p_\lambda&y_0+\delta&0
\end{array}
\label{C:eq:two-sides-a}
\end{equation}

\emph{Case (b)} (\(\beta\le1<\alpha\)): \(p_0=a\,\cf\), site jitter
\(\delta=at\Lambda_\lambda=\ja\Lambda_\lambda\),
\(y_\lambda(\cdot;p)=\jb\Lambda_\lambda\{1+\eta-p^2\}\) with the
fourth-moment correction
\begin{equation}
 \eta(x):=\ja^2G_h(x)^2\{3-2q_4(x)\}\ge0,
 \qquad
 \norm\eta_\infty\le3\ja^2,\quad \norm\eta_{C^\beta}\le C\ja^2\ell^{-\beta},\quad
 d_*\eta=\ja^3\jb\mu_4 ,
 \label{C:eq:eta}
\end{equation}
defined for every \(x\) and equal to \(\ja^2\mu_4(x)/G_h(x)^2\) wherever
\(G_h(x)\ne0\) (by \eqref{eq:rough-moments}); the last identity holds
everywhere, both sides vanishing where \(G_h=0\),
and
\begin{equation}
\begin{array}{c|ccc}
 &p=2\pi-1&y=2\mu_0-1&\tau\\ \hline
 P&p_0&y_\lambda(\cdot;p_0)-d_*(1+3p_0)&d_*\\
 Q&p_0&y_\lambda(\cdot;p_0)&0\\ \hline
 Q\text{ (virtual)}&p_0+\delta&y_\lambda(\cdot;p_0+\delta)&0
\end{array}
\label{C:eq:two-sides-b}
\end{equation}

In both tables the rows \(P\) and \(Q\) are the \emph{actual}
parameters: under both hypotheses the coefficients \(U_k\) of the carried
field are drawn from the laws of Proposition~\ref{C:prop:carrier}, the
rough signs from the same Rademacher prior, and no parameter contains
\(\delta\).  The row ``\(Q\) (virtual)'' is not a parameter of any law in
the mixture; it is the \emph{jittered} parameter used in the single-site
identities of Section~\ref{C:sec:bridge} and in the target of
Section~\ref{C:sec:target}.  Its role is the following.  The two carrier
laws differ only in the law of \(U_k\) given the design, and
Proposition~\ref{C:prop:carrier}(iv) states that this difference, tested
against any likelihood cell, equals the tapered, design-weighted moment
increment \(\cJ(\mathbb D)\) built from the jittered parameter.  This is a
moment identity with a taper and with the moment corrections of
Section~\ref{C:sec:target}; it does not say that the shifted coefficients
\eqref{C:eq:shift} are the values of a random vector under \(Q\).  In case~(b)
the \(Q\)-side outcome is a deterministic quadratic function of the
propensity and of the rough field: this is the cubic bridge of Part~B with the
rough field in place of one power of the propensity.

\subsection{The single-site identities}\label{C:sec:bridge}
Fix a site \(x\) and write \(\Lambda=\Lambda_\lambda(x)\), \(G=G_h(x)\),
\(d=d_*(x)\), \(\eta=\eta(x)\), \(M_j=\E\Lambda^j\).  The \emph{shift
variable} \(\zeta\) is the coded value of the smooth nuisance at the
site, \(\zeta=y\) in case~(a) and \(\zeta=p\) in case~(b); the unshifted
\(Q\)-likelihood, as a function of \(\zeta\) for a cell
\((R,T)\in\{\pm1\}^2\), and the \emph{real field} are
\begin{equation}
\begin{array}{c|cc}
 &L(\zeta)&\real(\zeta)\\ \hline
 \text{(a)}&(1+R\ja\Lambda)(1+T\zeta)&d\ja\Lambda\,T+d\,RT\\[2pt]
 \text{(b)}&(1+R\zeta)\bigl(1+T\jb\Lambda(1+\eta-\zeta^2)\bigr)&-2d\zeta\,T+d(1-3\zeta^2)\,RT
\end{array}
\label{C:eq:L-real}
\end{equation}
\(L\) is a polynomial of degree \(\dg\) in \(\zeta\) and affine in
\(\Lambda\), \(L=A+B\Lambda\) with \(B\) carrying the factor \(\jfi\).  By
\eqref{eq:walsh}, \(L(\zeta_0)\) is the \(Q\)-likelihood at the unshifted
parameter (\(\zeta_0=y_0\), resp.\ \(p_0\)), \(L(\zeta_0+\jsh\Lambda)\) the
\(Q\)-likelihood at the jittered parameter, and the \(P\)-likelihood is
\begin{equation}
 L^P=L(\zeta_0)+\real(\zeta_0):
 \label{C:eq:LP}
\end{equation}
in case~(a), \(y+\tau(1+p)=y_0+d\ja\Lambda\) and \(py+\tau(1+p)=\ja\Lambda y_0+d\);
in case~(b), \(y+\tau(1+p)=y_\lambda-d(1+3p_0)+d(1+p_0)=y_\lambda-2dp_0\) and
\(py+\tau(1+p)=p_0y_\lambda-dp_0(1+3p_0)+d(1+p_0)=p_0y_\lambda+d(1-3p_0^2)\).
Define the \emph{increments}
\begin{equation}
 \Delta^{(e)}(\zeta)=\frac{\jsh^e}{e!}\,\partial_\zeta^eL(\zeta),\quad e=1,\dots,\dg,
 \qquad\text{so that}\qquad
 L(\zeta+\jsh\Lambda)-L(\zeta)=\sum_{e=1}^{\dg}\Lambda^e\Delta^{(e)}(\zeta),
 \label{C:eq:increments}
\end{equation}
and the \emph{substitution polynomials} in \(\Lambda\), for
\(e\in\{1,\dots,\dg\}\), \(f\in\{0,\dots,\dg\}\):
\begin{equation}
\begin{array}{c|l}
 \text{(a)}&\K^{(1,0)}(\Lambda)=\Lambda,\qquad \K^{(1,1)}(\Lambda)=\varkappa:=\ja^2G^4
 \quad(\text{so that }\varkappa\jb=d\ja G^2),\\[2pt]
 \text{(b)}&\K^{(e,f)}(\Lambda)=M_f\Lambda^e:\quad
 \K^{(e,0)}=\Lambda^e,\quad \K^{(e,f)}=0\ (f\text{ odd}),\quad \K^{(e,2)}=G^2\Lambda^e .
\end{array}
\label{C:eq:K}
\end{equation}
In both cases every nonzero \(\K^{(e,f)}\) is a Walsh polynomial all of
whose monomials have degree \(\equiv e+f\pmod2\) (in case~(a):
\(\Lambda\) for \((1,0)\), the constant \(\varkappa\) for \((1,1)\)), and
\(\norm{\K^{(e,f)}}_W\le N_0^{\dg}\).

\begin{lemma}[Bridge with design exponents]\label{C:lem:bridge}
In either case, for every real \(\zeta\), every cell \((R,T)\), and every
\(f\in\{0,\dots,\dg\}\),
\begin{equation}
 \sum_{e=1}^{\dg}\E_\Lambda\bigl[\K^{(e,f)}(\Lambda)\,\Delta^{(e)}(\zeta)\bigr]
 =\E_\Lambda\bigl[\Lambda^f\,\real(\zeta)\bigr],
 \label{C:eq:bridge}
\end{equation}
where \(\E_\Lambda\) averages the rough field at the site (the other
quantities being fixed).  In particular (\(f=0\))
\(\E_\Lambda L(\zeta+\jsh\Lambda)=\E_\Lambda L(\zeta)+\E_\Lambda\real(\zeta)\).
\end{lemma}

\begin{proof}
\emph{Case (a).}  \(\Delta^{(1)}(y)=\jb\partial_yL=\jb T(1+R\ja\Lambda)\).
For \(f=0\): \(\E[\Lambda\jb T(1+R\ja\Lambda)]=\ja\jb G^2RT=dRT\), and
\(\E[\real]=\E[d\ja\Lambda T+dRT]=dRT\).  For \(f=1\):
\(\E[\varkappa\jb T(1+R\ja\Lambda)]=\varkappa\jb T=d\ja G^2T\), and
\(\E[\Lambda\real]=d\ja G^2T\).

\emph{Case (b).}  Write \(c=T\jb\Lambda\), so that
\(L(p)=1+c(1+\eta)+Rp+Rpc(1+\eta)-cp^2-Rcp^3\) and
\[
 \partial_pL=R+Rc(1+\eta)-2cp-3Rcp^2,\qquad
 \partial_p^2L=-2c-6Rcp,\qquad
 \partial_p^3L=-6Rc .
\]
With \(\E[\Lambda^jc]=M_{j+1}\jb T\), for every \(j\ge0\):
\begin{align}
 \E[\Lambda^j\Delta^{(1)}(p)]&=\ja\bigl[M_jR+M_{j+1}\jb T\{R(1+\eta)-2p-3Rp^2\}\bigr],
 \label{C:eq:E1}\\
 \E[\Lambda^j\Delta^{(2)}(p)]&=-\ja^2M_{j+1}\jb T(1+3Rp),
 \label{C:eq:E2}\\
 \E[\Lambda^j\Delta^{(3)}(p)]&=-\ja^3M_{j+1}\jb RT .
 \label{C:eq:E3}
\end{align}
For \(f=0\) the left side of \eqref{C:eq:bridge} is
\(\sum_e\E[\Lambda^e\Delta^{(e)}]\).  By \eqref{C:eq:E1} with \(j=1\)
(\(M_1=0\), \(M_2=G^2\)): \(\ja\jb G^2T\{R(1+\eta)-2p-3Rp^2\}
=d\{RT(1+\eta)-2pT-3RTp^2\}\).  By \eqref{C:eq:E2} with \(j=2\): \(M_3=0\),
so the term vanishes.  By \eqref{C:eq:E3} with \(j=3\):
\(-\ja^3\mu_4\jb RT=-d\eta RT\) by \eqref{C:eq:eta}.  The sum is
\(d(1-3p^2)RT-2dpT=\real(p)=\E[\Lambda^0\real(p)]\) (\(\real\) does not
contain \(\Lambda\)).  For \(f\ge1\), \(\K^{(e,f)}=M_f\Lambda^e\), so the
left side is \(M_f\) times the left side for \(f=0\), i.e.\
\(M_f\real(p)=\E[\Lambda^f\real(p)]\).
\end{proof}

The identity with \(f=0\) is the idea behind the construction: at a
single site, averaging the rough field, the jittered \(Q\)-likelihood
equals the \(P\)-likelihood; in case~(a) the \(Q\)-side \(RT\)-cell
produces \(d\) through \(\E[\ja\Lambda\cdot\jb\Lambda]=\ja\jb G^2\), in case~(b)
the fourth-moment correction \(\eta\) cancels the cubic term of the
jitter and the \(P\)-side control outcome absorbs the \(T\)-cell
\(-2dp_0\), which is legal because \(p_0\) is smooth.  The identities
with \(f\ge1\) are what the design weights of the carrier require
(Proposition~\ref{C:prop:matching}, Step~3): when the design density
carries a factor \(\Lambda^f\) at the site, the jitter \(\Lambda^e\) of
the site must be replaced by \(\K^{(e,f)}(\Lambda)\), and
Lemma~\ref{C:lem:bridge} says that this reproduces the design-weighted
\(P\)-side \(\E[\Lambda^f\real]\) exactly.  In case~(a) the real field
contains \(\Lambda\) and the correction replaces the jitter by the
constant \(\varkappa\); in case~(b) it does not, and the correction is a
rescaling of the jitter by \(\E[\Lambda^f]\).

\section{Case C: legality}\label{C:sec:legal}

\begin{lemma}\label{C:lem:legal}
For sufficiently small fixed \(c_a,c_b>0\), all sufficiently large \(n\),
every \(\lambda\), every \(U\) in the legal ball, and every design density
produced by the carrier, the two actual rows \(P\) and \(Q\) of
\eqref{C:eq:two-sides-a}, resp.\ \eqref{C:eq:two-sides-b}, define laws in
\(\cP\), and
\(\tau_P(x_0)-\tau_Q(x_0)=c_ac_b\Delta\).
\end{lemma}

\begin{proof}
\emph{The rough field.}  \(\Lambda_\lambda=G_h\cdot\sum_m\lambda_m
\phi((\cdot-x_0)/\ell-m)\) is \(G_h\) times a bounded-overlap sum of
translates of \(\phi(\cdot/\ell)\).  By Lemma~\ref{lem:holder-scaling}
with \(t=\alpha\wedge\beta\le1\), \(\norm{\Lambda_\lambda}_{C^{\alpha\wedge\beta}}\le C\ell^{-\alpha\wedge\beta}
+Ch^{-\alpha\wedge\beta}\le C\ell^{-\alpha\wedge\beta}\), and \(\abs{\Lambda_\lambda}\le N_0\).

\emph{The carried field.}  \(\cf=\sum_k\psi_kv_Q(\theta_k)^\top U_k\) is a
bounded-overlap sum of functions smooth at scale \(r\) with bounded
coefficients, so by Lemma~\ref{lem:holder-scaling} (any \(t>0\))
\(\norm{\cf}_{C^{\alpha\vee\beta}}\le C_Qr^{-\alpha\vee\beta}\) and \(\abs{\cf}\le C_Q\).
The carrier draws every \(U_k\) inside the fixed legal ball by
Proposition~\ref{C:prop:carrier}(iii), and legality of the actual
parameters involves only \(U_k\).  (The shifted coefficients
\eqref{C:eq:shift} are never drawn and are not parameters of any law.
The jittered parameter is used only in the single-site identities and in
Lemma~\ref{C:lem:walsh-bounds}, where the jitter is the first-order term
of \eqref{C:eq:shift} alone, of size
\(\norm{t\sum_im_k(x_i)c_{ki}\Lambda_i}\le C_QN_0t/t_n=C_Q\Lambda^{-B}\to0\) on
the support of \(\chi_Q\).  The base point there is
\(U\in\mathbb B_0=\operatorname{supp}\Pi_0\), and \(2\mathbb B_0\) lies
inside the legal ball, so the fixed margin between \(\mathbb B_0\) and
the boundary of the legal ball absorbs the shift for large \(n\); the
statement is not made for arbitrary \(U\) on the boundary of the legal
ball.  The moment corrections of
Section~\ref{C:sec:target} are terms of the algebra element
\(\cJ(\mathbb D)\), controlled through \(\norm{\mathbb D-1^{\otimes Q}}_\cA\),
and are not shifts of any parameter.)

\emph{Case (a).}  Propensity: \(\norm{p_\lambda}_{C^\alpha}\le Cc_a\) and
\(\abs{p_\lambda}\le N_0c_a\); choose \(c_a\) so that \(\pi=(1+p_\lambda)/2\)
lies in \(\cH^\alpha(L_\pi)\) (the constant \(1/2\) is an interior point)
and \(N_0c_a\le1-2\kappa\).  Outcome, \(Q\) side: \(y_0=b\cf\) has
\(\norm{y_0}_{C^\beta}\le C_Qbr^{-\beta}=C_Qc_b\), \(\abs{y_0}\le C_Qc_b\).
Outcome, \(P\) side: \(y_P=y_0-d_*\), and
\(\norm{d_*}_{C^\beta}\le C\Delta h^{-\beta}=Ch^{\gamma-\beta}\to0\) since
\(\gamma>A_0\ge\beta\).

\emph{Case (b).}  Propensity: \(\norm{p_0}_{C^\alpha}\le C_Qar^{-\alpha}=C_Qc_a\),
\(\abs{p_0}\le C_Qc_a\); choose \(c_a\) accordingly.  Outcome, \(Q\)
side: \(y=\jb\Lambda_\lambda(1+\eta-p_0^2)\); by \eqref{C:eq:eta} and the
bound on \(p_0\), \(\norm{1+\eta-p_0^2}_\infty\le2\) and
\(\norm{1+\eta-p_0^2}_{C^\beta}\le C\ja^2\ell^{-\beta}+C_Qc_a^2r^{-\beta}\le C\ell^{-\beta}\),
so \(\norm y_{C^\beta}\le C_Q\jb\ell^{-\beta}=C_Qc_b\) and \(\abs y\le2N_0c_b\).
Outcome, \(P\) side: \(y_P=y-d_*(1+3p_0)\), and
\(\norm{d_*(1+3p_0)}_{C^\beta}\le C\Delta h^{-\beta}+C_Q\Delta ar^{-\beta}
=Ch^{\gamma-\beta}+C_Qc_a\Delta r^{\alpha-\beta}\to0\), since
\(\gamma>A_0>\beta\) and \(\alpha>\beta\).

\emph{CATE, ranges, design, separation (both cases).}
\(\tau_P=d_*=c_ac_b\Delta G_h^2\) has
\(\norm{\tau_P}_{C^\gamma}\le Cc_ac_b\Delta h^{-\gamma}=Cc_ac_b\);
\(\tau_Q=0\).  \(\pi=(1+p)/2\), \(\mu_0=(1+y)/2\) and \(\mu_1=\mu_0+\tau\) lie
in \([\kappa,1-\kappa]\) because \(\abs p\le C_Qc_a\) and
\(\abs y+\abs\tau\le C_Qc_b+o(1)\), both \(<1-2\kappa\) after choosing
\(c_a,c_b\).  The
design density is \(f_H\) of Section~\ref{C:sec:carrier}, which lies in
\([\underline f,\overline f]\) by \eqref{C:eq:density-bounds}, the only
requirement that \(\cP\) places on the design.  Finally
\(\tau_P(x_0)-\tau_Q(x_0)=c_ac_b\Delta G(0)^2=c_ac_b\Delta\).
\end{proof}

\section{Case C: the carrier with exact design-weighted matching}\label{C:sec:carrier}

Fix an active block \(k\).  The chart \(\theta_k\) identifies \(S_k\) with
\(\Torus^d\); functions on \(S_k\) become functions on the torus only after
a cutoff, since a rough packet or the bump \(G_h\) may cross \(\partial S_k\)
and its restriction to \(S_k\), read periodically, would have a jump there.
Let \(\chi_0\in C_c^\infty\) on the open chart cube be a fixed cutoff with
\(0\le\chi_0\le1\), equal
to one on the fixed subcube containing \(\operatorname{supp}(\phi_k\circ\theta_k^{-1})\)
(hence on \(\operatorname{supp}m\), with \(m\) as below) and vanishing near
the boundary of the cube; \(\chi_0\) is the same function for every block.
In the chart write \(m(u)=m_k(\theta_k^{-1}(u))\) and
\[
 \Lambda(u)=\chi_0(u)\,\Lambda_\lambda(\theta_k^{-1}(u)),\qquad
 \psi_m(u)=\chi_0(u)\,\psi^\flat_m(\theta_k^{-1}(u)),\qquad
 G(u)=G_h(\theta_k^{-1}(u)),
\]
so that \(\Lambda(u)=\sum_{m\in\cJ_k}\lambda_m\psi_m(u)\), where
\(\cJ_k=\{m:\operatorname{supp}\psi^\flat_m\cap S_k\ne\varnothing\}\) is the set
of packets \emph{read by block \(k\)}.  Since \(\chi_0=1\) on
\(\operatorname{supp}m\), \(m\Lambda\) and \(m\psi_m\) are the chart
expressions of \(m_k\Lambda_\lambda\) and \(m_k\psi^\flat_m\), so the cutoff
does not change any quantity that carries the multiplier \(m\); it only
makes the packets smooth on the torus.  Then \(\abs{\cJ_k}\le Ct^{-d}\), at
most \(N_0\) of the \(\psi_m\) are nonzero at any point, and each
\(\psi_m\) is a smooth function supported in a cube of side \(\asymp t\)
with
\[
 \norm{\psi_m}_{\cW_d}\le C_\psi,\qquad \norm{\psi_m}_{L^1(\Torus^d)}\le Ct^d :
\]
by \eqref{eq:rough}, \(\psi^\flat_m(x)=G_h(x)\phi((x-x^\flat_m)/\ell)\) with
\(x^\flat_m=x_0+\ell m\), and the chart is affine with \(\ell\mapsto t'\asymp t\),
so \(\psi_m(u)=b_{k,m,n}((u-u_m)/t')\) with the scaled profile
\[
 b_{k,m,n}(v)=\chi_0(u_m+t'v)\,G_h(x^\flat_m+\ell v)\,\phi(v),\qquad
 u_m=\theta_k(x^\flat_m).
\]
This profile is not a fixed function, but the family is supported in
\(\operatorname{supp}\phi\) and has \(C^{d+1}\) norms bounded uniformly in
\(k\), \(m\) and \(n\): \(\chi_0\) is fixed and \(t'\le1\), and
\(\partial_v^\nu\{G_h(x^\flat_m+\ell v)\}=(\ell/h)^{\abs\nu}(\partial^\nu G)((x^\flat_m-x_0+\ell v)/h)\)
with \(\ell/h\le1\).  The Wiener norm of a dilate
\(v\mapsto b((u-u_m)/t')\) of a bump \(b\) is bounded by a constant times
the \(C^{d+1}\) norm of \(b\), uniformly in the dilation, which gives the
first bound; the second is \(\abs{\psi_m}\le\abs{\phi(\cdot/t')}\) with
\(0\le\chi_0\le1\) and \(\abs{G_h}\le1\).  \(G\) is not cut off here:
it enters the algebra only through the products \(\chi_0G^2\), \(\chi_0G^4\)
of \eqref{C:eq:coef-chart}.  All constants are uniform in \(k\) and in the
signs.  The likelihood cells are polynomials of degree
\(\le\dg\) in each site's evaluation of the carried field, so the carrier
must reproduce moments of degree \(\le\dg\) per site; accordingly the
algebra, the ideal increment, the polarization and the target below are
of degree \(\le\dg\) per slot (\(\dg=1\): multilinear).

\subsection{The \texorpdfstring{\(m\)}{m}-adapted algebra of degree \texorpdfstring{\(\dg\)}{d}}
Let \(X_1,\dots,X_Q\) be commuting formal symbols and let
\(\{\lambda_m\}_{m\in\cJ_k}\) be a second family of commuting symbols with
\(\lambda_m^2=1\); monomials are \(\lambda^S=\prod_{m\in S}\lambda_m\),
\(S\subseteq\cJ_k\), with \(\lambda^S\lambda^{S'}=\lambda^{S\triangle S'}\).
Put \(\cE=\{0,1,\dots,\dg\}^Q\) and
\begin{equation}
 \begin{gathered}
 N_m:=c_d^{-1}N_0\max\{1,\norm{m}_{\cW_d}\}
       \le C_dN_0\Lambda^d
       \qquad(\text{Lemma~\ref{C:lem:multiplier}}),\\
 \abs{m(u)\Lambda(u)}\le c_d^{-1}N_0\le N_m .
 \end{gathered}
 \label{C:eq:Nm}
\end{equation}
An element of \(\cA\) (the \emph{\(m\)-adapted representative algebra}) is
an array \(W=(W_{e,S})\), \(e\in\cE\), \(S\subseteq\cJ_k\), with
\(W_{e,S}\in\cW_{dQ}\).  Its \emph{realization} as a polynomial in the
bare symbols \(X_i\), and its norm, are
\begin{equation}
 \begin{aligned}
  \realize(W)(T;X,\lambda)
   &=\sum_{e\in\cE}\sum_S m^e(T)W_{e,S}(T)X^e\lambda^S,
   &m^e(T)&:=\prod_i m(u_i)^{e_i},\\
  \norm W_\cA
   &=\sum_{e,S}\norm{W_{e,S}}_{\cW_{dQ}}N_m^{\abs e}<\infty,
   &\abs e&:=\sum_i e_i .
 \end{aligned}
 \label{C:eq:A-norm}
\end{equation}
Equality and the norm in \(\cA\) are defined on the arrays, not on their
realizations.  No quotient by the kernel of \(\realize\) is taken.
In particular, \(\realize\) need not be injective: a nonzero
representative supported where \(m^e=0\) can have zero realization.
Thus \(\cA\) is not identified with a subspace of bare-coefficient
polynomials.  The notation \(m^eW_{e,S}X^e\lambda^S\) below can denote a
\emph{tagged representative term}; when read as an ordinary polynomial,
it denotes its realization only.

Addition is coefficientwise.  Multiplication is convolution of the
representatives, with \(\lambda^S\lambda^{S'}=\lambda^{S\triangle S'}\),
followed by the projection \(P_\dg\) discarding exponents greater than
\(\dg\) in any slot.  Thus the product of tagged terms uses
representative \(W_{e,S}V_{e',S'}\) at \((e+e',S\triangle S')\), before
projection.  Vector-valued elements carry the sum of the coordinate
norms.  All polarization identities and the fixed point are formed in
this representative space.

The \emph{evaluation} is
\(\ev_\lambda W=\realize(W)(T;\Lambda_\lambda(u_\cdot),\lambda)\).
The \emph{reflection} is the isometric involution
\((RW)_{e,S}=(-1)^{\abs e+\abs S}W_{e,S}\), corresponding after realization
to \(X_i\mapsto-X_i\), \(\lambda_m\mapsto-\lambda_m\).  Hence
\(\ev_{-\lambda}W=\ev_\lambda(RW)\).  For a symbol \(m\in\cJ_k\), the
\emph{symbol seminorm} is
\begin{equation}
 \norm W_{[m]}=\sum_e\sum_{S\ni m}
       \norm{W_{e,S}}_{\cW_{dQ}}N_m^{\abs e},
 \label{C:eq:symbol-seminorm}
\end{equation}
the mass of the representative terms containing the explicit symbol
\(\lambda_m\), as distinct from dependence entering through site values.

\begin{lemma}\label{C:lem:algebra}
\(\cA\) is a commutative Banach algebra with unit; \(R\) is an isometric
involution; moreover
\[
 \sup_T\norm{\lambda\mapsto\ev_\lambda W(T)}_W\le\norm W_\cA,
 \qquad \norm{\ev_\lambda W}_\infty\le\norm W_\cA
 \quad\text{for every }\lambda.
\]
Also \(\ev_\lambda(WV)=\ev_\lambda(W)\ev_\lambda(V)\) whenever the
projection \(P_\dg\) discards no term, i.e.\ whenever every monomial of
\(WV\) has exponents \(\le\dg\) (this is the case in every product used
below: the slot factors of an atom power
\(F^{\otimes Q}=\prod_s(1+\theta t_rg(u_s))\) involve only \(X_s\);
the ambient expression in \eqref{C:eq:target} also has degree
\(\le\dg\) per slot, and its representative lift is defined explicitly in
Section~\ref{C:sec:target}; and the
polarization coefficients \(a_m(A)\) are \(X\)-free); multiplying by a
scalar Wiener element \(w(u_i)\) of one slot has norm
\(\le\norm w_{\cW_d}\); and for every symbol \(m\),
\[
 \norm{WV}_{[m]}\le\norm W_{[m]}\norm V_\cA+\norm W_\cA\norm V_{[m]} .
\]
Moreover, in the \emph{frozen-field evaluation}
\(W(\lambda;\xi)=\sum_{e,S}m^e(T)W_{e,S}(T)\xi^e\lambda^S\), in which the site
symbols \(X_i\) are replaced by fixed numbers \(\xi_i\) with
\(\abs{m(u_i)\xi_i}\le N_m\) (for instance \(\xi_i=\Lambda_{\lambda_0}(u_i)\) for
a fixed sign vector \(\lambda_0\)), the Walsh monomials containing
\(\lambda_m\) have total mass \(\le\norm W_{[m]}\) uniformly in \(T\), so
that \(\sup_T\abs{W(\lambda;\xi)-W(\lambda^{(m)};\xi)}\le2\norm W_{[m]}\),
where \(\lambda^{(m)}\) is \(\lambda\) with the sign of \(\lambda_m\) reversed.
The full evaluation \(\ev_\lambda W=W(\lambda;\Lambda_\lambda(u_\cdot))\)
depends on \(\lambda_m\) in addition through the site values
\(\Lambda_\lambda(u_i)\) with \(m\in J(u_i)\); this dependence is not
controlled by the seminorm and is treated separately where it matters
(Lemma~\ref{C:lem:weights}).
\end{lemma}
\begin{proof}
Let \(\cA^{\rm full}\) be the same construction with \(e\in\mathbb N^Q\)
(and \(m^e=\prod m(u_i)^{e_i}\)); it is a commutative Banach algebra with
unit by submultiplicativity of \(\cW_{dQ}\), \(N_m^{\abs{e'}}N_m^{\abs{e''}}=N_m^{\abs{e'+e''}}\),
\(\norm{\lambda^S\lambda^{S'}}=1\), and completeness of weighted
\(\ell^1\)-sums.  \(\norm{P_\dg}\le1\).  The product on \(\cA\) is
\(W\cdot V=P_\dg(WV)\); it is commutative, bilinear, submultiplicative, and
associative because exponents add: a monomial of \(WVU\) exceeds degree
\(\dg\) in a slot iff one of the partial products already does or the
total does, so \(P_\dg(P_\dg(WV)U)=P_\dg(WVU)=P_\dg(WP_\dg(VU))\).
For the Walsh-norm evaluation bound, at each fixed tuple \(T\),
\(\norm{\Lambda_\lambda(u_i)}_W\le N_0\).  The Walsh algebra inequality
therefore gives
\[
 \norm{\ev_\lambda W(T)}_W
 \le\sum_{e,S}\abs{W_{e,S}(T)}(c_d^{-1}N_0)^{\abs e}
 \le\norm W_\cA .
\]
For the pointwise bound, \(\norm{W_{e,S}}_\infty\le\norm{W_{e,S}}_{\cW}\),
\(\abs{m^e\prod\Lambda(u_i)^{e_i}}\le(c_d^{-1}N_0)^{\abs e}\le N_m^{\abs e}\), \(\abs{\lambda^S}=1\);
\(\ev_\lambda\) is multiplicative on \(\cA^{\rm full}\) (substitution of
commuting values into absolutely convergent series), hence on \(\cA\)
whenever \(P_\dg(WV)=WV\).
A product monomial contains \(\lambda_m\) only if at least one factor does
(\(\lambda_m^2=1\) can only remove it), which gives the seminorm
inequality.  In \(W(\lambda;\xi)\) the monomials containing \(\lambda_m\) are
those with \(m\in S\), of total mass
\(\le\sum_e\sum_{S\ni m}\norm{W_{e,S}}_\infty N_m^{\abs e}\le\norm W_{[m]}\), and
reversing \(\lambda_m\) changes exactly their sign.
\end{proof}

An element is \emph{symmetric} if invariant under simultaneous
permutation of tensor slots and of the indices of \(e\).

\paragraph{Coefficient extraction on representatives.}
For \(S'\subseteq[Q]\) and \(f\in\{0,\dots,\dg\}^{S'}\), let
\(\widehat f\) be its extension by zero to \([Q]\).  Define the specific
representative lift \(W^{(f;S')}\) by
\[
 \bigl(W^{(f;S')}\bigr)_{a,S}
 =\begin{cases}
   m^{\widehat f}(T)W_{a+\widehat f,S}(T),&a|_{S'}=0,\\
   0,&\text{otherwise}.
  \end{cases}
\]
This is a linear operation on the given array and involves no division
by \(m\).  Its realization is the bare \(X_{S'}^f\)-coefficient of
\(\realize(W)\):
\begin{equation}
 \begin{gathered}
  \realize(W)=\sum_{f\in\{0,\dots,\dg\}^{S'}}
       X_{S'}^f\,\realize\bigl(W^{(f;S')}\bigr),
  \qquad X_{S'}^f=\prod_{i\in S'}X_i^{f_i},\\
  W^{(f;S')}\text{ has zero site degree on }S',
  \qquad \norm{W^{(f;S')}}_\cA\le\norm W_\cA .
 \end{gathered}
 \label{C:eq:decomp}
\end{equation}
Indeed, multiplication of the representative by \(m^{\widehat f}\)
costs at most \(\norm m_{\cW_d}^{\abs f}\le N_m^{\abs f}\), exactly
compensated by removing \(\abs f\) from its weight.  The same argument
gives \(\norm{W^{(f;S')}}_{[m]}\le\norm W_{[m]}\).
If \(RW=W\), then \(RW^{(f;S')}=(-1)^{\abs f}W^{(f;S')}\).
The ordinary coefficients of \(\realize(W)\) are unique, but their
representative lifts here are the ones just specified; no uniqueness of
an arbitrary lift, or of an expansion after rough-field evaluation, is
asserted.

\paragraph{Ambient module and bare symbols.}
Let \(\cA^{\rm amb}\) be the space of bare-coefficient polynomials
\(\sum_{e,S}V_{e,S}(T)X^e\lambda^S\), \(e\in\cE\), with norm
\(\sum_{e,S}\norm{V_{e,S}}_{\cW_{dQ}}(c_d^{-1}N_0)^{\abs e}\).
The realization map is bounded:
\[
 \norm{\realize(W)}_{\cA^{\rm amb}}\le\norm W_\cA,
 \qquad
 \ev_\lambda\realize(W)=\ev_\lambda W.
\]
Bare powers \(X_i\), the substitution polynomials
\(\K_i^{(e,f)}\), and coefficient vectors such as
\(J^{(e)}/X^e:=[X^e]\realize(J^{(e)})\) are interpreted in this ambient
module.  Products used below never exceed the allowed degree.
An ambient identity is not used to infer an identity in \(\cA\).
Instead, the target in Section~\ref{C:sec:target} is assigned an explicit
representative array, with surplus nonnegative powers of \(m\) absorbed
in its representatives.  This array is the object to which
polarization, the norm estimates, and the fixed-point theorem apply.

\subsection{The ideal increment}
Let \(\Psi=\Psi_{\dg Q}(U)\) list the nonconstant monomials of degree
\(\le\dg Q\) in \(U\in\R^M\), \(N=N_{\dg Q}\) its length.  For a tuple
\(T\) and \(0\ne e\in\cE\) define
\begin{equation}
 J^{(e)}(T)=\chi_Q(T)\,t^{\abs e}\prod_i\frac{m(u_i)^{e_i}}{e_i!}\;
 \E_{\Pi_0}\Bigl[\prod_i\bigl(c_i(T)\cdot\nabla_U\bigr)^{e_i}\Psi(U)\Bigr]\ X^e,
 \label{C:eq:ideal}
\end{equation}
extended by zero across the diagonals, and \(J=\sum_{e\ne0}J^{(e)}\in\cA^N\).
\(J\) is symmetric, \(RJ^{(e)}=(-1)^{\abs e}J^{(e)}\), and contains no
Walsh symbols.  By Taylor's formula for the polynomial \(\Psi\),
\(\Psi(U+w)=\sum_{e\in\mathbb N^Q}\prod_i(w_i\cdot\nabla)^{e_i}\Psi(U)/e_i!\)
for \(w=\sum_iw_i\) (the operators commute), so
\(\sum_{e\in\cE\setminus0}J^{(e)}\) is the part of degree \(\le\dg\) per
site of \(\chi_Q\{\E_{\Pi_0}\Psi(U+t\sum_im(u_i)X_ic_i)-\E_{\Pi_0}\Psi(U)\}\),
the moment increment of the virtual shift \eqref{C:eq:shift}.  For \(\dg=1\)
it is the multilinear part.

\begin{lemma}[Cells see only degree \(\dg\) per slot]\label{C:lem:cells}
For \(n\in\cE\) let \(\ell_{n,T}\) be the linear functional on \(\R^N\)
with \(\ell_{n,T}(\Psi(U))=\prod_i(v_Q(u_i)^\top U)^{n_i}\) (\(\ell_{0,T}=0\)).
Then
\[
 \ell_{n,T}\bigl(\ev_\lambda J(T)\bigr)
 =\chi_Q\Bigl\{\E_{\Pi_0}\prod_i\bigl(v_Q(u_i)^\top U+tm(u_i)\Lambda(u_i)\bigr)^{n_i}
 -\E_{\Pi_0}\prod_i(v_Q(u_i)^\top U)^{n_i}\Bigr\},
\]
and \(\ell_{n,T}(\ev_\lambda J^{(e)})=\chi_Q\prod_i\binom{n_i}{e_i}(tm(u_i)\Lambda(u_i))^{e_i}
\,\E_{\Pi_0}\prod_i(v_Q(u_i)^\top U)^{n_i-e_i}\), zero unless \(e\le n\)
coordinatewise.
\end{lemma}
\begin{proof}
\(v_Q(u_i)^\top c_j=\one\{i=j\}\) gives
\(\prod_i(v_Q(u_i)^\top(U+t\sum_jm(u_j)c_j\Lambda(u_j)))^{n_i}
=\prod_i(v_Q(u_i)^\top U+tm(u_i)\Lambda(u_i))^{n_i}\), a polynomial of
degree \(n_i\le\dg\) in each \(\Lambda(u_i)\); so the terms of the Taylor
expansion with some exponent \(>\dg\) lie in \(\ker\ell_{n,T}\), and the
\(e\)-th term is the binomial expansion.
\end{proof}

\begin{lemma}[Wiener estimate]\label{C:lem:wiener}
\(J^{(e)}\in\cA^N\) with representatives free of the factors \(m(u_i)\),
and
\[
 \norm{J^{(e)}}_\cA\le C_Q\Lambda^{E_Q}N_m^{\abs e}(t/t_n)^{\abs e};
\]
hence, with
\begin{equation}
 B>E_Q+\dg dQ+1,\qquad
 \norm{J}_\cA\le C_Q\Lambda^{E_Q+\dg dQ-B}=:\delta_n\le C_Q\Lambda^{-1}\to0 .
 \label{C:eq:wiener}
\end{equation}
\end{lemma}
\begin{proof}
Each coordinate of \eqref{C:eq:ideal} is \(m^{e}(T)\) times a constant times
\(\chi_Qt^{\abs e}\) times a product of \(\abs e\le\dg Q\) Lagrange coordinates
(with repetitions: \(e_i\) coordinates of \(c_i\)) times \(X^e\), so the
representative is the latter.  Insert the shell partition
\(\sum_{\bm m}\rho_{\bm m}=1\); by Corollary~\ref{cor:configuration-multipliers}
(repetitions of Lagrange coordinates allowed) and
Lemma~\ref{lem:admissible-shells}, each of the
\(\le C_Q\{1+\log(1/t_n)\}^{E_Q}\) shells meeting the taper contributes
\(\le C_Q(t/t_n)^{\abs e}\) to the \(\cW_{dQ}\) norm.  Sum over
shells, coordinates, and \(e\), multiply by the weight
\(N_m^{\abs e}\le C_Q\Lambda^{d\abs e}\), and use \(t/t_n=\Lambda^{-B}\),
\(\log(1/t_n)=O(\Lambda)\), \(\abs e\le\dg Q\).
\end{proof}

\subsection{One-slot algebra, centring, polarization}
Let \(\cB=\cB_k\) be the \emph{one-slot algebra}: the construction
\eqref{C:eq:A-norm} with \(Q=1\), i.e.\ elements
\[
 g=\sum_{e\le\dg}\sum_{S}m(u)^eg_{e,S}(u)X^e\lambda^S,\qquad g_{e,S}\in\cW_d,\qquad
 \norm g_\cB=\sum_{e,S}\norm{g_{e,S}}_{\cW_d}N_m^e ,
\]
with the truncated product, evaluation, and reflection.  For \(g\in\cB\) and a slot \(s\in[Q]\), \(g(u_s)\) denotes the
element of \(\cA\) obtained by placing \(g\) in slot \(s\) (with
\(X\mapsto X_s\)) and \(1\) in the other slots; \(\norm{g(u_s)}_\cA=\norm g_\cB\).
For \(g_1,\dots,g_Q\in\cB\), \(\bigotimes_sg_s:=\prod_sg_s(u_s)\in\cA\), formed
without truncation since the slots are distinct, and
\(g^{\otimes Q}=\bigotimes_sg\).

Let \(\{w_\nu\}\) be the real trigonometric basis of \(L^2(\Torus^d)\):
\(w_0=1\), \(\norm{w_\nu}_\infty\le1\), \(\int w_\nu=0\) for \(\nu\ne0\);
\(\norm\cdot_{\cW_{dQ}}\) is (up to a constant) the \(\ell^1\) norm of the
coefficients in \(\{w_{\nu_1}\otimes\cdots\otimes w_{\nu_Q}\}\).  Expanding
every representative \(W_{e,S}\) in this basis, every \(W\in\cA\) is an
absolutely convergent sum of \emph{basis tensors}
\(\lambda^S\bigotimes_ig_i\) with the \emph{\(m\)-adapted slot functions}
\[
 g_i=m(u)^{e_i}w_{\nu_i}(u)X^{e_i}\in\cB,\qquad e_i\in\{0,\dots,\dg\},
\]
with weight \(N_m^{\abs e}\), and \(\sum(\text{coefficient}\times\text{weight})\le C\norm W_\cA\).
The multiplier \(m\) is kept inside the slot function and is never
expanded into Fourier modes; this is what makes the atoms below
\(m\)-adapted.  For \(e_i\ge1\) the evaluated slot function
\(m^{e_i}w_{\nu_i}\Lambda^{e_i}\) is not mean zero; its mean is the Walsh
polynomial
\begin{equation}
 \overline g_i=\int_{\Torus^d}m(u)^{e_i}w_{\nu_i}(u)\Lambda(u)^{e_i}\,du
 =\sum_{S:\abs S\le e_i}\lambda^S\,\overline g_{i,S},
 \label{C:eq:centring}
\end{equation}
an element of \(\cB\) (free of \(X\): \(e=0\), constant in \(u\)), regarded
as an element of \(\cA\) through the slot embedding.  Expanding
\(\Lambda^{e_i}\) into products of packets and using \(\abs m\le c_d^{-1}\):
\(\norm{\overline g_i}_W\le\int c_d^{-e_i}(\sum_m\abs{\psi_m})^{e_i}\le(c_d^{-1}N_0)^{e_i}\le N_m^{e_i}\),
and for every symbol \(m'\),
\(\norm{\overline g_i}_{[m']}\le e_ic_d^{-e_i}N_0^{e_i-1}\int\abs{\psi_{m'}}\le Ct^d\).
Put \(g^\circ_i=(g_i-\overline g_i)/(2N_m^{e_i})\) for \(e_i\ge1\) and
\(g^\circ_i=g_i/2\) for \(e_i=0\), \(\nu_i\ne0\); then
\(\norm{g^\circ_i}_\cB\le1\), \(\norm{\ev_\lambda g^\circ_i}_\infty\le1\)
(by \eqref{C:eq:Nm}), \(\int\ev_\lambda g^\circ_i\,du=0\) for every
\(\lambda\), \(\norm{g^\circ_i}_{[m']}\le Ct^d\), and
\(Rg^\circ_i=(-1)^{e_i}g^\circ_i\).

\begin{lemma}[Polarization]\label{C:lem:pol}
Fix \(\eta_f\in(0,1)\), \(\theta=\eta_f/(2\max_r\abs{t_r})\), and a
finite-dimensional normed space \(V\).  There is a countable family of
one-slot atoms \(F_m=1+\theta t_{r(m)}\overline g_m\in\cB\), closed under
\(R\), with \(\norm{\overline g_m}_\cB\le1\), \(\norm{\ev_\lambda\overline g_m}_\infty\le1\),
\(\int\ev_\lambda\overline g_m=0\) for all \(\lambda\), and linear maps
\(a_m:\cA^{\rm sym}\otimes V\to\cA^{\rm sc}\otimes V\) (\(\cA^{\rm sc}\): Walsh
polynomials with constant coefficients) such that every symmetric
\(A\in\cA\otimes V\) satisfies \(A=\sum_ma_m(A)F_m^{\otimes Q}\) absolutely in
\(\cA\), and
\begin{equation}
 \sum_m\norm{a_m(A)}\le C_{\rm pol}\norm A,\qquad
 \sum_m\bigl|\norm{a_m(A)}-\norm{a_m(A')}\bigr|\le C_{\rm pol}\norm{A-A'} .
 \label{C:eq:pol-bounds}
\end{equation}
Each \(\ev_\lambda F_m\) is a probability density on \(\Torus^d\) with
values in \([1-\eta_f/2,1+\eta_f/2]\), and
\(\norm{F_m^{\otimes Q}-1^{\otimes Q}}_\cA\le(1+\eta_f/2)^Q-1=:C_F\).
\end{lemma}
\begin{proof}
Treat one basis tensor \(\lambda^S\bigotimes_ig_i\); keep \(\lambda^S\) in
the coefficient.  Multilinearity of the tensor product gives
\(\bigotimes_ig_i=\sum_{I\subseteq[Q]}\prod_{i\notin I}\overline g_i\prod_{i\in I}2N_m^{e_i}
\bigotimes_{i\in I}g^\circ_i\otimes\bigotimes_{i\notin I}1\) (with
\(\overline g_i=0\) for \(e_i=0\), \(\nu_i\ne0\), and the slot \(w_0=1\)
treated as the unit), \(2^Q\) terms whose coefficients are Walsh
polynomials of norm \(\le2^Q\prod_iN_m^{e_i}\), i.e.\ \(\le2^Q\) times the
weight of the basis tensor.  Symmetrise, and apply
Lemma~\ref{lem:polarization-identity} in the commutative algebra \(\cB\) with the
symmetrised tensor product on \(\cA\): it expresses
\(g^\circ_{i_1}\odot\cdots\odot g^\circ_{i_j}\odot1^{\odot(Q-j)}\) through the
tensor powers \((1+\theta t_rg_\epsilon)^{\otimes Q}=\prod_s(1+\theta t_rg_\epsilon(u_s))\),
\(g_\epsilon=j^{-1}\sum_{i\in I}\epsilon_ig^\circ_i\in\cB\); only sums within
\(\cB\) and products across distinct slots occur, so no truncation is
active and the identity holds in \(\cA\).  The atoms \(1+\theta t_rg_\epsilon\)
are sums of \(m\)-adapted slot functions, hence in \(\cB\) and of degree
\(\le\dg\) in \(X\), and satisfy the three bounds; the coefficients are
\(C_j\theta^{-j}\lambda_{jr}\prod\epsilon_i\) times the Walsh-polynomial
coefficient above.  Since \(Rg^\circ_i=\pm g^\circ_i\), \(Rg_\epsilon=g_{\epsilon'}\)
for \(\epsilon'_i=(-1)^{e_i}\epsilon_i\), so the family is \(R\)-closed.
Summing over the expansion of \(A\) gives \eqref{C:eq:pol-bounds};
linearity of \(a_m\) holds because the centring elements are attached to
the basis element, not to \(A\).  The density statement follows from
\(\abs{\theta t_r\ev_\lambda g_\epsilon}\le\eta_f/2\) and \(\int\ev_\lambda g^\circ_i=0\).
\end{proof}

\subsection{The coefficient function in the chart}
With the cutoff \(\chi_0\) fixed at the start of this section, put
\begin{equation}
 \text{(a)}\quad \varkappa_k(u)=\chi_0(u)\,\ja^2G(u)^4,\qquad\qquad
 \text{(b)}\quad G_k^2(u)=\chi_0(u)\,G(u)^2 .
 \label{C:eq:coef-chart}
\end{equation}
\begin{lemma}\label{C:lem:coef-wiener}
\(\norm{\varkappa_k}_{\cW_d}\le C\ja^2\le1\) and \(\norm{G_k^2}_{\cW_d}\le C\),
uniformly in \(k\) and \(n\).
\end{lemma}
\begin{proof}
\(G=G_h\circ\theta_k^{-1}\) has all derivatives bounded (it varies at scale
\(h/r\ge1\) in the chart), so \(\chi_0G^2\) and \(\chi_0G^4\) are smooth
functions on \(\Torus^d\) with bounded \(C^{d+1}\) norm, hence bounded
Wiener norm.
\end{proof}
Since \(\chi_0=1\) on \(\operatorname{supp}m\), multiplying a coefficient
that carries \(m(u_i)\) by \(\varkappa_k(u_i)\) or \(G_k^2(u_i)\) is the same as
multiplying by \(\varkappa\) or \(G^2\) at the site.

\subsection{Substitution operators and the target}\label{C:sec:target}
In the chart, the substitution polynomials \eqref{C:eq:K} at slot \(i\)
read
\begin{equation}
\begin{array}{c|l}
 \text{(a)}&\K_i^{(1,0)}=X_i,\qquad \K_i^{(1,1)}=\varkappa_k(u_i),\\[2pt]
 \text{(b)}&\K_i^{(e,0)}=X_i^{e},\qquad \K_i^{(e,f)}=0\ (f\text{ odd}),\qquad \K_i^{(e,2)}=G_k^2(u_i)X_i^{e},
\end{array}
\label{C:eq:subst}
\end{equation}
In both cases \(\K_i^{(e,f)}=\varkappa^{(e,f)}(u_i)X_i^{e'(e,f)}\) with
\(e'(e,f)\in\{0,e\}\) and \(\norm{\varkappa^{(e,f)}}_{\cW_d}\le C\)
(Lemma~\ref{C:lem:coef-wiener}).
For \(0\ne e\in\cE\) with \(S'=\operatorname{supp}e\) define the linear
map \(\cS_e:\cA\to\cA^{\rm amb}\) using the prescribed extraction
\eqref{C:eq:decomp}:
\begin{equation}
 \cS_e(W)=\sum_{f\in\{0,\dots,\dg\}^{S'}}
 \realize\bigl(W^{(f;S')}\bigr)\prod_{i\in S'}\K_i^{(e_i,f_i)} .
 \label{C:eq:S-op}
\end{equation}
In words: in the \(X_{S'}\)-expansion of \(W\), the monomial \(X_{S'}^f\)
is replaced by \(\prod_i\K_i^{(e_i,f_i)}\); for \(f_i=0\) this inserts the
jitter \(X_i^{e_i}\), and for \(f_i\ge1\) the corrected jitter of
\eqref{C:eq:K}.  The \emph{target} is the representative-valued linear map described by
\begin{equation}
 \cJ(\mathbb D)=\sum_{0\ne e\in\cE}\cJ^{(e)}(\mathbb D),\qquad
 \realize\bigl(\cJ^{(e)}(\mathbb D)\bigr)
   =\frac{J^{(e)}}{X^e}\,\cS_e(\mathbb D),
 \label{C:eq:target}
\end{equation}
where \(J^{(e)}/X^e=[X^e]\realize(J^{(e)})\) is the \(X\)-free vector
coefficient in \eqref{C:eq:ideal}.  The second equality describes only
the realization; the following formula \emph{defines} its lift.
Let \(B_e(T)\) denote the sole vector representative of \(J^{(e)}\), so
\(\realize(J^{(e)})=m^eB_eX^e\).  For an input representative
\(\mathbb D_{a,S}\), put \(S'=\supp e\), \(f=a|_{S'}\), and
\[
 g_i=\begin{cases}e'(e_i,f_i),&i\in S',\\a_i,&i\notin S'.\end{cases}
\]
If any substitution factor is zero, its contribution is zero.
Otherwise add to the \((g,S)\)-representative of
\(\cJ^{(e)}(\mathbb D)\) the vector
\[
 \left\{\prod_{i\in S'}
 m(u_i)^{e_i+f_i-g_i}\varkappa^{(e_i,f_i)}(u_i)\right\}
 B_e(T)\,\mathbb D_{a,S}(T).
\]
Sum these contributions over \((a,S)\).  Every exponent
\(e_i+f_i-g_i\) is nonnegative, because \(g_i\in\{0,e_i\}\), and
\(g\in\cE\).  Thus this defines the representative without division by
\(m\) or an inverse of \(\realize\), even on the zero set of \(m\).
Realizing the formula gives precisely \eqref{C:eq:target}.

In case~(a), \(e=\one_{S'}\); in the ambient notation the formula reads
\[
 \realize\bigl(\cJ^{(S')}(\mathbb D)\bigr)
 =\frac{J^{(S')}}{X_{S'}}
   \sum_{A\subseteq S'}\realize\bigl(\mathbb D^{(A;S')}\bigr)
       \prod_{i\in A}\varkappa_i\prod_{i\in S'\setminus A}X_i.
\]
The term \(A=\varnothing\) is the jitter weighted by the extracted
zero-degree design coefficient.  The remaining terms are the
\emph{moment corrections} in which the site field is replaced by
\(\varkappa_i\) together with its prescribed design coefficient.

\begin{lemma}\label{C:lem:target}
For symmetric \(R\)-invariant \(\mathbb D\in\cA\): \(\cJ(\mathbb D)\in\cA^N\)
is symmetric, linear in \(\mathbb D\), satisfies
\(R\cJ^{(e)}(\mathbb D)=(-1)^{\abs e}\cJ^{(e)}(\mathbb D)\), and
\begin{equation}
 \norm{\cJ(\mathbb D)}_\cA\le C_Q\delta_n\norm{\mathbb D}_\cA,\qquad
 \norm{\cJ(\mathbb D)-\cJ(\mathbb D')}_\cA\le C_Q\delta_n\norm{\mathbb D-\mathbb D'}_\cA .
 \label{C:eq:target-norm}
\end{equation}
\end{lemma}
\begin{proof}
\emph{Membership and norm.}  Use the representative rule preceding the
lemma.  For one input \((a,S)\), the output degree is \(g\in\cE\);
on \(S'\) its surplus power is \(e_i+f_i-g_i\ge0\), and outside
\(S'\) the original degree \(a_i\) is unchanged.  The Wiener algebra
inequality and Lemma~\ref{C:lem:coef-wiener} give
\[
 \begin{aligned}
 &N_m^{\abs g}
 \left\|\left\{\prod_{i\in S'}
 m(u_i)^{e_i+f_i-g_i}\varkappa^{(e_i,f_i)}(u_i)\right\}
 B_e\mathbb D_{a,S}\right\|_{\cW_{dQ}}\\
 &\hspace{12mm}\le
 C_QN_m^{\abs e+\abs a}\norm{B_e}_{\cW_{dQ}}
                            \norm{\mathbb D_{a,S}}_{\cW_{dQ}} .
 \end{aligned}
\]
Here \(\norm m_{\cW_d}\le N_m\) and
\(\abs g+\sum_{i\in S'}(e_i+f_i-g_i)=\abs e+\abs a\).
Since \(N_m^{\abs e}\norm{B_e}=\norm{J^{(e)}}_\cA\), summing first over
\((a,S)\), then over \(e\ne0\) and vector coordinates, gives
\eqref{C:eq:target-norm} by \eqref{C:eq:wiener}.  The representative
formula is linear, so the difference estimate follows identically.

\emph{Symmetry and parity.}  The rule for \(g\) and for the surplus powers
commutes with simultaneous permutations of slots and the increment
pattern; summing over \(e\) proves symmetry.  If \(\mathbb D\) is
\(R\)-invariant, each nonzero representative has
\(\abs a+\abs S\) even.  For every nonzero substitution,
\(g_i\equiv e_i+f_i\pmod2\) on \(S'\).  Thus the output representative
satisfies
\[
 \abs g+\abs S\equiv\abs a+\abs S+\abs e
       \equiv\abs e\pmod2.
\]
This proves \(R\cJ^{(e)}(\mathbb D)=(-1)^{\abs e}\cJ^{(e)}(\mathbb D)\)
\emph{in the representative space}, not merely after realization.
\end{proof}

\begin{proposition}[Carrier]\label{C:prop:carrier}
For every active block \(k\) and all large \(n\) there is a symmetric,
\(R\)-invariant \(\mathbb D=\mathbb D_k\in\cA\) with
\(\norm{\mathbb D-1^{\otimes Q}}_\cA\le C_0\delta_n\), and for every sign
vector \(\lambda\) laws \(\nu_{P,k}(\cdot\mid\lambda)\), \(\nu_{Q,k}(\cdot\mid\lambda)\)
of a pair \((H_k,U_k)\) such that:
\begin{enumerate}[label=(\roman*),leftmargin=*]
\item the law of \(H_k\) is the same under both and does not depend on
\(\lambda\);
\item each label \(H\) carries a density \(F_{H}(\cdot;\lambda)=\ev_\lambda F\) on
\(\Torus^d\), \(F\in\cB\), with \(\int F_H=1\), \(1-\eta_f\le F_H\le1+\eta_f\), and
\(D(T;\lambda):=\E F_{H_k}^{\otimes Q}(T;\lambda)=\ev_\lambda\mathbb D(T)\), which
is even in \(\lambda\);
\item under both laws \(U_k\) is supported on the fixed finite support of
\(\Pi_0\); under \(\nu_{P,k}\), \(U_k\sim\Pi_0\) independently of \(H_k\);
\item pointwise on \((\Torus^d)^Q\),
\begin{equation}
 \E_{\nu_{Q,k}}\bigl[F_{H_k}^{\otimes Q}(T)\Psi(U_k)\bigr]
 -\E_{\nu_{P,k}}\bigl[F_{H_k}^{\otimes Q}(T)\Psi(U_k)\bigr]
 =\ev_\lambda\cJ(\mathbb D)(T).
 \label{C:eq:master}
\end{equation}
\end{enumerate}
The kernels depend on \(\lambda\) only through \((\lambda_m)_{m\in\cJ_k}\)
and are measurable: the label set is countable, and for each fixed label
the kernel depends on finitely many signs and takes finitely many
values.  Moreover
\(\mathbb D\in\cA\), so for \(1\le e\le\dg\) the realized coefficient
of \(X_i^e\) carries \(m(u_i)^e\) and vanishes where \(m(u_i)=0\); and
\begin{equation}
 \norm{\mathbb D}_{[m]}\le C_Q\delta_nt^d\qquad\text{for every symbol }m\in\cJ_k ,
 \label{C:eq:D-symbol}
\end{equation}
so that, in the frozen-field evaluation of Lemma~\ref{C:lem:algebra},
the Walsh monomials of \(\mathbb D\) containing a given symbol
\(\lambda_m\) have total mass \(\le C_Q\delta_nt^d\) uniformly on
\((\Torus^d)^Q\).  (The dependence of \(\ev_\lambda\mathbb D\) on
\(\lambda_m\) through the site values \(\Lambda_\lambda(u_i)\) is a separate
matter, handled in Lemma~\ref{C:lem:weights}.)
\end{proposition}

\begin{proof}
For symmetric \(R\)-invariant \(\mathbb D\) put
\(A=\cJ(\mathbb D)\), decompose \(A=\sum_ma_m(A)F_m^{\otimes Q}\) and
\(RA=\sum_ma_m(RA)F_m^{\otimes Q}\) (Lemma~\ref{C:lem:pol}); applying \(R\)
to the second, \(A=R(RA)=\sum_m\bigl(Ra_m(RA)\bigr)(RF_m)^{\otimes Q}\).  Set
\[
 \omega_m=K_0\bigl(\norm{a_m(A)}+\norm{a_m(RA)}\bigr),\qquad K_0>4/\delta_Q,
\]
\[
 \Phi(\mathbb D)=\Bigl(1-\sum_m\omega_m\Bigr)1^{\otimes Q}
 +\sum_m\omega_m\,\tfrac12\bigl\{F_m^{\otimes Q}+(RF_m)^{\otimes Q}\bigr\}.
\]
\(\Phi(\mathbb D)\) is symmetric and \(R\)-invariant.  By
\eqref{C:eq:pol-bounds} and \eqref{C:eq:target-norm},
\(\norm{\Phi(\mathbb D)-\Phi(\mathbb D')}\le C_0\delta_n\norm{\mathbb D-\mathbb D'}\)
with \(C_0=2K_0C_FC_{\rm pol}C_Q\), and \(\sum_m\omega_m\le C_1\norm{\mathbb D}\delta_n\),
which is \(\le1\) for large \(n\) on the ball below (\(\delta_n\to0\)), so
that the label \(0\) introduced below has nonnegative probability.
For large \(n\), \(\Phi\) maps the ball \(\norm{\mathbb D-1}\le2C_0\delta_n\) into
itself with contraction constant \(\le1/2\); let \(\mathbb D\) be its fixed
point.

Labels: for each \(m\) with \(\omega_m>0\) (labels with \(\omega_m=0\) are
omitted), \(m^+\) with density \(F_m\) and \(m^-\) with density \(RF_m\),
each of probability \(\omega_m/2\), and a label \(0\) with density \(1\);
the \(\omega_m\) are \(\lambda\)-free.  This gives (i), (ii) (using
Lemma~\ref{C:lem:pol}, \(\ev_{-\lambda}F_m=\ev_\lambda RF_m\), and
\(\E F_H^{\otimes Q}=\ev_\lambda\Phi(\mathbb D)=\ev_\lambda\mathbb D\), even in
\(\lambda\) by \(R\)-invariance).  Tilts: under \(Q\) and label \(m^\pm\),
\(U_k\sim\Pi_{v^\pm_m(\lambda)}\) with
\(v^+_m(\lambda)=\ev_\lambda a_m(A)/\omega_m\), \(v^-_m(\lambda)=\ev_\lambda Ra_m(RA)/\omega_m\),
so \(\norm{v^\pm_m(\lambda)}\le K_0^{-1}<\delta_Q\) (the evaluation of a Walsh
polynomial is bounded by its norm) and Lemma~\ref{lem:moment-simplex}
applies (with \(\Psi\) embedded in \(\Psi_{4Q}\)); under \(P\), \(U_k\sim\Pi_0\).  Then
\[
 \begin{aligned}
 \E_{\nu_Q}[F^{\otimes Q}\Psi]-\E_{\nu_P}[F^{\otimes Q}\Psi]
 &=\sum_m\tfrac{\omega_m}2
   \bigl[\ev_\lambda F_m^{\otimes Q}v^+_m
          +\ev_\lambda(RF_m)^{\otimes Q}v^-_m\bigr]\\
 &=\tfrac12\ev_\lambda A+\tfrac12\ev_\lambda R(RA)
   =\ev_\lambda A,
 \end{aligned}
\]
which is (iv).  The kernels read only the symbols occurring in \(\cA\).
The support statement is the definition of \(\cA\): \(\cJ(\mathbb D)\in\cA\)
for \(\mathbb D\in\cA\) (Lemma~\ref{C:lem:target}), the atoms are products of
\(m\)-adapted slot functions, and \(\Phi\) maps \(\cA\) into \(\cA\); hence the
fixed point lies in \(\cA\).  Finally \eqref{C:eq:D-symbol} is
Lemma~\ref{C:lem:symbol} below.
\end{proof}

\begin{lemma}[Symbol seminorm of the fixed point]\label{C:lem:symbol}
For every atom \(F\) of Lemma~\ref{C:lem:pol} and every symbol \(m\),
\(\norm{F^{\otimes Q}}_{[m]}\le C_Qt^d\).  Consequently, for every
\(\mathbb D\) in the ball of Proposition~\ref{C:prop:carrier},
\(\norm{\Phi(\mathbb D)}_{[m]}\le C_Q\delta_nt^d\), and the fixed point
satisfies \eqref{C:eq:D-symbol}.
\end{lemma}
\begin{proof}
The only symbols in an atom \(F=1+\theta t_rg_\epsilon\) come from the
centring polynomials \(\overline g_i\) in \(g^\circ_i\), and
\(\norm{g^\circ_i}_{[m]}\le Ct^d\) by \eqref{C:eq:centring}; hence
\(\norm{g_\epsilon}_{[m]}\le Ct^d\) and \(\norm{1+\theta t_rg_\epsilon(u_s)}_{[m]}\le Ct^d\)
for each slot \(s\), while \(\norm{1+\theta t_rg_\epsilon(u_s)}_\cA\le1+\eta_f/2\).
By the derivation inequality of Lemma~\ref{C:lem:algebra} applied to the
product over the \(Q\) slots, \(\norm{F^{\otimes Q}}_{[m]}\le Q(1+\eta_f/2)^{Q-1}Ct^d\).
The same holds for \(RF\), since \(R\) preserves the seminorm.  Now
\(\Phi(\mathbb D)-(1-\sum\omega_{m'})1=\sum_{m'}\omega_{m'}\tfrac12\{F_{m'}^{\otimes Q}+(RF_{m'})^{\otimes Q}\}\)
with \(\lambda\)-free numbers \(\omega_{m'}\ge0\), \(\sum_{m'}\omega_{m'}\le C_1\norm{\mathbb D}\delta_n\),
and the atom family does not depend on \(\mathbb D\); therefore
\(\norm{\Phi(\mathbb D)}_{[m]}\le\sum_{m'}\omega_{m'}C_Qt^d\le C_Q\delta_nt^d\).
The fixed point \(\mathbb D=\Phi(\mathbb D)\) inherits the bound.  (No
contraction argument in the seminorm is needed, because \(\Phi\) depends on
its argument only through the numbers \(\omega_{m'}\).)
\end{proof}

\subsection{Projectivity and the integrated identity}\label{C:sec:projectivity}
For \(q\le Q\), \(T_q=(u_1,\dots,u_q)\) and \(Z\in(\Torus^d)^{Q-q}\), all
atoms integrate to one for every \(\lambda\), so
\(\E_{\nu_S}[F^{\otimes q}(T_q)\Psi]=\int\E_{\nu_S}[F^{\otimes Q}(T_q,Z)\Psi]dZ\)
and \(D_q(T_q;\lambda)=\int D(T_q,Z;\lambda)dZ\).  Write
\(\mathbb D_q=\int\mathbb D(\cdot,Z)dZ\) as shorthand for evaluated
slot integration: the integrated site variables are replaced by their
rough-field values and then integrated, giving Walsh-polynomial
coefficients.  The retained coefficient lifts are always those
specified in \eqref{C:eq:decomp} before this integration.  This notation
does not assert that slot integration is an algebra homomorphism on
\(\cA\), or require a uniformly bounded representative Wiener norm for
\(\mathbb D_q\); the pointwise Walsh bounds below are sufficient.
Integrating
\eqref{C:eq:master} and applying \(\ell_{n,T_q}\), \(n\in\{0,\dots,\dg\}^q\), by
Lemma~\ref{C:lem:cells}:
\begin{multline}
 \ell_{n,T_q}\Bigl(\E_{\nu_Q}[F^{\otimes q}\Psi]-\E_{\nu_P}[F^{\otimes q}\Psi]\Bigr)\\
 =\sum_{0\ne e\le n}\ \sum_{f}\ \overline{\mathbb D}^{(f;S'_e)}_q
 \prod_{i\in S'_e}\binom{n_i}{e_i}\bigl(tm(u_i)\bigr)^{e_i}\K_i^{(e_i,f_i)}(\Lambda(u_i))
 \ \E_{\Pi_0}\prod_i(v_Q(u_i)^\top U)^{n_i-e_i},
 \label{C:eq:proj}
\end{multline}
\(S'_e=\operatorname{supp}e\), \(f\in\{0,\dots,\dg\}^{S'_e}\),
\(\K_i^{(e,f)}(\Lambda(u_i))\) the evaluated substitution polynomial, and
the \emph{tapered integrated coefficients}
\begin{equation}
 \overline{\mathbb D}^{(f;S')}_q(T_q;\lambda)=\int\ev_\lambda\mathbb D^{(f;S')}(T_q,Z)\,\chi_Q(T_q,Z)\,dZ .
 \label{C:eq:tapered}
\end{equation}
Compare with the untapered decomposition of the design weight,
\(D_q=\sum_{f}\prod_{i\in S'}\Lambda(u_i)^{f_i}\,D_q^{(f;S')}\),
\(D_q^{(f;S')}=\int\ev_\lambda\mathbb D^{(f;S')}dZ\):
\begin{equation}
 \begin{aligned}
 D^{(f;S')}_q-\overline{\mathbb D}^{(f;S')}_q
   &=\int\ev_\lambda\mathbb D^{(f;S')}(T_q,Z)
                 \{1-\chi_Q(T_q,Z)\}\,dZ,\\
 \norm{D^{(f;S')}_q(T_q;\cdot)
              -\overline{\mathbb D}^{(f;S')}_q(T_q;\cdot)}_W
   &\le C_Qg_q(T_q),
 \end{aligned}
 \label{C:eq:ghost-diff}
\end{equation}
where \(g_q(T_q):=\int\{1-\chi_Q(T_q,Z)\}\,dZ\).
The bound is in Walsh norm, not only absolute value.  Indeed,
Lemma~\ref{C:lem:algebra} and \eqref{C:eq:decomp} give
\(\norm{\ev_\lambda\mathbb D^{(f;S')}(T_q,Z)}_W
 \le\norm{\mathbb D^{(f;S')}}_\cA\le C_Q\) uniformly in \((T_q,Z)\).
Since the taper is nonnegative and independent of \(\lambda\), the
triangle inequality under the integral gives
\[
 \norm{D^{(f;S')}_q-\overline{\mathbb D}^{(f;S')}_q}_W
 \le\int(1-\chi_Q)\norm{\ev_\lambda\mathbb D^{(f;S')}}_W\,dZ
 \le C_Qg_q(T_q).
\]
The leakage also satisfies
\begin{equation}
 \int_{(\Torus^d)^q}g_q(T_q)\,dT_q\le C_Qt^d\Lambda^{M_Q},\qquad M_Q=Bd+Q-2,
 \label{C:eq:ghost}
\end{equation}
because \(\{\chi_Q\ne1\}\subseteq\{\min_iP_i\le2t_n\}\), whose volume is
\(\le C_Qt_n^d\{1+\log(1/t_n)\}^{Q-2}\) by Lemma~\ref{lem:bad-volume}.

\begin{lemma}[Structure of the design weight]\label{C:lem:weights}
Let \(T_q\) be a tuple of distinct points and \(S'\subseteq[q]\).
\begin{enumerate}[label=(\alph*),leftmargin=*]
\item \(D_q\) is even in \(\lambda\); \(D_q^{(f;S')}\) has parity \((-1)^{\abs f}\)
in \(\lambda\); \(\norm{D_q^{(f;S')}(T_q;\cdot)}_W\le C_Q\) and
\(\norm{D_q^{(f;S')}(T_q;\cdot)}_W\le C_Q\delta_n\) for \(f\ne0\)
(Walsh norm at fixed \(T_q\)).
\item \(D_q^{(f;S')}\) does not depend on the site values \(\Lambda(u_i)\),
\(i\in S'\).  Regard \(D_q^{(f;S')}(T_q;\cdot)\) as a Walsh polynomial in
\(\lambda\) at fixed \(T_q\), and let \(J_{S'}=\bigcup_{i\in S'}J(u_i)\).  A
sign \(\lambda_{m'}\), \(m'\in J_{S'}\), can enter it only through the
explicit symbols of \(\mathbb D\) (the centring polynomials), through
integrated slots, and through retained slots \(u_j\), \(j\notin S'\), that
share a packet with some \(u_i\), \(i\in S'\).  Outside this last
(collision) case, the Walsh monomials of \(D_q^{(f;S')}\) containing
\(\lambda_{m'}\) have total mass \(\le C_Q\delta_nt^d\), uniformly in
\(T_q\).  Consequently, writing \(\E_{J_{S'}}\) for the average over the
signs in \(J_{S'}\) (at most \(N_0q\) of them), outside the collision case
\[
 \bigl\|D_q^{(f;S')}-\E_{J_{S'}}D_q^{(f;S')}\bigr\|_W\le C_Q\delta_nt^d ,
\]
\(\E_{J_{S'}}D_q^{(f;S')}\) does not depend on the signs in \(J_{S'}\), and
both terms have parity \((-1)^{\abs f}\).
\item If \(m(u_i)=0\) (site \(i\) not assigned to this block) then
\(D_q\) does not contain \(\Lambda(u_i)\) at all.
\end{enumerate}
\end{lemma}
\begin{proof}
(a) \(R\)-invariance of \(\mathbb D\) and \eqref{C:eq:decomp}; the norms are
\(\norm{\mathbb D}\) and \(\norm{\mathbb D-1}\) since every \(X\)-monomial of
\(\mathbb D-1\) has norm \(\le C_0\delta_n\).  (b) \(D_q^{(f;S')}\) is free of \(X_{S'}\) by the decomposition.  Fix
\(i\in S'\) and a sign \(\lambda_{m'}\), \(m'\in J(u_i)\).  It enters
\(D_q^{(f;S')}\) in three ways.  Write
\(\mathbb D^{(f;S')}=\sum_{e,S}m^eW_{e,S}X^e\lambda^S\) (\(e|_{S'}=0\)), expand
every evaluated site value into products of packets,
\[
 \Lambda(u_i)^{e_i}=\sum_{m_1,\dots,m_{e_i}}\lambda_{m_1}\cdots\lambda_{m_{e_i}}\,\psi_{m_1}(u_i)\cdots\psi_{m_{e_i}}(u_i),
\]
and integrate over \(Z\).  Then \(D_q^{(f;S')}\) is a sum, over \((e,S)\) and
packet choices, of Walsh monomials with coefficients
\(\int m^eW_{e,S}\prod\psi\,dZ\), and the mass of the monomials containing
\(\lambda_{m'}\) is at most the sum of the absolute coefficients of the
contributions in which \(\lambda_{m'}\) occurs explicitly, i.e.\ \(m'\in S\)
or \(m'\) is one of the chosen packets (a product may cancel
\(\lambda_{m'}\) through \(\lambda_{m'}^2=1\), which only decreases the mass).
\emph{Explicit symbol \(m'\in S\):} the packet expansion of the site
values contributes mass \(\le(c_d^{-1}N_0)^{\abs e}\le N_m^{\abs e}\) per term, so
these contributions have mass
\(\le\sum_e\sum_{S\ni m'}\norm{W_{e,S}}_\infty N_m^{\abs e}\le\norm{\mathbb D^{(f;S')}}_{[m']}\le\norm{\mathbb D}_{[m']}\le C_Q\delta_nt^d\)
by \eqref{C:eq:D-symbol}.
\emph{Chosen packet \(m'\) at an integrated slot \(u_z\):} such terms have
\(e_z\ge1\), hence belong to \(\mathbb D-1\), so
\(\sum_{e_z\ge1,S}\norm{W_{e,S}}_{\cW}N_m^{\abs e}\le\norm{\mathbb D-1}_\cA\le C_0\delta_n\).
For one such term the sum of the absolute coefficients over the packet
choices in which one of the \(e_z\) factors at slot \(z\) reads \(m'\) is
at most
\begin{multline*}
 c_d^{-\abs e}e_z\norm{W_{e,S}}_\infty\int\abs{\psi_{m'}(u_z)}\Bigl(\sum_{a}\abs{\psi_a(u_z)}\Bigr)^{e_z-1}
 \prod_{j\ne z}\Bigl(\sum_{a}\abs{\psi_a(u_j)}\Bigr)^{e_j}dZ\\
 \le c_d^{-\abs e}e_zN_0^{\abs e-1}\norm{\psi_{m'}}_{L^1(\Torus^d)}\norm{W_{e,S}}_{\cW}
 \le Ct^dN_m^{\abs e}\norm{W_{e,S}}_{\cW},
\end{multline*}
(the product over \(j\ne z\) runs over the integrated and the retained
slots, the latter contributing their sup norms),
where \(\norm{\psi_{m'}}_{L^1}\le Ct^d\) because \(\psi_{m'}\) is bounded by one
and supported in a set of diameter \(\asymp t\) in the chart
(\(\psi^\flat_{m'}=G_h\phi((\cdot-x_0)/\ell-m')\) and \(\ell/r=t\)).  Summing
over the \(\le Q\) integrated slots and over \((e,S)\) with \(e_z\ge1\), these
contributions have mass \(\le C_Qt^d\cdot C_0\delta_n\).
\emph{Chosen packet \(m'\) at a retained slot \(u_j\), \(j\notin S'\):}
\(\psi_{m'}(u_j)\ne0\) requires \(J(u_j)\cap J(u_i)\ne\varnothing\), i.e.\ the
collision case, which is excluded wherever (b) is used.
Altogether the mass of the monomials containing \(\lambda_{m'}\) is
\(\le C_Q\delta_nt^d\) outside the collision case.  Finally,
\(D_q^{(f;S')}-\E_{J_{S'}}D_q^{(f;S')}\) consists exactly of the monomials of
\(D_q^{(f;S')}\) containing at least one sign of \(J_{S'}\) (averaging kills
these and keeps the others), so its Walsh norm is at most the sum over the
\(\le N_0q\) signs \(m'\in J_{S'}\) of the masses just bounded; the parity
statement holds because averaging over a set of signs removes monomials
without changing the degrees of the remaining ones.
(c) The coefficient of \(X_i^e\) in \(\mathbb D\) carries \(m(u_i)^e\)
(Proposition~\ref{C:prop:carrier}), and slot integration does not affect
the retained slot \(u_i\).
\end{proof}

\begin{lemma}[Local-sign removal across all incident blocks]\label{C:lem:cross-block}
Let \(\mathcal I\) be a set of at most \(Q\) distinct physical sites,
all interior for the assignment partition, with pairwise disjoint
packet sets \(J(x_i)\), \(i\in\mathcal I\).  Fix
\(S'\subseteq\mathcal I\), and put
\(J_{S'}=\bigcup_{i\in S'}J(x_i)\).
For any incident active block \(k\), let
\[
 \mathcal I_k=\{i\in\mathcal I:x_i\in S_k\},\quad
 T_k=(\theta_k(x_i))_{i\in\mathcal I_k},\quad q_k=\abs{\mathcal I_k},
 \quad S'_k=\{i\in S':k(i)=k\}.
\]
Identify physical site labels with their slots in \(T_k\).
For any \(f_k\in\{0,\dots,\dg\}^{S'_k}\), set
\(A_k=D_{k,q_k}^{(f_k;S'_k)}(T_k;\cdot)\), with
\(A_k=D_{k,q_k}\) when \(S'_k=\varnothing\).  Then
\[
 \norm{A_k-\E_{J_{S'}}A_k}_W\le C_Q\delta_nt^d .
\]
Both \(A_k\) and \(\E_{J_{S'}}A_k\) have parity
\((-1)^{\abs{f_k}}\), and the latter is independent of all signs in
\(J_{S'}\).  The conclusion includes blocks with \(S'_k=\varnothing\)
and signs read at sites assigned to other blocks.
\end{lemma}
\begin{proof}
Fix a packet index \(a\in J_{S'}\).  If \(a\notin\cJ_k\), the block
kernel and all its coefficients are independent of \(\lambda_a\).
Otherwise the collision-free assumption gives a unique
\(i\in S'\) with \(a\in J(x_i)\).  We first check that no retained
site-value factor in \(A_k\) can read this sign.

If \(i\in\mathcal I_k\) and \(k(i)=k\), the slot belongs to \(S'_k\),
and coefficient extraction removes its site variable.  If
\(i\in\mathcal I_k\) and \(k(i)\ne k\), interior assignment gives
\(\phi_k(x_i)=0\), hence \(m_k(x_i)=0\); every positive site-degree
term vanishes by the representative formula.  If
\(i\notin\mathcal I_k\), there is no such retained slot in this block.
Every other retained point \(x_j\) has \(a\notin J(x_j)\) by
collision-freeness.  The chart cutoff can only remove packet support,
not introduce a new packet at a retained point.

Consequently \(\lambda_a\) can enter only as an explicit symbol or
through a ghost slot.  For the explicit-symbol contributions, the
coefficient-extraction bound and Lemma~\ref{C:lem:symbol} give mass at
most
\[
 \norm{\mathbb D_k^{(f_k;S'_k)}}_{[a]}
 \le\norm{\mathbb D_k}_{[a]}
 \le C_Q\delta_nt^d.
\]
For a ghost-slot contribution, expand its rough-field power into
packets as in Lemma~\ref{C:lem:weights}(b).
A term using \(a\) at a ghost slot has positive degree there and hence
comes from \(\mathbb D_k-1\).  For each such representative, integrate
one factor \(\abs{\psi_a}\), using
\(\int\abs{\psi_a}\le Ct^d\), and bound the remaining packet sums by
\(N_0\), the multipliers by \(c_d^{-1}\), and the representative by
its Wiener norm.  Summing over the at most \(Q\) ghost slots costs
\[
 C_Qt^d\norm{\mathbb D_k-1}_\cA\le C_Q\delta_nt^d.
\]
Coefficient extraction is contractive and introduces only
deterministic powers of \(m\), so the same estimate applies for every
\(f_k\), including the empty extraction.  This proves that the
Walsh monomials of \(A_k\) containing \(\lambda_a\) have total mass
at most \(C_Q\delta_nt^d\).

Averaging over \(J_{S'}\) deletes exactly the monomials containing at
least one of those signs.  There are at most \(N_0Q\) such signs, so
summing the last bound proves the assertion.  Averaging removes
monomials without changing the parity of any monomial that remains;
Lemma~\ref{C:lem:weights}(a) therefore gives the parity statement.
\end{proof}

\subsection{Design density and tilt}
Given \(\lambda\) and labels \(H=(H_k)\), put
\begin{equation}
 \widetilde f_H(x)=\prod_{k:x\in S_k}F_{H_k}(\theta_k(x);\lambda),\qquad
 Z(H,\lambda)=\int\widetilde f_H,\qquad f_H=\widetilde f_H/Z .
 \label{C:eq:density}
\end{equation}
At most \(N_1\) blocks meet a point, so
\begin{equation}
 \Bigl(\frac{1-\eta_f}{1+\eta_f}\Bigr)^{N_1}\le f_H\le\Bigl(\frac{1+\eta_f}{1-\eta_f}\Bigr)^{N_1},
 \label{C:eq:density-bounds}
\end{equation}
and \(\eta_f\) is fixed so that these lie in \([\underline f,\overline f]\).
Let \(\nu_S=\bigl(\bigotimes_k\nu_{S,k}(\cdot\mid\lambda)\bigr)\otimes\Rad(\lambda)\)
be the law of \((H,\lambda,U)\) on side \(S\); the \((H,\lambda)\)-marginals
coincide (Proposition~\ref{C:prop:carrier}(i)).  Apply
Lemma~\ref{lem:tilt}: here the density depends on \(\lambda\) as well as on
\(H\), the tilted priors \eqref{eq:z-tilt} induce the same law of \(f_H\)
and the same \(X^n\)-marginal, and the mixed joint density is
\eqref{eq:z-cancel}, in which the untilted law has independent Rademacher packet signs
and conditionally independent block draws given those signs.
Unconditional independence is asserted only for components with
disjoint enlarged dependence sets, as verified in
Lemma~\ref{C:lem:links}.

\section{Case C: proof}\label{C:sec:proof}

\subsection{Mixtures and the two-point reduction}
Let \(\PP_n\) and \(\mathbb Q_n\) be the laws of \((X_i,A_i,Y_i)_{i\le n}\)
obtained by drawing \((H,U,\lambda)\sim\Pi_{S,n}\) (\(S=P,Q\)) and then
\(n\) iid triples from the \(P\)-row, respectively the \(Q\)-row, of
\eqref{C:eq:two-sides-a} (case~(a)) or \eqref{C:eq:two-sides-b} (case~(b)).
By Lemma~\ref{C:lem:legal} every law in either mixture belongs to \(\cP\),
and \(\tau(x_0)\) equals \(c_ac_b\Delta\) on the \(P\)-side and \(0\) on
the \(Q\)-side.  By Lemma~\ref{lem:two-point}, Theorem~\ref{thm:main}
in case~(C) follows from \(\TV(\PP_n,\mathbb Q_n)\le1/2\) for
\(n\ge n_\varepsilon\), since \(\Delta=n^{-z}=n^{-(\rho+\varepsilon)}\).

\subsection{Components and factorisation}
By \eqref{eq:z-cancel} the \(X^n\)-marginals coincide, so
\(\TV(\PP_n,\mathbb Q_n)\le\E\,\TV(\PP_n(\cdot\mid X^n),\mathbb Q_n(\cdot\mid X^n))\).
Fix \(x^n\).  For a site \(i\) let \(\cK_i=\{k\text{ active}:x_i\in S_k\}\)
and let the \emph{dependence set} be
\[
 \widetilde J_i=J(x_i)\cup\bigcup_{k\in\cK_i}\cJ_k ,
\]
the packets read at \(x_i\) or by any block meeting \(x_i\).  Declare
\(i\sim j\) if \(\cK_i\cap\cK_j\ne\varnothing\) or
\(\widetilde J_i\cap\widetilde J_j\ne\varnothing\), and let \(\{C\}\) be
the components of the transitive closure; \(q_C=\abs C\),
\(\cK(C)=\bigcup_{i\in C}\cK_i\), \(\widetilde J(C)=\bigcup_{i\in C}\widetilde J_i\).

\begin{lemma}[Geometry of links]\label{C:lem:links}
If \(i\sim j\) then \(\norm{x_i-x_j}\le C_Qr\).  A shared original packet
(\(J(x_i)\cap J(x_j)\ne\varnothing\)) forces \(\norm{x_i-x_j}\le R_0\ell\).
Distinct components have disjoint block sets and disjoint dependence
sets, and the primitive variables \((H_k,U_k)_{k\in\cK(C)}\),
\((\lambda_m)_{m\in\widetilde J(C)}\) of distinct components are
independent under \(\nu_S(\cdot\mid\lambda)\otimes\Rad\).
\end{lemma}
\begin{proof}
Blocks have diameter \(O(r)\) and a packet has diameter \(O(\ell)\le O(r)\);
a packet in \(\cJ_k\) lies within \(R_0\ell\) of \(S_k\).  So two sites
whose blocks share a packet, or whose packets or blocks intersect, are
within \(C_Qr\).  Disjointness follows from the definition of \(\sim\), and
independence from the product structure \eqref{eq:z-cancel}.
\end{proof}

The site likelihood \(L_i\) is a function of \((U_k)_{k\in\cK_i}\),
\((\lambda_m)_{m\in J(x_i)}\) and \(w_i=(R_i,T_i)\) (the functions \(G_h\),
\(\eta\), \(d_*\) are deterministic); the kernel of block \(k\) reads
\((\lambda_m)_{m\in\cJ_k}\) (Proposition~\ref{C:prop:carrier}).  With
\eqref{eq:z-cancel} and \eqref{C:eq:density}, the mixed joint density of
\((x^n,w^n)\) under side \(S\) equals
\[
 \mathcal Z_n^{-1}4^{-n}\int\prod_{k}\prod_{i:k\in\cK_i}F_{H_k}(\theta_k(x_i);\lambda)
 \prod_iL^S_i(w_i)\,d\nu_S(H,U\mid\lambda)\,d\Rad(\lambda),
\]
and by Lemma~\ref{C:lem:links} and Fubini the integral factorises over
components:
\begin{equation}
 \PP_n(w^n\mid x^n)=\prod_C\frac{N^P_C(w_C)}{N_C},\qquad
 \mathbb Q_n(w^n\mid x^n)=\prod_C\frac{N^Q_C(w_C)}{N_C},
 \label{C:eq:factor}
\end{equation}
\[
 N^S_C(w_C)=4^{-q_C}\int\prod_{k\in\cK(C)}\prod_{i\in C:k\in\cK_i}F_{H_k}(\theta_k(x_i);\lambda)
 \prod_{i\in C}L^S_i(w_i)\,d\nu_S\,d\Rad,
 \qquad N_C=\sum_{w_C}N^S_C(w_C),
\]
where \(N_C\) does not depend on \(S\): summing a cell over \(w_i\) gives
\(\sum_{w_i}L_i(w_i)=4\) for any parameter, so
\(N_C=\E_\lambda\E_{\nu}\prod_k\prod_iF_{H_k}=\E_\lambda\prod_{k\in\cK(C)}D_{q_k}(T_k;\lambda)\),
with \(T_k\) the tuple of the \(q_k\) sites of \(C\) in \(S_k\) (in the
chart of \(k\)).  Sites in no active block are singleton components with
\(L^P_i=L^Q_i=1\).

Hence, by Lemma~\ref{lem:hellinger}(a),(b), on the event that all
components have \(q_C\le Q\),
\begin{equation}
 \TV(\PP_n(\cdot\mid x^n),\mathbb Q_n(\cdot\mid x^n))
 \le\Bigl\{\sum_CH^2_C\Bigr\}^{1/2},\qquad
 H^2_C=H^2\Bigl(\frac{N^P_C}{N_C},\frac{N^Q_C}{N_C}\Bigr).
 \label{C:eq:tv-hell}
\end{equation}

\subsection{The component identity}\label{C:sec:component-identity}
Fix a component \(C\) with \(q_C\le Q\).  For \(i\in C\) write
\(\Lambda_i=\Lambda_\lambda(x_i)\), \(G_i=G_h(x_i)\), \(d_i=\ja\jb G_i^2\),
\(e_{k,i}=v_Q(\theta_k(x_i))^\top U_k\) for \(k\in\cK_i\), and
\(\cf_i=\cf(x_i)=\sum_{k\in\cK_i}\psi_k(x_i)e_{k,i}\); the shift variable at
the site is \(\zeta_i=\mathfrak a\,\cf_i\) with the smooth amplitude
\(\mathfrak a=b\) (case~(a)) or \(\mathfrak a=a\) (case~(b)), so that
\(\mathfrak a\psi_k(x_i)\,t\,m_k(x_i)=\jsh\phi_k(x_i)\).  Let \(L_i(\zeta)\),
\(\real_i(\zeta)\), \(\Delta^{(e)}_i(\zeta)\), \(\K_i^{(e,f)}\) be the objects
of Section~\ref{C:sec:bridge} at the site \(x_i\) with the cell
\(w_i=(R_i,T_i)\).  Then
\begin{equation}
 L^{Q,0}_i=L_i(\zeta_i),\qquad L^P_i=L_i(\zeta_i)+\real_i(\zeta_i),
 \label{C:eq:LQ0-LP}
\end{equation}
where \(L^{Q,0}_i\) is the \(Q\)-likelihood at the actual (unshifted)
parameter.  \(L_i(\zeta)\) is a polynomial of degree \(\dg\) in \(\zeta\),
hence of degree \(\le\dg\) in each \(e_{k,i}\), and affine in
\(\Lambda_i\), \(L_i=A_i+B_i\Lambda_i\) with \(B_i\) carrying the factor
\(\jfi\); the same holds for every \(\partial_\zeta^eL_i\).

\emph{\(P\)-side.}  Under \(\nu_P\), \(U\sim\Pi_0^{\otimes}\) independently of
\(H\), so
\begin{equation}
 N^P_C=4^{-q_C}\E_\lambda\Bigl[D_C\,\E_{\Pi_0^\otimes}\prod_{i\in C}\bigl(L_i(\zeta_i)+\real_i(\zeta_i)\bigr)\Bigr],
 \qquad D_C:=\prod_{k\in\cK(C)}D_{q_k}(T_k;\lambda).
 \label{C:eq:NP}
\end{equation}

\emph{\(Q\)-side.}  \(\prod_{i\in C}L_i(\zeta_i)\) is a polynomial in the
variables \(\{e_{k,i}\}\) of degree \(\le\dg\) in each, and for each block
\(k\) of degree \(\le\dg q_k\le\dg Q\) in \(U_k\).  Expand it into monomials
\(\prod_k\prod_ie_{k,i}^{n_{k,i}}\); the blocks are independent given
\(\lambda\), so the \(\nu_Q\)-expectation of a monomial times the design
factors is the product over \(k\) of
\(\E_{\nu_{Q,k}}[F^{\otimes q_k}\prod_ie_{k,i}^{n_{k,i}}]\), each given by
\eqref{C:eq:proj}.  Substituting and re-summing over the monomials, the
binomial coefficients reassemble into Taylor coefficients
(\(\sum_n(\text{coefficient})\binom{n}{e}x^{n-e}=\partial_x^eg(x)/e!\) for a
polynomial \(g\)), and \(\partial_{e_{k,i}}=\mathfrak a\psi_k(x_i)\partial_{\zeta_i}\).
This gives the exact expansion
\begin{multline}
 N^Q_C=4^{-q_C}\E_\lambda\sum_{(e_k,f_k)_k}\ \prod_{k\in\cK(C)}\Omega_k(f_k,e_k)\\
 \times\E_{\Pi_0^\otimes}\Bigl[\prod_{i\in C\setminus S'}L_i(\zeta_i)\prod_{i\in S'}
 \Bigl\{\prod_{k\in\cK_i:e_{k,i}\ge1}\frac{\phi_k(x_i)^{e_{k,i}}\K_i^{(e_{k,i},f_{k,i})}(\Lambda_i)}{e_{k,i}!}\Bigr\}
 \jsh^{e(i)}\partial_\zeta^{e(i)}L_i(\zeta_i)\Bigr],
 \label{C:eq:NQ}
\end{multline}
where \(e_k\in\{0,\dots,\dg\}^{T_k}\) is the increment pattern of block
\(k\) on its sites, \(f_k\in\{0,\dots,\dg\}^{\operatorname{supp}e_k}\) the
design exponents, \(e(i)=\sum_ke_{k,i}\), \(S'=\{i:e(i)\ge1\}\), and the
weights are the tapered integrated coefficients \eqref{C:eq:tapered} of
block \(k\),
\[
 \Omega_k(f_k,e_k)=\overline{\mathbb D}^{(f_k;\operatorname{supp}e_k)}_{k,q_k}(T_k;\lambda),
 \qquad \Omega_k(\varnothing,0)=D_{q_k}(T_k;\lambda).
\]
(If \(e(i)>\dg\) the factor \(\partial_\zeta^{e(i)}L_i\) vanishes.)
Which blocks can act at a site is read off from the explicit factor
\(\prod_{i:e_{k,i}\ge1}\phi_k(x_i)^{e_{k,i}}\) in \eqref{C:eq:NQ}: a
summand with \(e_{k,i}\ge1\) at a site with \(\phi_k(x_i)=0\) vanishes.
(The weight \(\Omega_k(f_k,e_k)\) itself need not vanish: for \(f_k=0\) it
is the integrated \(X\)-free coefficient, which is bounded but in
general not small, Lemma~\ref{C:lem:walsh-bounds}.  When \(f_{k,i}\ge1\) the weight does
vanish, because by Proposition~\ref{C:prop:carrier} the coefficient of
\(X_i^{f}\) in \(\mathbb D_k\) carries \(m_k(u_i)^{f}=0\); this is consistent,
since \(f_{k,i}\ge1\) requires \(e_{k,i}\ge1\).)  Hence every nonzero
summand has, at each site with \(e_{k,i}\ge1\), \(\phi_k(x_i)>0\), i.e.\
the site is assigned to \(k\) or lies in the transition layer; for an
interior site assigned to \(k\), \(\phi_k(x_i)=1\) and \(k\) is the only
block that can carry a nonzero increment at \(i\).  In case~(a) each \(e_{k,i}\in\{0,1\}\),
a site has only one factor \(\zeta_i\) in the multilinear cell, and
\eqref{C:eq:NQ} is the sum over pairwise disjoint \(S'_k\subseteq C\cap S_k\)
and \(A_k\subseteq S'_k\) (the sites with \(f=1\)) of
\[
 \prod_k\Omega_k(A_k,S'_k)\
 \E_{\Pi_0^\otimes}\Bigl[\prod_{i\notin S'}L_i\prod_k\Bigl\{\prod_{i\in S'_k\setminus A_k}\jb\phi_k(x_i)\Lambda_i\,\partial_yL_i
 \prod_{i\in A_k}\jb\phi_k(x_i)\varkappa_i\,\partial_yL_i\Bigr\}\Bigr].
\]

\subsection{Matching}
Let \(\mathrm{cl}(C)=\one\{\exists i\ne j\in C: J(x_i)\cap J(x_j)\ne\varnothing\}\),
\(\mathrm{tr}(C)=\one\{C\cap\mathcal L\ne\varnothing\}\), and
\(g_C=\sum_{k\in\cK(C)}g_{q_k}(T_k)\) with \(g_q\) from \eqref{C:eq:ghost-diff}.

\begin{lemma}[Walsh bounds]\label{C:lem:walsh-bounds}
For \(U\) in the legal ball and \(n\) large:
\(\norm{L_i(\zeta_i)}_W,\ \norm{\partial_\zeta^eL_i(\zeta_i)}_W\le C\) (\(e\le\dg\));
\(\norm{\real_i(\zeta_i)}_W\le C\Delta\); \(\jsh\le\jfi\), \(\jsh^2\le c_ac_b\Delta\);
\(\norm{\K_i^{(e,f)}(\Lambda_i)}_W\le C\);
\(\norm{\Omega_k(f,e)}_W\le C_Q\) for all \((f,e)\), and \(\le C_Q\delta_n\)
when \(f\ne0\) (for \(f=0\ne e\) the weight
\(\Omega_k(0,e)=\int\ev_\lambda\mathbb D^{(0;S')}\chi_Q\,dZ=\int\chi_Q\,dZ+O(\delta_n)\)
is uniformly \(O(1)\) and need not be \(O(\delta_n)\); it can vanish, e.g.\
when the tuple lies where the taper is zero, and only the upper bound is
used); the same for the untapered \(D^{(f;S')}_{q_k}\); and \(\Omega_k(f,e)\),
\(D^{(f;S')}_{q_k}\) have parity \((-1)^{\abs f}\) in \(\lambda\) (the taper
does not affect parity or norms).
\end{lemma}
\begin{proof}
\eqref{C:eq:L-real}, \(\abs{\zeta_i}\le C_Qc\), \(\norm{\Lambda_i}_W\le N_0\),
\(\ja,\jb\le1\), \(\abs{\eta_i}\le3\ja^2\), \eqref{C:eq:jsh-square},
\(\varkappa_i\le\ja^2\le1\), \(G_i^2\le1\), and Lemma~\ref{C:lem:weights}(a).
\end{proof}

\begin{lemma}[Parity count]\label{C:lem:parity}
Let \(\Theta\) be a Walsh polynomial in \(\lambda\) of parity
\((-1)^{\abs f}\), \(f=(f_k)_k\), and consider a term of the form
appearing in \eqref{C:eq:NQ},
\[
 \Xi=\Theta\cdot\prod_{i\in C\setminus S'}L_i(\zeta_i)\prod_{i\in S'}
 \Bigl\{\prod_{k:e_{k,i}\ge1}\phi_k(x_i)^{e_{k,i}}\K_i^{(e_{k,i},f_{k,i})}(\Lambda_i)\Bigr\}
 \jsh^{e(i)}\partial_\zeta^{e(i)}L_i(\zeta_i),\qquad S'\ne\varnothing .
\]
Then \(\abs{\E_\lambda\Xi}\le C_Q\norm\Theta_W\Delta\), uniformly in \(U\) in
the legal ball and in the cells.
\end{lemma}
\begin{proof}
Expand each cell factor \(L_i\) or \(\partial_\zeta^{e}L_i\) as
\(A+B\Lambda_i\) and choose, for each, either the \(\lambda\)-free part or
the part \(B\Lambda_i\) (which carries the factor \(\jfi\)); let \(\#y\) be
the number of factors \(B\Lambda_i\) chosen.  Every monomial of \(\Theta\)
has degree \(\equiv\abs f\), every monomial of a nonzero
\(\K_i^{(e,f)}(\Lambda_i)\) has degree \(\equiv e+f\pmod2\) (Section~\ref{C:sec:bridge}),
and \(\Lambda_i\) has degree one.  Since the parity of the degree of a
product of Walsh monomials is the sum of the parities, every monomial of
the chosen term has degree \(\equiv\abs f+\sum_{k,i}(e_{k,i}+f_{k,i})+\#y
\equiv\abs e+\#y\pmod2\), where \(\abs e=\sum_ie(i)\ge1\).  A Walsh polynomial
all of whose monomials have odd degree has zero \(\lambda\)-mean.  Hence
only the choices with \(\abs e+\#y\) even contribute, and for these the
Walsh norm of the term is \(\le C_Q\norm\Theta_W\jsh^{\abs e}\jfi^{\#y}\) by
Lemma~\ref{C:lem:walsh-bounds}.  If \(\abs e=1\) then \(\#y\ge1\) and
\(\jsh\jfi=c_ac_b\Delta\); if \(\abs e\ge2\) then \(\jsh^2\le c_ac_b\Delta\).
Summing over the \(\le2^{Q}\) choices proves the claim.
\end{proof}

\begin{proposition}[Component matching]\label{C:prop:matching}
For every component \(C\) with \(q_C\le Q\) and every cell \(w_C\),
\begin{equation}
 \abs{N^Q_C(w_C)-N^P_C(w_C)}\le4^{-q_C}C_Q\Delta\bigl[g_C+\mathrm{cl}(C)+\mathrm{tr}(C)+\delta_nt^d\bigr].
 \label{C:eq:matching}
\end{equation}
\end{proposition}

\begin{proof}
Write \(\widehat N^Q_C\) for \eqref{C:eq:NQ} with every tapered weight
\(\Omega_k(f_k,e_k)\), \((f_k,e_k)\ne(\varnothing,0)\), replaced by the
untapered \(D^{(f_k;\operatorname{supp}e_k)}_{q_k}\) (and
\(\Omega_k(\varnothing,0)=D_{q_k}\) unchanged).

\emph{Step 1: taper.}  \(N^Q_C-\widehat N^Q_C\) is a sum of \(C_Q\) terms,
each with \(S'\ne\varnothing\) and containing at least one factor
\(\Omega_k-D^{(f_k;\cdot)}_{q_k}\), of Walsh norm \(\le C_Qg_{q_k}(T_k)\) and
parity \((-1)^{\abs{f_k}}\) by \eqref{C:eq:ghost-diff} and
Lemma~\ref{C:lem:weights}(a).  The product \(\Theta\) of the weights and
differences has parity \((-1)^{\abs f}\) and Walsh norm
\(\le C_Qg_{q_k}(T_k)\), so Lemma~\ref{C:lem:parity} bounds each term by
\(C_Q\Delta g_{q_k}(T_k)\): \(\abs{N^Q_C-\widehat N^Q_C}\le4^{-q_C}C_Q\Delta g_C\).

\emph{Step 2: crude bound when \(\mathrm{cl}(C)=1\) or \(\mathrm{tr}(C)=1\).}
In \eqref{C:eq:NP}, every term with at least one \(\real_i\) has Walsh norm
\(\le C_Q\Delta\), and the term with none equals the
\((f_k,e_k)=(\varnothing,0)\) term of \(\widehat N^Q_C\).  Every other term
of \(\widehat N^Q_C\) has \(S'\ne\varnothing\) and is bounded by
Lemma~\ref{C:lem:parity} with \(\Theta=\prod_kD^{(f_k;\cdot)}_{q_k}\) by
\(C_Q\Delta\).  Hence \(\abs{\widehat N^Q_C-N^P_C}\le4^{-q_C}C_Q\Delta\).

\emph{Step 3: exact matching when \(\mathrm{cl}(C)=\mathrm{tr}(C)=0\).}
All sites are interior; site \(i\) is assigned to \(k(i)\).  By
Lemma~\ref{C:lem:weights}(c) and the remark after \eqref{C:eq:NQ}, only
\(k(i)\) can have a nonzero increment or positive retained site degree
at \(i\).  Other blocks can still read its signs through symbols or
ghost slots; this is controlled below.  If \(k(i)\) is inactive, then
\(G_i=\Lambda_i=d_i=0\), so every term choosing an increment or
\(\real_i\) at that site is zero.  Such zero terms may be omitted.  Hence in
\eqref{C:eq:NQ} the increment patterns are \(e=(e_i)_{i\in S'}\) with
\(e_i\in\{1,\dots,\dg\}\), \(\phi_{k(i)}(x_i)=1\), the exponents are
\(f=(f_i)_{i\in S'}\), and the untapered weights multiply to
\[
 \prod_kD^{(f_k;S'_k)}_{q_k}=D_C^{(f;S')},\qquad
 D_C=\sum_{f\in\{0,\dots,\dg\}^{S'}}\Lambda_{S'}^f\,D_C^{(f;S')},
\]
the coefficients prescribed by the representative extraction
\eqref{C:eq:decomp}, multiplied across blocks and then evaluated.
Only \(k(i)\) contributes a positive retained site degree at \(i\).
This is an identity for the specified expansion, not a claim of
uniqueness after rough-field evaluation.  Thus
\begin{equation}
 \widehat N^Q_C=4^{-q_C}\E_\lambda\sum_{S'\subseteq C}\sum_{f}D_C^{(f;S')}
 \E_{\Pi_0^\otimes}\Bigl[\prod_{i\notin S'}L_i(\zeta_i)\prod_{i\in S'}
 \sum_{e_i=1}^{\dg}\K_i^{(e_i,f_i)}(\Lambda_i)\,\Delta_i^{(e_i)}(\zeta_i)\Bigr],
 \label{C:eq:NQ-interior}
\end{equation}
with \(\Delta_i^{(e)}(\zeta)=\jsh^e\partial_\zeta^eL_i(\zeta)/e!\) as in
\eqref{C:eq:increments}, and by \eqref{C:eq:NP}
\begin{equation}
 N^P_C=4^{-q_C}\E_\lambda\sum_{S'\subseteq C}\sum_{f}\Lambda_{S'}^fD_C^{(f;S')}
 \E_{\Pi_0^\otimes}\Bigl[\prod_{i\notin S'}L_i(\zeta_i)\prod_{i\in S'}\real_i(\zeta_i)\Bigr].
 \label{C:eq:NP-interior}
\end{equation}
The terms with \(S'=\varnothing\) agree exactly and need no
replacement.  Fix \(S'\ne\varnothing\) and \(f\) for a remaining pair of
terms, and let \(J_{S'}=\bigcup_{i\in S'}J(x_i)\).  For every incident
block, including one with \(S'_k=\varnothing\), define
\[
 A_k=D_{k,q_k}^{(f_k;S'_k)}(T_k;\cdot),\qquad
 \widetilde D_k=\E_{J_{S'}}A_k,\qquad
 \widetilde D_C=\prod_{k\in\cK(C)}\widetilde D_k.
\]
Lemma~\ref{C:lem:cross-block}, applied with \(\mathcal I=C\), gives
\[
 \norm{A_k-\widetilde D_k}_W\le C_Q\delta_nt^d,
 \qquad \norm{A_k}_W+\norm{\widetilde D_k}_W\le C_Q,
\]
and both coefficients have parity \((-1)^{\abs{f_k}}\).
In particular, \(\widetilde D_C\) is independent of every sign in
\(J_{S'}\).  Averaging is performed factor by factor; no factorization
of \(\E_{J_{S'}}\) over the dependent block coefficients is assumed.

\emph{Exact matching of the replaced terms.}
Replace \(D_C^{(f;S')}\) by the same \(\widetilde D_C\) in
\eqref{C:eq:NQ-interior} and \eqref{C:eq:NP-interior}.
Since \(\mathrm{cl}(C)=0\), the fields \(\Lambda_i\), \(i\in S'\), are
independent of each other, of all remaining retained fields, of
\(\widetilde D_C\), and of the baseline coefficients
\(U\sim\Pi_0^\otimes\).  At fixed \(U\) and fixed remaining signs,
average these fields one site at a time.  The \(Q\)-side site factor is
\[
 \sum_{e=1}^{\dg}\E_{\Lambda_i}
       [\K_i^{(e,f_i)}(\Lambda_i)\Delta_i^{(e)}(\zeta_i)],
\]
whereas the \(P\)-side factor is
\(\E_{\Lambda_i}[\Lambda_i^{f_i}\real_i(\zeta_i)]\).
They agree by Lemma~\ref{C:lem:bridge} for every baseline value
\(\zeta_i\) and every binary cell.  All other factors are the same,
so the replaced \((S',f)\)-terms agree exactly.

\emph{Residual dependence.}
The difference
\(D_C^{(f;S')}-\widetilde D_C=\prod_kA_k-\prod_k\widetilde D_k\)
is a telescoping sum of at most \(\abs{\cK(C)}\le N_1Q\) products.
Each product contains one factor \(A_k-\widetilde D_k\) of Walsh norm
at most \(C_Q\delta_nt^d\), all other factors have Walsh norm at most
\(C_Q\), and the product has parity \((-1)^{\abs f}\).
Lemma~\ref{C:lem:parity} therefore bounds the replacement error in the
\(Q\)-side term by \(C_Q\Delta\delta_nt^d\).
On the \(P\)-side, \(S'\ne\varnothing\) guarantees at least one
\(\real_i\) factor of Walsh norm at most \(C\Delta\);
\(\norm{\Lambda_{S'}^f}_W\le C_Q\), so its replacement error has the
same bound directly.  There are only \(C_Q\) choices of \((S',f)\),
whence
\[
 \abs{\widehat N^Q_C-N^P_C}
       \le4^{-q_C}C_Q\Delta\delta_nt^d
 \quad\text{when }\mathrm{cl}(C)=\mathrm{tr}(C)=0.
\]

Combining the three steps gives \eqref{C:eq:matching}.
\end{proof}

\begin{lemma}\label{C:lem:HC}
For every component with \(q_C\le Q\),
\begin{equation}
 H^2_C\le C_Q\Delta^2\bigl[g_C+\mathrm{cl}(C)+\mathrm{tr}(C)+\delta_nt^d\bigr].
 \label{C:eq:HC}
\end{equation}
\end{lemma}
\begin{proof}
\emph{Normalisation.}  \(N_C=\E_\lambda D_C\in[(1-\eta_f)^{QK},(1+\eta_f)^{QK}]\)
with \(K=\abs{\cK(C)}\le N_1Q\).
\emph{Cell lower bound.}  For every \(\lambda\), \(U\) in the legal ball and
\(w_i\), \eqref{eq:walsh} gives
\(L^P_i\ge1-\abs{p}-\abs{y+\tau(1+p)}-\abs{py+\tau(1+p)}
\ge1-C_Qc_a-C_Qc_b-o(1)\ge\tfrac12\) for small \(c_a,c_b\) and large \(n\)
(Lemma~\ref{C:lem:legal}); hence \(\E_{\Pi_0^\otimes}\prod_iL^P_i(w_i)\ge2^{-q_C}\)
and \(N^P_C(w_C)/N_C\ge8^{-Q}=:c_Q\) by \eqref{C:eq:NP}.
\emph{Conclusion.}  By Lemma~\ref{lem:hellinger}(c) and
\eqref{C:eq:matching}, with \(\varrho_C=g_C+\mathrm{cl}+\mathrm{tr}+\delta_nt^d\le C_Q\),
\(H^2_C\le c_Q^{-1}4^{Q}\,N_C^{-2}(4^{-q_C}C_Q\Delta\varrho_C)^2\le C_Q\Delta^2\varrho_C^2\le C_Q\Delta^2\varrho_C\).
\end{proof}

\subsection{Hellinger sums}
Let \(m_I\) be the marginal density of \((X_i)_{i\in I}\) under the common
design law; \(m_I\le\overline f^{\,\abs I}\) since \(f_H\le\overline f\).
Let \(\mathcal A_n=\bigcup_{k\ \rm active}S_k\) be the active region;
\(\abs{\mathcal A_n}\le C h^d\) (there are \(O(h^dr^{-d})\) active blocks of
volume \(Cr^d\)), and \(\operatorname{supp}G_h\subseteq\mathcal A_n\).  A
site outside \(\mathcal A_n\) has \(\cK_i=\varnothing\) and
\(J(x_i)=\varnothing\), hence is a singleton component with
\(L^P_i=L^Q_i=1\) and \(H^2_C=0\); only components meeting
\(\mathcal A_n\) contribute.  Call \(I\) \emph{active} if it is connected
(for \(\sim\)) and \(x_I\cap\mathcal A_n\ne\varnothing\); this is a property
of \(x_I\), and the right side of \eqref{C:eq:HC} for \(C=I\) depends only
on \(x_I\).  Since components are connected,
\begin{equation}
 \E\sum_{C:q_C\le Q}H^2_C
 \le C_Q\Delta^2\sum_{q=1}^Q\binom nq\int\one\{I\text{ active}\}\bigl[g_I+\mathrm{cl}(I)+\mathrm{tr}(I)+\delta_nt^d\bigr]m_I\,dx_I .
 \label{C:eq:hell-sum}
\end{equation}

\begin{lemma}\label{C:lem:sums}
For \(\abs I=q\le Q\), with \(V_q:=\int\one\{I\text{ active}\}m_I\,dx_I\le C_Qh^d(C_Qr)^{d(q-1)}\),
\begin{align*}
 n^q\int\one\{I\text{ active}\}g_I\,m_I\,dx_I&\le C_Q\,nh^dt^d\Lambda^{M_Q}(nr^d)^{q-1},\\
 n^q\int\one\{I\text{ active}\}\mathrm{cl}(I)\,m_I\,dx_I&\le C_Q\,n^2h^d\ell^d(nr^d)^{q-2}\quad(q\ge2),\\
 n^q\int\one\{I\text{ active}\}\mathrm{tr}(I)\,m_I\,dx_I&\le C_Q\,nh^dt^d(nr^d)^{q-1},\\
 n^q\delta_nt^dV_q&\le C_Q\,nh^dt^d\delta_n(nr^d)^{q-1}.
\end{align*}
Consequently \(\E\sum_{C:q_C\le Q}H_C^2\le C_Q\Delta^2\{nh^dt^d\Lambda^{M_Q}+n^2h^d\ell^d\}\to0\).
\end{lemma}
\begin{proof}
By Lemma~\ref{C:lem:links} every link has length \(\le C_Qr\).  Order the
sites of an active \(I\) so that the first site lies in \(\mathcal A_n\) and
each later site is linked to an earlier one (\(C_Q\) orderings; a
connected set meeting \(\mathcal A_n\) admits such an ordering); the first
site has volume \(\abs{\mathcal A_n}\le Ch^d\), each later one
\((C_Qr)^d\); this gives \(V_q\) (for \(q=1\) it is \(\abs{\mathcal A_n}\)).

\emph{Ghost.}  Fix the block \(k_0\) in \(g_I\) and its sites \(I_0\),
\(q_0=\abs{I_0}\ge1\); order \(I\) with \(I_0\) first.  There are
\(O(h^dr^{-d})\) active blocks; given \(k_0\), the sites of \(I_0\) lie in
\(S_{k_0}\) and contribute \(\int_{(\Torus^d)^{q_0}}g_{q_0}\,dT\le C_Qt^d\Lambda^{M_Q}\)
by \eqref{C:eq:ghost}, i.e.\ \(C_Qr^{dq_0}t^d\Lambda^{M_Q}\) in physical
units; the other \(q-q_0\) sites contribute \((C_Qr)^d\) each.

\emph{Collision.}  Choose the pair with a shared packet: the first site
in \(\operatorname{supp}G_h\) (\(Ch^d\)), the second within \(R_0\ell\)
(\(C\ell^d\)), the rest \((C_Qr)^d\) each.

\emph{Transition.}  Choose the site in \(\mathcal L\); since \(I\) is
active, connected, and all links have length \(\le C_Qr\), this site lies
within \(C_QQr\) of \(\mathcal A_n\), i.e.\ in the transition layer of one
of the \(O(h^dr^{-d})\) blocks within that distance of the active region,
a set of volume \(\le C_dt^d\cdot O(h^d)\) in total; the rest
\((C_Qr)^d\) each.

\emph{Residual.}  \(V_q\) times \(\delta_nt^d\).

\emph{Conclusion.}  \(nr^d=n^{-cd\varepsilon}\le1\) and \(\delta_n\le1\), so all
four bounds are \(\le C_Q\{nh^dt^d\Lambda^{M_Q}+n^2h^d\ell^d\}\); apply
\eqref{eq:singleton}--\eqref{eq:collision}.
\end{proof}

\subsection{High components}\label{C:sec:high}
Let \(\mathcal E_n\) be the event that some component has more than
\(Q\) sites.  Such a component meets \(\mathcal A_n\) (components outside
\(\mathcal A_n\) are singletons), so some active set \(I\) of exactly
\(Q+1\) sites exists, and by the ordering argument
\[
 \PP(\mathcal E_n)\le\binom n{Q+1}\int\one\{I\text{ active}\}m_I\,dx_I
 \le C_Q\overline f^{\,Q+1}n^{Q+1}h^d(C_Qr)^{dQ}\to0
\]
by \eqref{eq:high}.

\subsection{Conclusion}
By Lemma~\ref{lem:components} (the conditional laws factor over the
components on \(\mathcal E_n^c\)), Lemma~\ref{C:lem:sums} and
Section~\ref{C:sec:high},
\[
 \TV(\PP_n,\mathbb Q_n)\le\PP(\mathcal E_n)+\Bigl\{\E\sum_{C:q_C\le Q}H^2_C\Bigr\}^{1/2}\longrightarrow0 ,
\]
so \(\TV(\PP_n,\mathbb Q_n)\le1/2\) for \(n\ge n_\varepsilon\), and the
two-point reduction (Lemma~\ref{lem:two-point}) proves
Theorem~\ref{thm:main} in case~(C) for all sufficiently small fixed
\(\varepsilon>0\).  For an arbitrary prescribed
\(\varepsilon>0\), apply the construction with a fixed
\(0<\varepsilon_0\le\varepsilon\) in that range; since
\(n^{-(\rho+\varepsilon_0)}\ge n^{-(\rho+\varepsilon)}\), the stated bound
follows after changing the constant.  \qed

\section{Case C: remark}\label{C:sec:remarks}

\begin{remark}[What the carrier does]
The common design density may depend on the rough field: a label law
independent of \(\lambda\) does not make its evaluated density atoms
independent of \(\lambda\).  The construction chooses the \(Q\)-side
conditional law of the smooth coefficients to reproduce the
\emph{design-weighted} moment increments.  The substitution polynomials
\eqref{C:eq:K} supply the design-weighted part of the \(P\)-side real
field.  For an interior, collision-free component, the untapered cells
match exactly after the block coefficients have been averaged over the
selected local signs.  Lemma~\ref{C:lem:cross-block} controls the error
of this replacement, including in blocks with no increment.

Only upper bounds on design coefficients and their residual local-sign
dependence are used: Lemma~\ref{C:lem:weights} and
Lemma~\ref{C:lem:cross-block} give the required bounds.  No sharp
asymptotic order for a particular pair-design coefficient is claimed.
The remaining discrepancy is charged to the taper, collision,
transition, and residual terms in Proposition~\ref{C:prop:matching}.
The decisive properties are bounded site degree, assignment of each
interior site's jitter to one block, and reflection symmetry for the
parity count.  The representative-space formulation keeps these
properties valid without identifying representative arrays that have
the same realization.  None of these corrections changes the scales
or proves the endpoint \(\varepsilon=0\) in case~(C).
\end{remark}


\begin{thebibliography}{9}
\bibitem{KBRW} E.~H. Kennedy, S.~Balakrishnan, J.~M. Robins, L.~Wasserman.
Minimax rates for heterogeneous causal effect estimation.
\emph{Annals of Statistics} 52(2), 793--816, 2024.
\bibitem{Tsybakov} A.~B. Tsybakov. \emph{Introduction to Nonparametric
Estimation}. Springer, 2009.
\bibitem{ODonnell} R.~O'Donnell. \emph{Analysis of Boolean Functions}.
Cambridge University Press, 2014.
\end{thebibliography}

\end{document}
